%% toc: P2 10 10.1,10.2,10.5,10.6
\section{Proofs for Lindstr\"om's theorem with function symbols and for Barwise's survey claims}
\label{sec:annex-p02-w02c-p02-003}

The notation is that of \cref{sec:p02-w02c-p02-003-general-lindstrom}.
For an assignment \(s\), a variable \(x\) and an element \(a\),
\(s(a/x)\) is the assignment that agrees with \(s\) except that it sends
\(x\) to \(a\).

\begin{annexitem}{annex:p02-w02c-p02-003-translations}{thm:p02-w02c-p02-003-translations}{Translations to the relational image}
\begin{proof}
\(\mathfrak A\) and \(\mathfrak A^{r}\) have the same domain, so
assignments in one are assignments in the other.

\emph{Part~1}, by induction on \(\chi\).  If \(\chi\) is
\(R_f(x_1,\ldots,x_n,y)\), then \(\mathfrak A^{r}\models\chi[s]\) iff
\((s(x_1),\ldots,s(x_n),s(y))\in R_f^{\mathfrak A^{r}}\) iff
\(f^{\mathfrak A}(s(x_1),\ldots,s(x_n))=s(y)\) iff
\(\mathfrak A\models(f(x_1,\ldots,x_n)=y)[s]\), and the last formula is
\(\chi^{-}\).  If \(\chi\) is \(P_c(x)\), then
\(\mathfrak A^{r}\models\chi[s]\) iff \(s(x)=c^{\mathfrak A}\) iff
\(\mathfrak A\models(c=x)[s]\).  If \(\chi\) is \(R(x_1,\ldots,x_n)\) with
\(R\in L\) or \(x=y\), then \(\chi^{-}=\chi\) and \(R\) has the same
interpretation in both structures.  The operation \({}^{-}\) commutes with
the connectives, including the countable \(\bigwedge\) and \(\bigvee\) of
\(L_{\omega_1\omega}\), and with the quantifiers; the satisfaction clauses for them refer
only to the satisfaction of the immediate subformulas at the same or at
modified assignments over the common domain, and the induction hypothesis
holds for all assignments; so the equivalence passes to compound formulas.

\emph{Part~2}, by induction on \(t\).  For a variable \(x\),
\(E_x(y)\) is \(x=y\), true at \(s\) iff \(s(y)=s(x)=x^{\mathfrak A}(s)\).
For a constant \(c\), \(E_c(y)\) is \(P_c(y)\), true iff
\(s(y)=c^{\mathfrak A}\).  Let \(t=f(t_1,\ldots,t_n)\) and let
\(z_1,\ldots,z_n\) be the new variables of \(E_t(y)\); they are distinct
from \(y\) and from the variables of \(t\).  Then
\(\mathfrak A^{r}\models E_t(y)[s]\) iff there are \(b_1,\ldots,b_n\) with
\(\mathfrak A^{r}\models E_{t_i}(z_i)[s']\) for each \(i\) and
\((b_1,\ldots,b_n,s(y))\in R_f^{\mathfrak A^{r}}\), where
\(s'=s(b_1/z_1)\cdots(b_n/z_n)\).  The variable \(z_i\) does not occur in
\(t_i\), and \(s'\) agrees with \(s\) on the variables of \(t_i\), so by
the induction hypothesis the first condition says \(b_i=t_i^{\mathfrak A}(s)\).
The second says \(f^{\mathfrak A}(b_1,\ldots,b_n)=s(y)\).  Together they hold
for some \(b_1,\ldots,b_n\) iff
\(s(y)=f^{\mathfrak A}(t_1^{\mathfrak A}(s),\ldots,t_n^{\mathfrak A}(s))=t^{\mathfrak A}(s)\).

\emph{Part~3}, by induction on \(\chi\).  For an atomic
\(R(t_1,\ldots,t_n)\), \(\mathfrak A\models\chi[s]\) iff
\((t_1^{\mathfrak A}(s),\ldots,t_n^{\mathfrak A}(s))\in R^{\mathfrak A}\) iff
there are \(b_1,\ldots,b_n\) with \(b_i=t_i^{\mathfrak A}(s)\) for each
\(i\) and \((b_1,\ldots,b_n)\in R^{\mathfrak A^{r}}\); by part~2, applied at
the assignment that sends the new variables \(z_i\) to \(b_i\), this is
\(\mathfrak A^{r}\models\chi^{u}[s]\).  An atomic \(t=t'\) is treated in
the same way with one new variable.  For \(\neg\), \(\land\) and the
countable \(\bigwedge\), \(\bigvee\) of \(L_{\omega_1\omega}\), the clause
for \(\chi\) refers only to the satisfaction of its immediate subformulas
at \(s\), and \({}^{u}\) commutes with the operation.  For \(\exists x\)
and \(\forall x\), the clause refers to the subformula at the assignments
\(s(a/x)\) over the common domain.  For \(Qx\,\psi\), with \(Q\) read as
in \(L(Q_0)\) or \(L(Q)\), the clause asks whether
\(\{a:\mathfrak A\models\psi[s(a/x)]\}\) is infinite, respectively
uncountable; by the induction hypothesis this set equals
\(\{a:\mathfrak A^{r}\models\psi^{u}[s(a/x)]\}\), so the two clauses
agree.  The unnesting of an atomic formula is a finite first-order formula
of \(L^{r}\), and \({}^{u}\) changes nothing else, so \(\chi^{u}\) is a
formula of \(L^{r}\) in the same logic, and a countable conjunction or
disjunction stays countable.
\end{proof}
\end{annexitem}

\begin{annexitem}{annex:p02-w02c-p02-003-general-lindstrom}{thm:p02-w02c-p02-003-general-lindstrom}{Barwise's 5.7 for general languages}
\begin{proof}
\emph{The relational part satisfies the hypotheses of
\cref{thm:p02-w02b-p02-001-barwise-lindstrom}.}  By condition~1 of
\cref{def:p02-w02c-p02-003-general-reading}, the restriction
\(\mathcal L^{\ast}_{\mathrm{rel}}\) of \(\mathcal L^{\ast}\) to relational
languages is a regular abstract logic in the relational sense.  A
relational signature is a language, so the Compactness and
L\"owenheim--Skolem hypotheses, which hold for every language, hold in
particular for every relational signature; these are parts~1 and~2 of
\cref{def:p02-w02b-p02-001-lindstrom-reading} for
\(\mathcal L^{\ast}_{\mathrm{rel}}\).  So
\cref{thm:p02-w02b-p02-001-barwise-lindstrom} applies to
\(\mathcal L^{\ast}_{\mathrm{rel}}\): every sentence \(\theta\) of it, with
finite occurrence signature \(\tau\), has a relational first-order sentence
\(\theta'\) of \(\tau\) such that \(\mathfrak C\models\theta\) iff
\(\mathfrak C\models\theta'\) for every \(\tau\)-structure \(\mathfrak C\).
This is the only use of the Compactness and L\"owenheim--Skolem
hypotheses.

\emph{The sentence \(\psi\).}  Let \(\varphi\in\mathcal L^{\ast}(L)\) have
occurrence language \(L_0\), so \(\varphi\in\mathcal L^{\ast}(L_0)\) by
condition~2.  Condition~3 for \(L_0\) gives
\(\varphi^{r}\in\mathcal L^{\ast}(L_0^{r})\) with
\(\mathfrak A_0\models\varphi\) iff \(\mathfrak A_0^{r}\models\varphi^{r}\)
for every structure \(\mathfrak A_0\) for \(L_0\).  The language
\(L_0^{r}\) is relational, so \(\varphi^{r}\) is a sentence of
\(\mathcal L^{\ast}_{\mathrm{rel}}\); let \(\tau\subseteq L_0^{r}\) be its
finite occurrence signature, and let \(\psi_1\) be the relational
first-order sentence of \(\tau\) given by the previous paragraph.  Put
\(\psi=\psi_1^{-}\) (\cref{def:p02-w02c-p02-003-relational-image}).  Each
symbol of \(\psi\) is a relation symbol of \(\tau\cap L_0\), a function
symbol \(f\) with \(R_f\in\tau\), or a constant symbol \(c\) with
\(P_c\in\tau\); all of these lie in \(L_0\), so \(\psi\) is a first-order
sentence of \(L_0\).

\emph{Equivalence.}  Let \(\mathfrak A\) be a structure for \(L\) and
\(\mathfrak A_0=\mathfrak A\restriction L_0\).  Then
\begin{align*}
\mathfrak A\models\varphi
 &\iff \mathfrak A_0\models\varphi
   &&\text{(i)}\\
 &\iff \mathfrak A_0^{r}\models\varphi^{r}
   &&\text{(ii)}\\
 &\iff \mathfrak A_0^{r}\restriction\tau\models\varphi^{r}
   &&\text{(iii)}\\
 &\iff \mathfrak A_0^{r}\restriction\tau\models\psi_1
   &&\text{(iv)}\\
 &\iff \mathfrak A_0^{r}\models\psi_1
   &&\text{(v)}\\
 &\iff \mathfrak A_0\models\psi
   &&\text{(vi)}\\
 &\iff \mathfrak A\models\psi
   &&\text{(vii)}.
\end{align*}
The reasons: (i) condition~2; (ii) condition~3; (iii) \(\tau\) is the
occurrence signature of \(\varphi^{r}\), condition~1; (iv) the choice of
\(\psi_1\); (v) first-order satisfaction depends only on the symbols that
occur; (vi) part~1 of \cref{thm:p02-w02c-p02-003-translations}, applied to
the language \(L_0\) and to \(\psi_1\), a first-order sentence of
\(\tau\subseteq L_0^{r}\) and hence of \(L_0^{r}\); (vii) \(\psi\) is a
sentence of \(L_0\).

\emph{Converse.}  Condition~4 of
\cref{def:p02-w02c-p02-003-general-reading}.
\end{proof}
\end{annexitem}

\begin{annexitem}{annex:p02-w02c-p02-003-b74-lindstrom}{thm:p02-w02c-p02-003-b74-lindstrom}{Lindstr\"om's theorem for logics in Barwise's sense}
\begin{proof}
Throughout, \(\mathcal L^{\ast}\) is a logic on \(\mathcal C\) with
conditions 4--8 of \cref{def:p02-w02c-p02-003-b74-logic}, and the numbers
(i)--(v) refer to the clauses of part~4 of
\cref{def:p02-w02c-p02-003-b74-languages}.  By the remark after that
definition, clauses (iii) and (i) give, for any \(K\) in \(\mathcal C\) and
any finite list of kinds and arities, new symbols of those kinds and
arities, none in \(K\), whose adjunction to \(K\) is a language of
\(\mathcal C\); ``new'' always refers to this.

\emph{Step 1: three consequences of the axioms.}
(1a) If \(\varphi\in L^{\ast}\) and \(L_\varphi\subseteq K\) with \(K\) in
\(\mathcal C\), then \(\varphi\in K^{\ast}\): the inclusion
\(i:L_\varphi\subseteq K\) is a morphism by (ii), and by the Occurrence Axiom
\(i^{\ast}\) is the inclusion of \(L_\varphi^{\ast}\) in \(K^{\ast}\).
(1b) If \(\varphi\in K^{\ast}\) and \(\mathfrak M\) is a \(K\)-structure,
then \(\mathfrak M\models\varphi\) iff
\(\mathfrak M\restriction L_\varphi\models\varphi\).  Indeed
\(\varphi^{i}=i^{\ast}(\varphi)=\varphi\) for the inclusion \(i\) of (1a);
\(i\) is simple, so every \(K\)-structure is \(i\)-invertible with
pull-back its reduct to \(L_\varphi\); and the Translation Axiom gives the
claim.  So two structures for languages containing \(L_\varphi\) that agree
on \(L_\varphi\) agree on \(\varphi\).
(1c) If \(\alpha:L\to K\) is a morphism of \(\mathcal C\),
\(\varphi\in L^{\ast}\) and the \(K\)-structure \(\mathfrak N\) is
\(\alpha\)-invertible, then \(\mathfrak N\models\varphi^{\alpha}\) iff
\(\mathfrak N^{-\alpha}\models\varphi\).  Let \(i:L_\varphi\subseteq L\) and
\(\beta=\alpha i\), the restriction of \(\alpha\) to \(L_\varphi\) with the
same terms.  By functoriality and the Occurrence Axiom,
\(\varphi^{\beta}=\alpha^{\ast}(i^{\ast}(\varphi))=\alpha^{\ast}(\varphi)=\varphi^{\alpha}\).
Since \(\mathbf V^{\beta}=\mathbf V^{\alpha}\) and the terms agree, the
pull-backs along \(\beta\) and \(\alpha\) have the same universe, and the
interpretation of each symbol of \(L_\varphi\) is the same in both; so
\(\mathfrak N^{-\beta}=\mathfrak N^{-\alpha}\restriction L_\varphi\), and it
is an \(L_\varphi\)-structure because \(\mathfrak N^{-\alpha}\) is an
\(L\)-structure.  The Translation Axiom for \(\beta\) and (1b) give
\(\mathfrak N\models\varphi^{\alpha}\) iff
\(\mathfrak N^{-\alpha}\restriction L_\varphi\models\varphi\) iff
\(\mathfrak N^{-\alpha}\models\varphi\).

\emph{Step 2: first-order logic inside \(\mathcal L^{\ast}\).}  Claim: let
\(K\) be in \(\mathcal C\), let \(\chi\) be a first-order formula of \(K\),
let \(x_1,\ldots,x_n\) be distinct variables including the free variables
of \(\chi\), and let \(d_1,\ldots,d_n\) be distinct constant symbols not in
\(K\) with \(K(\bar d)=K\cup\{d_1,\ldots,d_n\}\) in \(\mathcal C\).  Then
there is \(\theta\in K(\bar d)^{\ast}\) such that for every \(K\)-structure
\(\mathfrak M\) and \(\bar a\) in its universe,
\((\mathfrak M,\bar a)\models\theta\) iff \(\mathfrak M\models\chi[\bar a]\),
where \((\mathfrak M,\bar a)\) interprets \(d_i\) by \(a_i\) and
\(\mathfrak M\models\chi[\bar a]\) is first-order satisfaction at the
assignment \(x_i\mapsto a_i\).  We argue by induction on \(\chi\).  If
\(\chi\) is atomic, let \(\theta\) be \(\chi\) with \(d_i\) put for
\(x_i\); it is an atomic sentence of \(K(\bar d)\), in \(K(\bar d)^{\ast}\)
by condition~4, and an induction on terms shows that each term takes in
\((\mathfrak M,\bar a)\) the value it takes in \(\mathfrak M\) at the
assignment.  For \(\neg\chi_1\) and \(\chi_1\land\chi_2\) use condition~5;
put \(\neg(\neg\theta_1\land\neg\theta_2)\) for \(\chi_1\lor\chi_2\), and
treat \(\forall y\,\chi_1\) as \(\neg\exists y\,\neg\chi_1\).  Let
\(\chi=\exists y\,\chi_1\).  If \(y\) is \(x_i\), let \(\theta_1\) serve
\(\chi_1\) with the same list and put \(\theta=q(d_i,\theta_1)\in
K(\bar d)^{\ast}\); by condition~6, \((\mathfrak M,\bar a)\models\theta\) iff
some \(b\) gives \((\mathfrak M,\bar a')\models\theta_1\), where \(\bar a'\)
is \(\bar a\) with \(b\) in place \(i\), iff some \(b\) gives
\(\mathfrak M\models\chi_1[\bar a']\), iff \(\mathfrak M\models\chi[\bar a]\).
If \(y\) is not in the list, take a new constant symbol \(e\) with
\(K(\bar d,e)\) in \(\mathcal C\), let \(\theta_1\in K(\bar d,e)^{\ast}\)
serve \(\chi_1\) with the list \(x_1,\ldots,x_n,y\), and put
\(\theta=q(e,\theta_1)\in K(\bar d,e)^{\ast}\).  Its occurrence language
lies inside \(K(\bar d,e)\) and does not contain \(e\), so it lies inside
\(K(\bar d)\), and \(\theta\in K(\bar d)^{\ast}\) by (1a).  By condition~6,
\((\mathfrak M,\bar a)\models\theta\) iff some \(b\) gives
\((\mathfrak M,\bar a,b)\models\theta_1\), iff \(\mathfrak M\models\chi[\bar a]\).
No case uses an element of the universe, so the claim holds for empty
\(\mathfrak M\) when \(n=0\).  With \(n=0\) the claim gives, for each
first-order sentence \(\chi\) of \(K\), a sentence \(\chi^{\bullet}\in
K^{\ast}\) with the same \(K\)-structures as models.  Below, a first-order
sentence written as a conjunct of an \(\mathcal L^{\ast}\)-sentence stands
for this \(\chi^{\bullet}\).

\emph{Step 3: rank types of relational images.}  Fix the
\(\mathcal L^{\ast}\)-sentence \(\varphi\) and put \(L=L_\varphi\), finite by
condition~7.  Let \(v_1,v_2,\ldots\) be the variables in a fixed order.  For
nonempty \(L\)-structures \(\mathfrak A,\mathfrak B\), \(\bar a\in A^{n}\) and
\(\bar b\in B^{n}\) define \((\mathfrak A,\bar a)\sim_m(\mathfrak B,\bar b)\)
by recursion on \(m\): \(\sim_0\) holds when \(\mathfrak A^{r}\models
\gamma[\bar a]\) iff \(\mathfrak B^{r}\models\gamma[\bar b]\) for every
atomic formula \(\gamma\) of \(L^{r}\) with variables among
\(v_1,\ldots,v_n\) (read at \(v_i\mapsto a_i\), \(v_i\mapsto b_i\)); and
\(\sim_{m+1}\) holds when \(\sim_0\) holds, for every \(a\in A\) there is
\(b\in B\) with \((\mathfrak A,\bar aa)\sim_m(\mathfrak B,\bar bb)\), and for
every \(b\in B\) there is \(a\in A\) with the same.  An induction on \(m\)
shows that each \(\sim_m\) is an equivalence relation and that
\(\sim_{m+1}\) implies \(\sim_m\).
\begin{enumerate}
\item[(3a)] \emph{Finitely many classes, each defined by a formula.}  For
      each \(m,n\) there are finitely many \(\sim_m\)-classes of pairs
      \((\mathfrak A,\bar a)\) with \(\bar a\) of length \(n\), and each class
      \(\kappa\) has a first-order formula \(\xi_\kappa\) of \(L^{r}\), of
      quantifier rank at most \(m\) and with free variables among
      \(v_1,\ldots,v_n\), such that \(\mathfrak B^{r}\models\xi_\kappa[\bar
      b]\) iff \((\mathfrak B,\bar b)\in\kappa\).  For \(m=0\): \(L^{r}\) is
      finite and has only variables as terms, so there are finitely many
      atomic formulas with variables among \(v_1,\ldots,v_n\); a class is
      fixed by the set of those true at the tuple, and \(\xi_\kappa\) is the
      conjunction of those atomic formulas and the negations of the others
      (for \(n=0\) there are none, and \(\xi_\kappa\) is
      \(\forall v_1\,v_1=v_1\)).  For \(m+1\): the class of
      \((\mathfrak A,\bar a)\) under \(\sim_{m+1}\) is fixed by its
      \(\sim_0\)-class and the set
      \(X=\{[(\mathfrak A,\bar aa)]_m:a\in A\}\) of \(\sim_m\)-classes of
      length \(n+1\), a subset of a finite set; so there are finitely many
      classes.  Put
      \[
       \xi_\kappa=\xi_{\kappa_0}\land\bigwedge_{\lambda\in X}\exists
       v_{n+1}\,\xi_\lambda\land\forall v_{n+1}\bigvee_{\lambda\in X}\xi_\lambda ,
      \]
      a finite formula of rank at most \(m+1\) (\(X\) is nonempty because
      \(A\) is).  If \((\mathfrak B,\bar b)\in\kappa\) the three conjuncts
      hold at \(\bar b\) by the definition of \(\sim_{m+1}\) and the
      induction hypothesis.  Conversely, if \(\mathfrak B^{r}\models\xi_\kappa
      [\bar b]\), the first conjunct gives \(\sim_0\); the second gives, for
      each \(\lambda\in X\), some \(b\) with \((\mathfrak B,\bar bb)\in
      \lambda\), so every \(a\in A\) is matched; the third gives, for each
      \(b\), some \(\lambda\in X\) containing \((\mathfrak B,\bar bb)\), and
      an \(a\) with \((\mathfrak A,\bar aa)\in\lambda\).  The coincidence lemma
      \cref{lem:fol-coincidence} justifies reading each \(\xi_\lambda\) at the
      extended assignment.
\item[(3b)] \emph{Partial isomorphisms.}  A \emph{finite partial
      isomorphism} from \(\mathfrak A\) to \(\mathfrak B\) is a one-one map
      \(p\) from a finite subset of \(A\) into \(B\) such that, for all
      elements of its domain, \(R^{\mathfrak A}(\bar a)\) iff
      \(R^{\mathfrak B}(p\bar a)\) for relation symbols \(R\in L\),
      \(f^{\mathfrak A}(\bar a)=a\) iff \(f^{\mathfrak B}(p\bar a)=p(a)\) for
      function symbols, and \(c^{\mathfrak A}=a\) iff \(c^{\mathfrak B}=p(a)\)
      for constant symbols.  If \((\mathfrak A,\bar a)\sim_0(\mathfrak
      B,\bar b)\), then \(\{(a_i,b_i):i\leq n\}\) is a finite partial
      isomorphism: the atomic formulas \(v_i=v_j\) make it a one-one map,
      and \(R(\ldots)\), \(R_f(\ldots)\), \(P_c(v_i)\) give the three
      conditions, by the definition of \(\mathfrak A^{r}\).  A
      \emph{back-and-forth family} from \(\mathfrak A\) to \(\mathfrak B\) is
      a nonempty set \(\mathcal I\) of finite partial isomorphisms such that
      for each \(p\in\mathcal I\) and \(a\in A\) some \(q\in\mathcal I\)
      extends \(p\) with \(a\) in its domain, and for each \(b\in B\) some
      \(q\in\mathcal I\) extends \(p\) with \(b\) in its range.  If
      \(\mathcal I\) is one, so is the set \(\mathcal I^{\downarrow}\) of
      restrictions of its members, and a restriction of a finite partial
      isomorphism is one.
\end{enumerate}

\emph{Step 4: a back-and-forth family preserves \(\varphi\).}  Claim: if
\(\mathfrak A,\mathfrak B\) are nonempty \(L\)-structures with a
back-and-forth family \(\mathcal I\), then \(\mathfrak A\models\varphi\) iff
\(\mathfrak B\models\varphi\).  Suppose \(\mathfrak A\models\varphi\) and
\(\mathfrak B\not\models\varphi\).  By (iii) and (i) take a copy
\(L'=\{s':s\in L\}\) of \(L\) in \(\mathcal C\) meeting \(L\) only in
\(\mathbf V=\mathbf V'\), symbols of the same kinds and arities, and new
symbols: unary \(U,W,Q_U,Q_W\), constants \(e_U,e_W\), ternary
\(A_U,A_W\), a \(4\)-ary \(H\) and a binary \(I\).  Let \(K_0\) be the union,
in \(\mathcal C\).  Let \(\delta_U,\delta_W:L\to K_0\) be the \(0\)-term
interpretations with \(\mathbf V^{\delta_U}=U\), \(s^{\delta_U}=s\) and
\(\mathbf V^{\delta_W}=W\), \(s^{\delta_W}=s'\) for \(s\neq\mathbf V\).
Neither sends a symbol to \(\mathbf V\), and \(L\) is finite, so both are
morphisms of \(\mathcal C\) by (iv).  For a \(K_0\)-structure \(\mathfrak
N\), the pull-back \(\mathfrak N^{-\delta_U}\) has universe \(U^{\mathfrak
N}\) and the symbols of \(L\) restricted to it, and \(\mathfrak
N^{-\delta_W}\) has universe \(W^{\mathfrak N}\) and interprets each
\(s\in L\) by the restriction of \(s'^{\mathfrak N}\).  Write \(h_{s,t}\)
for \(\{(x,y):H(s,t,x,y)\}\).  Consider the first-order sentences of
\(K_0\):
\begin{itemize}
\item \(\mathrm{Srt}\): \(U,W,Q_U,Q_W\) are pairwise disjoint,
      \(\exists x\,U(x)\) and \(\exists x\,W(x)\);
\item \(\mathrm{Clo}\): for each \(n\)-ary function symbol \(f\in L\),
      \(\forall\bar x\,(\bigwedge_jU(x_j)\to U(f(\bar x)))\) and
      \(\forall\bar x\,(\bigwedge_jW(x_j)\to W(f'(\bar x)))\); for each
      constant symbol \(c\in L\), \(U(c)\land W(c')\);
\item \(\mathrm{Sq}\): \(Q_U(e_U)\), \(Q_W(e_W)\),
      \(\forall x\,y\,\neg H(e_U,e_W,x,y)\); \(A_U(s,a,s')\to Q_U(s)\land
      U(a)\land Q_U(s')\) and \(Q_U(s)\land U(a)\to\exists
      s'\,A_U(s,a,s')\), universally closed, and the same for \(A_W\), \(W\),
      \(Q_W\); and the update axiom
      \[
      \begin{aligned}
       &\forall s\,t\,a\,b\,s'\,t'\,\bigl(A_U(s,a,s')\land A_W(t,b,t')\to{}\\
       &\qquad\forall x\,y\,(H(s',t',x,y)\leftrightarrow(H(s,t,x,y)\lor(x=a\land
       y=b)))\bigr);
      \end{aligned}
      \]
\item \(\mathrm{Piso}(s,t)\), a formula: \(H(s,t,x,y)\to U(x)\land W(y)\);
      \(h_{s,t}\) is functional and one-one; for each relation symbol
      \(R\in L\) other than \(\mathbf V\), \(\bigwedge_jH(s,t,x_j,y_j)\to
      (R(\bar x)\leftrightarrow R'(\bar y))\); for each \(n\)-ary function
      symbol \(f\),
      \[\textstyle\bigwedge_jH(s,t,x_j,y_j)\land H(s,t,x,y)\to
      (f(\bar x)=x\leftrightarrow f'(\bar y)=y);\]
      for each constant symbol
      \(c\), \(H(s,t,x,y)\to(c=x\leftrightarrow c'=y)\); all universally
      closed except in \(s,t\);
\item \(\mathrm{BF}\): \(I(e_U,e_W)\); \(I(s,t)\to Q_U(s)\land Q_W(t)\land
      \mathrm{Piso}(s,t)\); \(I(s,t)\land U(a)\to\exists b\,s'\,t'\,(W(b)\land
      A_U(s,a,s')\land A_W(t,b,t')\land I(s',t'))\); and
      \(I(s,t)\land W(b)\to\exists a\,s'\,t'\,(U(a)\land A_U(s,a,s')\land
      A_W(t,b,t')\land I(s',t'))\), universally closed.
\end{itemize}
Let \(\Theta_0\) be the conjunction of these with \(\varphi^{\delta_U}\)
and \(\neg(\varphi^{\delta_W})\); \(\varphi^{\delta_U},\varphi^{\delta_W}\in
K_0^{\ast}\), so \(\Theta_0\in K_0^{\ast}\) by condition~5.

\(\Theta_0\) has a model \(\mathfrak N\).  Its universe is the disjoint
union of \(A\), a copy of \(B\) (identified with \(B\)), and the finite
sequences from \(A\) and from \(B\) as new objects; \(\mathbf V\) is the
universe.  Put \(U=A\), \(W=B\), \(Q_U\), \(Q_W\) the two sets of
sequences, \(e_U,e_W\) the empty ones, \(A_U(s,a,s')\) iff \(s'\) is \(s\)
with \(a\) appended, and likewise \(A_W\).  Let \(H(s,t,x,y)\) hold iff
there are \(i<\lvert s\rvert\) and \(j<\lvert t\rvert\) with
\(\lvert s\rvert-i=\lvert t\rvert-j\), \(s(i)=x\) and \(t(j)=y\): the two
sequences are aligned at their ends, so appending \(a\) and \(b\) adds
exactly the pair \((a,b)\), and the update axiom holds for all \(s,t\);
when \(\lvert s\rvert=\lvert t\rvert\), \(h_{s,t}\) is
\(\{(s(i),t(i))\}\).  Let \(I(s,t)\) hold iff \(\lvert s\rvert=\lvert
t\rvert\) and \(h_{s,t}\) is a function belonging to
\(\mathcal I^{\downarrow}\).  Interpret each relation symbol \(R\in L\) by
\(R^{\mathfrak A}\) and \(R'\) by \(R^{\mathfrak B}\); fix \(a_0\in A\),
\(b_0\in B\), and let \(f(\bar x)=f^{\mathfrak A}(\bar x)\) for \(\bar x\)
in \(A\) and \(a_0\) otherwise, \(f'(\bar y)=f^{\mathfrak B}(\bar y)\) for
\(\bar y\) in \(B\) and \(b_0\) otherwise, \(c=c^{\mathfrak A}\),
\(c'=c^{\mathfrak B}\).  Then \(\mathrm{Srt}\), \(\mathrm{Clo}\) and
\(\mathrm{Sq}\) hold.  \(\mathrm{BF}\) holds: the empty map is in
\(\mathcal I^{\downarrow}\); members of \(\mathcal I^{\downarrow}\) are
finite partial isomorphisms, which is \(\mathrm{Piso}\); and if
\(h_{s,t}\in\mathcal I^{\downarrow}\) and \(a\in A\), take \(p\in\mathcal
I\) extending \(h_{s,t}\), \(q\in\mathcal I\) extending \(p\) with \(a\) in
its domain, \(b=q(a)\), and append \(a\), \(b\): the new map
\(h_{s,t}\cup\{(a,b)\}\) is a restriction of \(q\); the back clause is
symmetric.  By \(\mathrm{Clo}\), \(\mathfrak N\) is \(\delta_U\)- and
\(\delta_W\)-invertible, with \(\mathfrak N^{-\delta_U}=\mathfrak A\) and
\(\mathfrak N^{-\delta_W}=\mathfrak B\); by (1c),
\(\mathfrak N\models\varphi^{\delta_U}\) and \(\mathfrak
N\not\models\varphi^{\delta_W}\).

By the L\"owenheim property, \(\Theta_0\) has a model \(\mathfrak N^{\ast}\)
of cardinality at most \(\aleph_0\).  By \(\mathrm{Clo}\) and (1c),
\(\mathfrak A^{\ast}=\mathfrak N^{\ast-\delta_U}\models\varphi\) and
\(\mathfrak B^{\ast}=\mathfrak N^{\ast-\delta_W}\not\models\varphi\), and by
\(\mathrm{Srt}\) both are nonempty and countable.  Enumerate
\(U^{\mathfrak N^{\ast}}=\{u_0,u_1,\ldots\}\) and
\(W^{\mathfrak N^{\ast}}=\{w_0,w_1,\ldots\}\), with repetitions if finite.
Put \((s_0,t_0)=(e_U,e_W)\); given \(I(s_k,t_k)\), use the forth clause
with \(u_j\) when \(k=2j\) and the back clause with \(w_j\) when
\(k=2j+1\) to get \(s_{k+1},t_{k+1}\) with \(I(s_{k+1},t_{k+1})\), paired
appends of a pair \((a,b)\).  By \(\mathrm{Sq}\), \(h_{s_0,t_0}=\varnothing\)
and \(h_{s_{k+1},t_{k+1}}=h_{s_k,t_k}\cup\{(a,b)\}\), so the maps increase;
their union \(h\) contains every \(u_j\) in its domain and every \(w_j\) in
its range, and by \(\mathrm{Piso}\) it is a one-one map of
\(U^{\mathfrak N^{\ast}}\) onto \(W^{\mathfrak N^{\ast}}\).  Any finitely
many elements lie in one \(h_{s_k,t_k}\), so \(\mathrm{Piso}\) gives
\(R^{\mathfrak A^{\ast}}(\bar a)\leftrightarrow R^{\mathfrak B^{\ast}}(h\bar a)\),
\(f^{\mathfrak B^{\ast}}(h\bar a)=h(f^{\mathfrak A^{\ast}}(\bar a))\) (take
\(x=f^{\mathfrak A^{\ast}}(\bar a)\), which lies in \(U\) by
\(\mathrm{Clo}\)), and \(c^{\mathfrak B^{\ast}}=h(c^{\mathfrak A^{\ast}})\).
So \(h\) is an isomorphism of \(\mathfrak A^{\ast}\) onto \(\mathfrak
B^{\ast}\), and the Isomorphism Axiom contradicts \(\mathfrak
A^{\ast}\models\varphi\), \(\mathfrak B^{\ast}\not\models\varphi\).

\emph{Step 5: the case of a bounded rank.}  Say that \(m\) \emph{bounds}
\(\varphi\) if \(\mathfrak A\sim_m\mathfrak B\) (with \(n=0\)) implies that
\(\mathfrak A\models\varphi\) iff \(\mathfrak B\models\varphi\), for all
nonempty \(L\)-structures.  If \(m\) bounds \(\varphi\), let \(\psi_0\) be
the disjunction of the formulas \(\xi_\kappa\) of (3a) over the finitely
many \(\sim_m\)-classes \(\kappa\) of nonempty structures that contain a
model of \(\varphi\) (the sentence \(\neg\forall v_1\,v_1=v_1\) if there
is none).  For nonempty \(\mathfrak M\), \(\mathfrak M^{r}\models\psi_0\) iff
\(\mathfrak M\) is \(\sim_m\)-equivalent to a model of \(\varphi\), iff
\(\mathfrak M\models\varphi\).  By part~1 of
\cref{thm:p02-w02c-p02-003-translations}, \(\psi_0^{-}\) is a first-order
sentence of \(L\) with \(\mathfrak M\models\psi_0^{-}\) iff
\(\mathfrak M^{r}\models\psi_0\).  At most one \(L\)-structure is empty, and
only when \(L\) has no constant symbol.  Let \(\varepsilon\) be
\(\forall v_1\,v_1=v_1\) if it exists and satisfies \(\varphi\), and
\(\neg\forall v_1\,v_1=v_1\) otherwise, and put
\[
 \psi=(\exists v_1\,v_1=v_1\land\psi_0^{-})\lor(\neg\exists v_1\,v_1=v_1\land\varepsilon).
\]
Then \(\psi\) is a first-order sentence of \(L=L_\varphi\), and
\(\mathfrak M\models\varphi\) iff \(\mathfrak M\models\psi\) for every
\(L\)-structure \(\mathfrak M\).  This is part~1 of the theorem.  It
remains to show that some \(m\) bounds \(\varphi\); suppose not, and for
each \(m\) choose nonempty \(\mathfrak A_m\sim_m\mathfrak B_m\) with
\(\mathfrak A_m\models\varphi\) and \(\mathfrak B_m\not\models\varphi\).

\emph{Step 6: the coding sentence \(\Theta\).}  Take \(L'\), \(\delta_U\),
\(\delta_W\) and the symbols of Step~4 other than \(I\), and further new
symbols: unary \(N\), binary \(<\) and \(S\), constants \(z_0,z_\ast\), and a
ternary \(G\); let \(K_1\) be the union, in \(\mathcal C\), with
\(\delta_U,\delta_W:L\to K_1\) as before.  Let \(\mathrm{Disc}\) say:
\(<\) is a strict linear order of \(N\) (irreflexive, transitive, total on
\(N\), and holding only between elements of \(N\)); \(N(z_0)\) and no
element of \(N\) is below \(z_0\); \(S(r,r')\) iff \(N(r)\), \(N(r')\),
\(r<r'\) and no element of \(N\) lies strictly between them; every element
of \(N\) other than \(z_0\) has some \(r'\) with \(S(r',r)\).  Let
\(\Theta\) be the conjunction of \(\mathrm{Srt}\), \(\mathrm{Clo}\),
\(\mathrm{Sq}\), \(\mathrm{Disc}\), \(N(z_\ast)\),
\(\forall r\,(N(r)\to r<z_\ast\lor r=z_\ast)\), \(\varphi^{\delta_U}\),
\(\neg(\varphi^{\delta_W})\), and
\[
\begin{aligned}
 &G(z_\ast,e_U,e_W),\\
 &\forall r\,s\,t\,\bigl(G(r,s,t)\to N(r)\land Q_U(s)\land Q_W(t)\land\mathrm{Piso}(s,t)\bigr),\\
 &\forall r\,r'\,s\,t\,\bigl(S(r,r')\land G(r',s,t)\to{}\\
 &\qquad\forall a\,(U(a)\to\exists b\,s'\,t'\,(W(b)\land A_U(s,a,s')\land A_W(t,b,t')\land G(r,s',t')))\\
 &\qquad\land\forall b\,(W(b)\to\exists a\,s'\,t'\,(U(a)\land A_U(s,a,s')\land A_W(t,b,t')\land G(r,s',t')))\bigr).
\end{aligned}
\]
So \(\Theta\in K_1^{\ast}\).  \emph{Models of \(\Theta\).}  For \(m\geq1\)
let \(\mathfrak N_m\) be built from \(\mathfrak A_m,\mathfrak B_m\) as in
Step~4, without \(I\), with \(N=\{0,\ldots,m\}\) as new objects, \(<\) and
\(S\) the order and successor relation of these numbers, \(z_0=0\),
\(z_\ast=m\), and \(G(r,s,t)\) iff \(r\in N\), \(\lvert s\rvert=\lvert
t\rvert\) and \((\mathfrak A_m,s)\sim_r(\mathfrak B_m,t)\), the sequences
read as tuples.  \(G(z_\ast,e_U,e_W)\) is \(\mathfrak A_m\sim_m\mathfrak
B_m\); \(G\) implies \(\sim_0\), hence \(\mathrm{Piso}\) by (3b); and if
\(S(r,r')\), so \(r'=r+1\), and \(G(r',s,t)\), the definition of
\(\sim_{r+1}\) supplies the forth and back extensions with \(G(r,s',t')\).
So \(\mathfrak N_m\models\Theta\).  \emph{A descending sequence gives a
contradiction.}  Suppose \(\mathfrak N\models\Theta\) and there are
\(r_0=z_\ast,r_1,r_2,\ldots\) in \(N^{\mathfrak N}\) with
\(S(r_{k+1},r_k)\) for every \(k\).  Let \(\mathcal I_k\) be the set of maps
\(h_{s,t}\) with \(G(r_k,s,t)\) such that \((s,t)\) is reached from
\((e_U,e_W)\) by \(k\) paired appends, and \(\mathcal I=\bigcup_k\mathcal
I_k\).  Then \(\varnothing=h_{e_U,e_W}\in\mathcal I_0\); by \(\mathrm{Sq}\)
each member of \(\mathcal I_k\) has at most \(k\) pairs, and by
\(\mathrm{Piso}\) it is a finite partial isomorphism from \(\mathfrak
A^{\ast}=\mathfrak N^{-\delta_U}\) to \(\mathfrak B^{\ast}=\mathfrak
N^{-\delta_W}\) (both \(L\)-structures by \(\mathrm{Clo}\), nonempty by
\(\mathrm{Srt}\)); and if \(h_{s,t}\in\mathcal I_k\) and \(a\in U\), the
last conjunct with \(S(r_{k+1},r_k)\) gives \(b,s',t'\) with
\(G(r_{k+1},s',t')\), and \(h_{s',t'}=h_{s,t}\cup\{(a,b)\}\in\mathcal
I_{k+1}\) by the update axiom; the back clause is symmetric.  So
\(\mathcal I\) is a back-and-forth family, and Step~4 gives
\(\mathfrak A^{\ast}\models\varphi\) iff \(\mathfrak B^{\ast}\models\varphi\),
while by \(\mathrm{Clo}\) and (1c) \(\mathfrak A^{\ast}\models\varphi\) and
\(\mathfrak B^{\ast}\not\models\varphi\).  Hence no model of \(\Theta\) has
such a sequence.

\emph{Step 7: hypothesis (b).}  Let \(\eta_k\) be the first-order sentence
\(\exists x_0\cdots\exists x_k\,(x_0=z_0\land\bigwedge_{i<k}S(x_i,x_{i+1})
\land x_k<z_\ast)\) of \(K_1\), and \(\Gamma=\{\Theta\}\cup\{\eta_k:k<\omega\}\),
a countable subset of \(K_1^{\ast}\).  A finite subset of \(\Gamma\) whose
indices are below \(m\) holds in \(\mathfrak N_m\), where the chain
\(0,1,\ldots,k\) witnesses \(\eta_k\).  By (b), \(\Gamma\) has a model
\(\mathfrak N\).  In \(\mathfrak N\) an element has at most one immediate
successor and at most one immediate predecessor, since \(<\) is linear; so
\(\eta_k\) fixes elements \(n_0=z_0,n_1,\ldots,n_k\), the same for all
\(k\geq j\) at each \(n_j\), with \(S(n_j,n_{j+1})\) and \(n_k<z_\ast\).  Put
\(r_0=z_\ast\).  If \(r_k\) is defined and \(n_j<r_k\) for all \(j\), then
\(r_k\neq z_0\), so \(\mathrm{Disc}\) gives an immediate predecessor
\(r_{k+1}\), with \(S(r_{k+1},r_k)\); for each \(j\), \(n_{j+1}<r_k\) and
immediacy give \(n_{j+1}\leq r_{k+1}\), so \(n_j<r_{k+1}\).  This defines
the sequence of Step~6 in a model of \(\Theta\), a contradiction.

\emph{Step 8: hypothesis (a).}  First, \emph{every model \(\mathfrak N\) of
\(\Theta\) has \(N^{\mathfrak N}\) finite}.  Put \(r_0=z_\ast\) and, while
\(r_k\neq z_0\), let \(r_{k+1}\) be the immediate predecessor of \(r_k\).
By Step~6 this stops at some \(r_k=z_0\).  Then \(N^{\mathfrak
N}=\{r_0,\ldots,r_k\}\): an element \(x\in N\) outside this set has
\(z_0<x<z_\ast\), and for the least \(j\) with \(r_j<x\), which exists and
is at least \(1\), \(r_j<x<r_{j-1}\), against \(S(r_j,r_{j-1})\).

Next, \emph{countable choices}.  For each \(m\geq1\), let \(\kappa\) be the
\(\sim_m\)-class of \(\mathfrak A_m\) and \(\xi_\kappa\) its sentence from
(3a).  The \(\mathcal L^{\ast}\)-sentence \(\varphi\land(\xi_\kappa^{-})^{\bullet}
\land(\exists v_1\,v_1=v_1)^{\bullet}\) of \(L\) holds in \(\mathfrak A_m\), so
by the L\"owenheim property it has a model \(\mathfrak A_m'\) of cardinality
at most \(\aleph_0\); it is nonempty, satisfies \(\varphi\), and
\(\mathfrak A_m'\sim_m\mathfrak A_m\) by part~1 of
\cref{thm:p02-w02c-p02-003-translations} and (3a).  Likewise
\(\neg\varphi\land(\xi_\kappa^{-})^{\bullet}\land(\exists v_1\,v_1=v_1)^{\bullet}\)
holds in \(\mathfrak B_m\) (the class of \(\mathfrak B_m\) is \(\kappa\)) and
gives a countable nonempty \(\mathfrak B_m'\sim_m\mathfrak B_m\) with
\(\mathfrak B_m'\not\models\varphi\).  So \(\mathfrak A_m'\sim_m\mathfrak
B_m'\), and the structure \(\mathfrak N_m'\) built from them as in Step~6
is a countable model of \(\Theta\); choose its universe so that
\(N=\{0,\ldots,m\}\) consists of the natural numbers themselves and the
rest of the universe is new, distinct for distinct \(m\).  The axiom of
choice is used to choose \(\mathfrak A_m',\mathfrak B_m'\) for all \(m\).

\emph{The indexed union.}  By (v) there are a constant symbol \(d\notin
K_1\), a language \(K_2\) in \(\mathcal C\) containing \(d\), and a one-one
\(1\)-term interpretation \(\alpha:K_1\to K_2\) with term \(d\) that is a
morphism of \(\mathcal C\).  So \(\mathbf V^{\alpha}\) is binary, each
\(n\)-ary relation or function symbol \(s\) of \(K_1\) goes to an
\((n+1)\)-ary \(s^{\alpha}\), and each constant symbol goes to a unary
function symbol.  For a \(K_2\)-structure \(\mathfrak M\) the pull-back has
universe \(\{x:\mathbf V^{\alpha}(x,d)\}\).  Let \(\mathrm{Inv}\) be the
first-order sentence of \(K_2\) saying that this set is closed under each
\(f^{\alpha}(\cdot,d)\) and contains each \(c^{\alpha}(d)\), for the finitely
many function and constant symbols of \(K_1\); a \(K_2\)-structure is
\(\alpha\)-invertible iff it satisfies \(\mathrm{Inv}\).  Add new symbols: a
unary \(J\), a binary \(\prec\), a unary function symbol \(g\) and a
constant symbol \(x_0\), giving \(K_3\) in \(\mathcal C\).  Let
\[
 \theta=\mathrm{Inv}\land\Theta^{\alpha}\land
 \forall z\,\bigl(z\prec x_0\to\mathbf V^{\alpha}(z,d)\land N^{\alpha}(z,d)\bigr)\in K_3^{\ast},
\]
using (1a) for \(\Theta^{\alpha}\in K_2^{\ast}\).  Let
\(\sigma=\neg(D(x_0)\land\neg q(d,\theta))\) and \(\lambda_2=\neg
q(x_0,\neg\sigma)\), both in \(K_3^{\ast}\), and let \(\lambda_1\) be the
first-order sentence of \(K_3\) saying that \(J\) is nonempty, \(\prec\) is a
strict linear order of \(J\) with no greatest element, \(g\) maps every
element into \(D\), and \(g\) is one-one.  Put
\(\Xi=\lambda_1\land\lambda_2\).  By conditions 5 and~6,
\(\mathfrak M\models\lambda_2\) iff for every \(x\in J^{\mathfrak M}\) there
is \(i\) with \((\mathfrak M,x,i)\models\theta\), interpreting \(x_0\) by
\(x\) and \(d\) by \(i\).

\emph{Every model of \(\Xi\) is countably infinite.}  Let
\(\mathfrak M\models\Xi\).  \(J\) is nonempty with no greatest element, so it
is infinite.  For \(x\in J\) take \(i\) with \((\mathfrak M,x,i)\models
\theta\).  By \(\mathrm{Inv}\) the \(K_2\)-reduct of \((\mathfrak M,x,i)\) is
\(\alpha\)-invertible; by (1b) and (1c) its pull-back \(\mathfrak P\)
satisfies \(\Theta\), so \(N^{\mathfrak P}\) is finite; and
\(N^{\mathfrak P}=\{z:\mathbf V^{\alpha}(z,i)\land N^{\alpha}(z,i)\}\)
contains every \(\prec\)-predecessor of \(x\).  So every element of \(J\) has
finitely many predecessors, the number of predecessors is a one-one map
of \(D\) into \(\omega\), \(J\) is countable, and \(g\) makes the universe
countable.

\emph{\(\Xi\) has an infinite model.}  Let the universe of \(\mathfrak M\)
be \(\omega\), the universes of all \(\mathfrak N_m'\), and new elements
\(\iota_m\) (\(m\geq1\)); it is countable, and fix a one-one \(g\) into
\(\omega\).  Put \(\mathbf V^{\alpha}(x,\iota_m)\) iff \(x\) is in the
universe of \(\mathfrak N_m'\); \(s^{\alpha}(\bar x,\iota_m)\) iff
\(s^{\mathfrak N_m'}(\bar x)\) for relation symbols \(s\) of \(K_1\);
\(f^{\alpha}(\bar x,\iota_m)=f^{\mathfrak N_m'}(\bar x)\) for \(\bar x\) in
\(\mathfrak N_m'\) and \(c^{\alpha}(\iota_m)=c^{\mathfrak N_m'}\); every other
value of a function is \(0\), every other instance of a relation is false,
and the symbols of \(K_2\) outside the range of \(\alpha\) other than
\(\mathbf V\) are interpreted by empty relations and the constant value
\(0\).  Let \(J=\omega\) with \(\prec\) its order, and \(x_0=0\), \(d=\iota_1\).
For each \(m\), \((\mathfrak M,\iota_m)\) (with \(d\) read as \(\iota_m\)) is
\(\alpha\)-invertible with pull-back \(\mathfrak N_m'\models\Theta\); and for
\(x=m\geq1\) in \(J\), \(i=\iota_m\) witnesses \(\theta\), since the
predecessors \(0,\ldots,m-1\) lie in \(N^{\mathfrak N_m'}=\{0,\ldots,m\}\);
for \(x=0\) take \(i=\iota_1\).  So \(\mathfrak M\models\Xi\) and is
infinite.  By (a), \(\Xi\) has an uncountable model, which contradicts the
previous paragraph.

So under (a) or (b) some \(m\) bounds \(\varphi\), and Step~5 gives part~1.

\emph{Part~2.}  Step~2 with \(n=0\) gives, for each first-order sentence
\(\psi\) of \(L\) in \(\mathcal C\), the sentence \(\theta=\psi^{\bullet}\in
L^{\ast}\) with the same \(L\)-structures as models.  When \(L\) is the
occurrence language of \(\psi\) in \(L_{\omega\omega}\), the occurrence
language \(L_\theta\) is the least language of \(\mathcal C\) containing
\(\theta\), so \(L_\theta\subseteq L\), and every symbol occurring in
\(\theta\) occurs in \(\psi\).  Part~1 gives the other half of
\(\mathcal L^{\ast}\equiv L_{\omega\omega}\), since \(\psi\) is a sentence of
\(L_\varphi\) and so its occurrence language in \(L_{\omega\omega}\) lies
inside \(L_\varphi\).
\end{proof}
\end{annexitem}

\begin{annexitem}{annex:p02-w02c-p02-003-b74-five-seven}{thm:p02-w02c-p02-003-b74-five-seven}{Barwise's 5.7 for logics in his sense}
\begin{proof}
A countable subset \(T\) of \(L^{\ast}\) whose finite subsets have models is
a set of \(\mathcal L^{\ast}\)-sentences of \(L\), so it has a model by the
Compactness Theorem as read in the statement; this is hypothesis~(b).  A
single sentence with a model is a set of at most \(\aleph_0\) sentences with
a model, so it has a model of cardinality at most \(\aleph_0\); this is
condition~8.  \Cref{thm:p02-w02c-p02-003-b74-lindstrom} applies.
\end{proof}
\end{annexitem}

\begin{annexitem}{annex:p02-w02c-p02-003-or-not}{thm:p02-w02c-p02-003-or-not}{Connectives from \(\lor\) and \(\neg\)}
\begin{proof}
By \cref{thm:p02-w02-p02-008-neg-and}, \(\langle\{\neg,\land\}\rangle\) is
the set \(\mathcal O\) of all operations of positive arity.  The
conjunction operation is generated by negation and disjunction, since
\(x\land y\) and \(\neg(\neg x\lor\neg y)\) take the same value at each of
the four arguments (both are \(1\) exactly at \((1,1)\)).  So
\(\{\neg,\land\}\subseteq\langle\{\neg,\lor\}\rangle\), and part~3 of
\cref{thm:p02-w02-p02-008-generated} gives
\(\mathcal O=\langle\{\neg,\land\}\rangle\subseteq\langle\{\neg,\lor\}\rangle\).
By \cref{thm:p02-w02-p02-008-formulas-clone}, every \(p\) of arity
\(n\geq1\) is represented by a formula \(\pi\) over \(\{\neg,\lor\}\) in
\(p_1,\ldots,p_n\).  For the second claim, fix \(\mathfrak M\), \(s\), and
let \(e_i\) be the truth value of \(\varphi_i\) at \(s\).  By induction on
subformulas \(\pi'\) of \(\pi\), the formula \(\pi'(\varphi_1,\ldots,
\varphi_n)\) is true in \(\mathfrak M\) at \(s\) iff \(\pi'\) is true at the
valuation \(p_i\mapsto e_i\): for \(p_i\) this is the definition of
\(e_i\), and the clauses for \(\neg\) and \(\lor\) are the same in both
semantics.  For \(\pi'=\pi\) the right side is \(p(e_1,\ldots,e_n)=1\),
which is the clause for \(\#\).
\end{proof}
\end{annexitem}

\begin{annexitem}{annex:p02-w02c-p02-003-small-model}{thm:p02-w02c-p02-003-small-model}{Small set-theoretic models}
\begin{proof}
Suppose \(T\entails\neg\forall x\,x=x\) in the sense of set-theoretic
structures.  By \cref{thm:fol-completeness}, available for every language
under the axiom of choice, \(T\derives\neg\forall x\,x=x\), so by the
hypothesis \(\neg\forall x\,x=x\) is true in \(\mathcal M\), which it is
not.  Hence some set-theoretic structure \(\mathfrak A\) satisfies \(T\) and
not \(\neg\forall x\,x=x\); in particular \(\mathfrak A\models T\).  By
\cref{thm:downward-ls} with \(A=\varnothing\) there is
\(\mathfrak B\preccurlyeq\mathfrak A\) whose universe has cardinality at
most \(\max(\lvert L\rvert,\aleph_0)\); as an elementary substructure,
\(\mathfrak B\) satisfies every sentence of \(T\).
\end{proof}
\end{annexitem}

\begin{annexitem}{annex:p02-w02c-p02-003-hf}{thm:p02-w02c-p02-003-hf}{Numbers, sequences, admissibility in \(\mathrm{HF}(\mathfrak M)\)}
\begin{proof}
Write \(\mathrm{HF}_n\) for \(\mathrm{HF}_n(M)\) and \(A=M\cup\mathrm{HF}(M)\).
By the reading of \cref{def:p02-w02-p02-005-bw-weak-so}, elements of \(M\)
have no members and are not elements of \(\mathrm{HF}(M)\).

\emph{Preliminaries.}  (P1) \(\mathrm{HF}_n\subseteq\mathrm{HF}_{n+1}\):
true for \(n=0\); if \(\mathrm{HF}_{n-1}\subseteq\mathrm{HF}_n\), every
finite subset of \(M\cup\mathrm{HF}_{n-1}\) is a finite subset of
\(M\cup\mathrm{HF}_n\).  (P2) Every finite \(F\subseteq A\) lies in
\(\mathrm{HF}(M)\): each member of \(F\) lies in \(M\cup\mathrm{HF}_{n}\)
for some \(n\), so by (P1) \(F\subseteq M\cup\mathrm{HF}_N\) for the
largest such \(n\), and \(F\in\mathrm{HF}_{N+1}\).  (P3) \(A\) is
transitive: members of elements of \(M\) do not exist, and if
\(a\in\mathrm{HF}_{n+1}\) then \(a\subseteq M\cup\mathrm{HF}_n\subseteq A\).
(P4) Put \(\rho(u)=0\) for \(u\in M\) and, for \(a\in\mathrm{HF}(M)\), let
\(\rho(a)\) be the least \(n\) with \(a\in\mathrm{HF}_n\); then
\(\rho(a)\geq1\), every member of \(a\) lies in
\(M\cup\mathrm{HF}_{\rho(a)-1}\), and so has \(\rho\)-value less than
\(\rho(a)\).  Hence induction on \(\rho\) proves statements about all
members of \(A\), and membership in \(A\) is well-founded.

\emph{Part~1.}  For the von Neumann numbers, \(0=\emptyset\) is a finite
subset of \(M\cup\mathrm{HF}_0\), so \(0\in\mathrm{HF}_1\); if
\(k\in\mathrm{HF}_{k+1}\) for all \(k<n\), then by (P1)
\(n=\{0,\ldots,n-1\}\) is a finite subset of \(\mathrm{HF}_n\), so
\(n\in\mathrm{HF}_{n+1}\).  The four conjuncts of \(\mathrm{Nat}(n)\) hold:
\(n\) has no member in \(M\); the second conjunct is then vacuous for its
members, which are numbers; \(c\in k\in n\) gives \(c<k<n\), so \(c\in n\);
and \(d\in c\in k\) gives \(d<c<k\), so \(d\in k\).

Conversely, suppose \(\mathrm{HF}(\mathfrak M)\models\mathrm{Nat}[a]\).  By
the first conjunct, every member of \(a\) is a set.  Let \(b\in a\).  Then
\(\mathrm{Nat}[b]\): if \(x\in b\) with \(x\in M\), the second conjunct
gives \(x\in a\), contrary to the first; if \(c\in b\) then \(c\in a\) by
the third conjunct, so \(c\) has no member in \(M\) and the second conjunct
for \(b\) is vacuous; the third conjunct for \(b\) is the fourth conjunct
for \(a\); and if \(c\in b\), \(d\in c\), \(e\in d\), then \(c\in a\) by the
third conjunct for \(a\), and the fourth conjunct for \(a\) applied to
\(c\in a\), \(d\in c\), \(e\in d\) gives \(e\in c\), which is the fourth
conjunct for \(b\).  By induction on \(\rho(a)\) (P4), each \(b\in a\) is
therefore a von Neumann number, so \(a\) is a finite set of numbers.  By
the third conjunct, \(k\in a\) and \(j<k\) give \(j\in a\).  A finite set
of numbers closed downward is \(\{0,\ldots,m-1\}=m\) for \(m\) its number
of elements.

\emph{Part~2.}  For \(i,v\in A\), the sets \(\{i\}\) and \(\{i,v\}\) are
finite subsets of \(A\), hence in \(\mathrm{HF}(M)\) by (P2), and so is
\(\langle i,v\rangle\); a finite set of such pairs is in
\(\mathrm{HF}(M)\) by (P2) again.  For the formulas, let
\begin{align*}
\mathrm{Sg}(e,i)\colon\;& i\in e\land\forall w\in e\,(w=i)
 \land\neg\exists y\,(y\in e),\\
\mathrm{Pair}(p,i,x)\colon\;&\exists e\in p\,\exists f\in p\,
 \bigl[\forall q\in p\,(q=e\lor q=f)\land\neg\exists y\,(y\in p)
 \land\mathrm{Sg}(e,i)\\
 &\land i\in f\land x\in f\land\forall w\in f\,(w=i)
 \land\forall y\,(y\in f\to y=x)\bigr],\\
\mathrm{Pair}'(p,i,b)\colon\;&\exists e\in p\,\exists f\in p\,
 \bigl[\forall q\in p\,(q=e\lor q=f)\land\neg\exists y\,(y\in p)
 \land\mathrm{Sg}(e,i)\\
 &\land i\in f\land b\in f\land\forall w\in f\,(w=i\lor w=b)
 \land\neg\exists y\,(y\in f)\bigr].
\end{align*}
Here \(x,y\) range over \(M\) and \(p,e,f,q,w,i,b\) over
\(\mathrm{HF}(M)\).  Since \(A\) is transitive (P3) and its sets are real
sets with the real membership relation, two sets of \(A\) with the same
members of \(A\) are equal.  So \(\mathrm{Pair}(p,i,x)\) holds iff
\(p=\{\{i\},\{i,x\}\}\), and \(\mathrm{Pair}'(p,i,b)\) iff
\(p=\{\{i\},\{i,b\}\}\), where \(\mathrm{Sg}(e,i)\) says \(e=\{i\}\); when \(b=i\) both members of \(p\) are \(\{i\}\)
and the formula still holds, with \(e=f\).  Let
\(\mathrm{P}(p,i)\) be \(\exists x\,\mathrm{Pair}(p,i,x)\lor
\exists b\,\mathrm{Pair}'(p,i,b)\), and let \(\mathrm{Seq}(s,n)\) be
\[
\begin{gathered}
\mathrm{Nat}(n)\land\neg\exists y\,(y\in s)
\land\forall p\in s\,\exists i\in n\,\mathrm{P}(p,i)
\land\forall i\in n\,\exists p\in s\,\mathrm{P}(p,i)\\
{}\land\forall p\in s\,\forall q\in s\,\forall i\in n\,
 (\mathrm{P}(p,i)\land\mathrm{P}(q,i)\to p=q).
\end{gathered}
\]
By part~1 and the reading of the pair formulas, \(\mathrm{Seq}(s,n)\)
holds iff \(n\) is a natural number, every member of \(s\) is a pair with
first coordinate in \(n\), every \(i\in n\) is the first coordinate of a
member of \(s\), and no two members share a first coordinate; that is, iff
\(s\) is a function with domain \(n\).  Put
\(\mathrm{Val}(s,i,x)=\exists p\in s\,\mathrm{Pair}(p,i,x)\) and
\(\mathrm{Val}'(s,i,b)=\exists p\in s\,\mathrm{Pair}'(p,i,b)\).  When
\(\mathrm{Seq}(s,n)\) holds, \(\langle i,x\rangle\in s\) iff \(i\in n\) and
\(s(i)=x\), and likewise for \(b\).

\emph{Part~3.}  Every element of \(A\) is in the universe \(V_M\) of sets
over the urelements \(M\), and \(A\) is transitive (P3).  We check the
axioms of \(\mathrm{KPU}\) in \((A,\in,M)\)
\cite[ch.~A.7, \S2, 2.3, p.~239]{Makkai1977}.  \emph{Urelements}: members
of \(M\) have no members.  \emph{Empty set}: \(\emptyset\in\mathrm{HF}_1\)
and \(\emptyset\notin M\).  \emph{Extensionality}: shown in part~2.
\emph{Foundation}: suppose \(\forall x\,(\forall y\in x\,\phi(y)\to\phi(x))\)
holds in \(A\) and \(\phi(x_0)\) fails for some \(x_0\in A\); take such an
\(x_0\) with \(\rho(x_0)\) least; by (P4) every \(y\in x_0\) has smaller
\(\rho\), so \(\phi(y)\) holds for all \(y\in x_0\), and then \(\phi(x_0)\)
holds, a contradiction.  \emph{Pair}: \(\{x,y\}\in A\) by (P2).
\emph{Union}: for \(a\in A\), the set of members of members of \(a\) is a
finite subset of \(A\) (empty if \(a\in M\)), so it lies in \(A\) by (P2).
\emph{Separation}: for \(a\in A\) and any formula \(\phi\), possibly with
parameters and with relation symbols for the relations of \(\mathfrak M\),
\(\{x\in a:A\models\phi[x]\}\) is a finite subset of \(A\), hence in \(A\).
\emph{Collection}: if every \(x\in a\) has some \(y\in A\) with
\(A\models\phi[x,y]\), choose one such \(y_x\) for each of the finitely many
members \(x\) of \(a\); then \(b=\{y_x:x\in a\}\in A\) by (P2) and every
\(x\in a\) has a witness in \(b\).  These hold for every formula \(\phi\),
in particular for \(\Delta_0\)-formulas, so \(A\) is a transitive model of
\(\mathrm{KPU}\), an admissible set
\cite[ch.~A.7, \S2, 2.4, p.~239]{Makkai1977}.
\end{proof}
\end{annexitem}

\begin{annexitem}{annex:p02-w02c-p02-003-hf-recursion}{thm:p02-w02c-p02-003-hf-recursion}{Recursion on the natural numbers in \(\mathrm{HF}(\mathfrak M)\)}
\begin{proof}
Let \(A=M\cup\mathrm{HF}(M)\).  As in the proof of
\cref{thm:p02-w02c-p02-003-hf}, two sets of \(A\) with the same members are
equal, and every finite subset of \(A\) lies in \(\mathrm{HF}(M)\) (facts
(P2), (P3) there).  Hence \(\mathrm{Emp}(e)\) holds exactly for
\(e=\emptyset\in\mathrm{HF}(M)\), and \(\mathrm{Succ}(i,j)\) exactly when
\(j=i\cup\{i\}\).  For a von Neumann number \(k\), \(k\cup\{k\}=k+1\).

\emph{If \(n\) is a natural number and \(w=F_{v_0}(n)\), the formula
holds.}  Let \(m=n+1\), a natural number in \(\mathrm{HF}(M)\) by part~1 of
\cref{thm:p02-w02c-p02-003-hf}, and
\(s=\{\langle i,F_{v_0}(i)\rangle:i\leq n\}\), which lies in
\(\mathrm{HF}(M)\) by part~2 since every value lies in \(A\).  Then
\(\mathrm{Nat}(n)\), \(\mathrm{Succ}(n,m)\), and \(\mathrm{Seq}(s,m)\)
hold (part~2: \(s\) is a function with domain \(m\)).  For
\(e=\emptyset=0\), \(\mathrm{Val}^{\ast}(s,e,v_0)\) holds as
\(s(0)=v_0\).  Let \(i\in n\) and \(j\in m\) with \(\mathrm{Succ}(i,j)\),
so \(j=i+1\).  If \(\mathrm{Val}^{\ast}(s,i,v)\) and
\(\mathrm{Val}^{\ast}(s,j,w')\), then by part~2 \(v=F_{v_0}(i)\) and
\(w'=F_{v_0}(i+1)=G(i,v)\), so \(\gamma(i,v,w',\bar p)\) holds by the
definition of \(G\).  Finally \(\mathrm{Val}^{\ast}(s,n,w)\) holds as
\(s(n)=F_{v_0}(n)=w\).

\emph{If the formula holds, \(n\) is a natural number and
\(w=F_{v_0}(n)\).}  By \(\mathrm{Nat}(n)\) and part~1, \(n\) is a natural
number.  Let \(m\), \(s\) be witnesses.  Then \(m=n+1\), \(s\) is a
function with domain \(n+1\) (part~2), and the witness \(e\) is
\(\emptyset=0\), so \(s(0)=v_0\).  By induction on \(i\leq n\),
\(s(i)=F_{v_0}(i)\): this holds for \(i=0\); if \(i<n\) and
\(s(i)=F_{v_0}(i)\), then \(i\in n\), \(j=i+1\in m\) and
\(\mathrm{Succ}(i,j)\), and \(\mathrm{Val}^{\ast}(s,i,s(i))\),
\(\mathrm{Val}^{\ast}(s,j,s(j))\) hold, so
\(\gamma(i,s(i),s(j),\bar p)\) holds.  As \(i\) is a natural number and
\(s(i)\) is of the value sort, the uniqueness hypothesis gives
\(s(j)=G(i,s(i))=G(i,F_{v_0}(i))=F_{v_0}(i+1)\).  Finally
\(\mathrm{Val}^{\ast}(s,n,w)\) gives \(w=s(n)=F_{v_0}(n)\).
\end{proof}
\end{annexitem}

\begin{annexitem}{annex:p02-w02c-p02-003-hf-sigma-recursion}{thm:p02-w02c-p02-003-hf-sigma-recursion}{\(\Sigma\)-recursion on membership in \(\mathrm{HF}(\mathfrak M)\)}
\begin{proof}
\emph{Stages.}  For \(a\in A\) let \(\rho(a)=0\) if \(a\in M\), and let
\(\rho(a)\) be the least \(k\) with \(a\in\mathrm{HF}_k(M)\) if
\(a\in\mathrm{HF}(M)\); then \(\rho(a)\geq1\) for every set, since
\(\mathrm{HF}_0(M)=\emptyset\).  If \(a\) is a set and \(u\in a\), then
\(a\) is a finite subset of \(M\cup\mathrm{HF}_{\rho(a)-1}(M)\), so
\(u\in A\) and \(\rho(u)<\rho(a)\).  The stages increase, since
\(\mathrm{HF}_k(M)\subseteq\mathrm{HF}_{k+1}(M)\) by induction on \(k\)
(\(\emptyset\subseteq\mathrm{HF}_1(M)\); and if
\(\mathrm{HF}_{k-1}(M)\subseteq\mathrm{HF}_k(M)\), a finite subset of
\(M\cup\mathrm{HF}_{k-1}(M)\) is one of \(M\cup\mathrm{HF}_k(M)\)).  So
any finite set \(\{a_1,\ldots,a_m\}\subseteq A\) lies in
\(\mathrm{HF}_{k+1}(M)\) with \(k=\max_j\rho(a_j)\), hence in
\(\mathrm{HF}(M)\); in particular Kuratowski pairs of elements of \(A\),
and finite sets of such pairs, lie in \(\mathrm{HF}(M)\).  Every claim
below about all \(a\in A\) is proved by induction on \(\rho(a)\).

\emph{Part~1: the transitive closure.}  For a set \(a\) put
\(C(a)=a\cup\bigcup\{C(u):u\in a\}\), and \(C(p)=\emptyset\) for
\(p\in M\); this is a definition by induction on \(\rho\), since members
of \(a\) have smaller stage.  By induction, \(C(a)\) is a finite subset
of \(A\): \(a\) is finite, each \(C(u)\) is finite, and a finite union of
finite subsets of \(A\) is one.  So \(C(a)\in\mathrm{HF}(M)\).  \(C(a)\)
is transitive: if \(b\in C(a)\) is a set and \(v\in b\), then either
\(b\in a\), and \(v\in b\subseteq C(b)\subseteq C(a)\), or \(b\in C(u)\)
for some \(u\in a\), and \(v\in C(u)\) by the induction hypothesis that
\(C(u)\) is transitive, so \(v\in C(a)\).  \(a\subseteq C(a)\) by
definition.  If \(b\) is transitive and \(a\subseteq b\), then
\(C(a)\subseteq b\): every \(u\in a\) lies in \(b\); if \(u\) is a set,
\(u\subseteq b\) by transitivity, so \(C(u)\subseteq b\) by the induction
hypothesis; if \(u\in M\), \(C(u)=\emptyset\).  Hence
\(C(a)=\mathrm{TC}(a)\), which therefore exists, and
\[
 \mathrm{TC}(a)=a\cup\textstyle\bigcup_{u\in a}\mathrm{TC}(u)
 \quad\text{for every set }a,\qquad
 \mathrm{TC}(p)=\emptyset\text{ for }p\in M.\tag{\(\ast\)}
\]
If \(u\in\mathrm{TC}(a)\) then \(\rho(u)<\rho(a)\): either \(u\in a\), or
\(u\in\mathrm{TC}(u')\) for some \(u'\in a\), and \(\rho(u)<\rho(u')<\rho(a)\)
by induction.  If \(u\in\mathrm{TC}(a)\) is a set, then
\(u\subseteq\mathrm{TC}(a)\), so \(\mathrm{TC}(u)\subseteq\mathrm{TC}(a)\).

\emph{Part~1: existence and uniqueness of \(F\).}  Fix \(\bar x\).  Define
\(F(\bar x,y)\) by induction on \(\rho(y)\): when \(F(\bar x,z)\) has been
defined for all \(z\) with \(\rho(z)<\rho(y)\), in particular for all
\(z\in\mathrm{TC}(y)\), the set \(F_{\bar x}\restriction\mathrm{TC}(y)\)
is a finite set of pairs of elements of \(A\), hence a member of
\(\mathrm{HF}(M)\subseteq A\), and we put \(F(\bar x,y)=G(\bar x,y,F_{\bar
x}\restriction\mathrm{TC}(y))\).  This \(F\) satisfies the identity.  If
\(F'\) also satisfies it, then \(F(\bar x,y)=F'(\bar x,y)\) by induction
on \(\rho(y)\): by the induction hypothesis \(F_{\bar
x}\restriction\mathrm{TC}(y)=F'_{\bar x}\restriction\mathrm{TC}(y)\), and
both values are \(G\) of the same arguments.

\emph{The \(\Delta_0\)-formulas.}  For \(p,z,v\in A\),
\(\mathrm{Pr}(p,z,v)\) holds iff \(p=\{\{z\},\{z,v\}\}\): if \(p\) is that
set, take \(q=\{z\}\), \(q'=\{z,v\}\); conversely the clauses say
\(q=\{z\}\), \(q'=\{z,v\}\) (sets with the same members are equal, and
\(q,q'\) are sets since they have members) and \(p=\{q,q'\}\) (\(p\) has a
member, so it is a set).  In
\(\exists^{f}v\,(z,v)\,\chi\) the bounded quantifiers range over all
\(v\) with \(\langle z,v\rangle\in f\), because \(v\in\{z,v\}\in\langle
z,v\rangle\); so \(\exists^{f}v\,(z,v)\,\chi\) holds iff some \(v\) with
\(\langle z,v\rangle\in f\) satisfies \(\chi\).  By the same remark,
\(\mathrm{Fun}(f,t)\) holds iff \(f\) is a set of pairs with first
coordinates in \(t\), every \(z\in t\) is a first coordinate, and first
coordinates determine second ones: that is, \(f\) is a function with
domain \(t\).  \(\mathrm{Tr}(t)\) says that \(t\) is a transitive set.
When \(\mathrm{Fun}(h,t)\) holds, \(\mathrm{Cl}(h,t)\) says that for
every \(z\in t\), \(h(z)\) is a set whose members are exactly the members
of \(z\) together with the members of the sets \(h(u)\), \(u\in z\); when
\(t\) is transitive, these \(u\) lie in \(t\), so \(h(u)\) is defined.
When \(\mathrm{Fun}(H,t)\), \(\mathrm{Fun}(h,t)\) and \(\mathrm{Fun}(f,t)\)
hold and \(h(z)\subseteq t\) for \(z\in t\), \(\mathrm{Rs}(H,f,h,t)\)
says that \(H(z)\) is the set of pairs in \(f\) whose first coordinate is
in \(h(z)\), that is \(H(z)=f\restriction h(z)\).

\emph{Part~2, left to right.}  Suppose \(F(\bar x,y)=w\).  Let
\(t=\{y\}\cup\mathrm{TC}(y)\), a finite subset of \(A\) and so a member of
\(\mathrm{HF}(M)\); it is transitive, since a set \(b\in t\) is \(y\),
with \(b\subseteq\mathrm{TC}(y)\), or lies in \(\mathrm{TC}(y)\), with
\(b\subseteq\mathrm{TC}(y)\) by transitivity.  Let
\(f=\{\langle z,F(\bar x,z)\rangle:z\in t\}\),
\(h=\{\langle z,\mathrm{TC}(z)\rangle:z\in t\}\) and
\(H=\{\langle z,F_{\bar x}\restriction\mathrm{TC}(z)\rangle:z\in t\}\),
all finite sets of pairs of elements of \(A\), so members of
\(\mathrm{HF}(M)\).  Then \(\mathrm{Fun}\) holds for \(f,h,H\) with domain
\(t\), \(y\in t\) and \(\mathrm{Ap}(f,y,w)\).  \(\mathrm{Cl}(h,t)\) holds
by \((\ast)\) (for \(z\in M\), \(z\) has no members and
\(h(z)=\emptyset\) is a set).  For \(z\in t\),
\(\mathrm{TC}(z)\subseteq t\) (if \(z\) is a set, \(z\subseteq t\) and
\(t\) is transitive), so \(f\restriction h(z)=F_{\bar
x}\restriction\mathrm{TC}(z)=H(z)\) and \(\mathrm{Rs}(H,f,h,t)\) holds.
Finally \(f(z)=F(\bar x,z)=G(\bar x,z,H(z))\), so \(\gamma(\bar
x,z,H(z),f(z),\bar r)\) holds for every \(z\in t\).

\emph{Part~2, right to left.}  Let \(t,f,h,H\) witness
\(\Phi_\gamma[\bar x,y,w,\bar r]\).  First, \(h(z)=\mathrm{TC}(z)\) for
every \(z\in t\), by induction on \(\rho(z)\): the members \(u\) of \(z\)
lie in \(t\) and have smaller stage, so by \(\mathrm{Cl}\) and the
induction hypothesis \(h(z)\) is the set with members those of \(z\) and
of the \(\mathrm{TC}(u)\), \(u\in z\), which is \(\mathrm{TC}(z)\) by
\((\ast)\) (for \(z\in M\) both are the empty set).  So
\(h(z)\subseteq t\) and, by \(\mathrm{Rs}\),
\(H(z)=f\restriction\mathrm{TC}(z)\).  Next, \(f(z)=F(\bar x,z)\) for
every \(z\in t\), by induction on \(\rho(z)\): by the last conjunct,
\(\gamma(\bar x,z,H(z),f(z),\bar r)\) holds, so \(f(z)=G(\bar x,z,H(z))\);
every \(u\in\mathrm{TC}(z)\) lies in \(t\) and has smaller stage, so
\(f(u)=F(\bar x,u)\) and \(H(z)=F_{\bar x}\restriction\mathrm{TC}(z)\);
hence \(f(z)=G(\bar x,z,F_{\bar x}\restriction\mathrm{TC}(z))=F(\bar
x,z)\).  With \(z=y\) and \(\mathrm{Ap}(f,y,w)\), \(w=f(y)=F(\bar x,y)\).

\emph{Part~3.}  \(\mathrm{Pr}\), \(\mathrm{Ap}\), \(\mathrm{Fun}\),
\(\mathrm{Tr}\), \(\mathrm{Cl}\) and \(\mathrm{Rs}\) are
\(\Delta_0\): all their quantifiers are bounded, and
\(\exists^{f}v\,(z,v)\) is a string of bounded quantifiers.  If
\(\gamma\) is \(\Sigma\), then
\(\forall z\in t\,\exists^{H}g\,(z,g)\,\exists^{f}v\,(z,v)\,\gamma\) is
obtained from \(\Sigma\)-formulas by conjunction and bounded
quantification, and \(\Phi_\gamma\) from it and \(\Delta_0\)-formulas
by conjunction and existential quantification; so \(\Phi_\gamma\) is
\(\Sigma\) by the closure clauses of Makkai's 2.5.  If \(G\) is
\(\Sigma\), its graph is defined by such a \(\gamma\) with parameters, and
part~2 shows that \(\Phi_\gamma\) defines the graph of \(F\).

\emph{Part~4.}  By induction on \(\varphi\), for every assignment of
sorts to its free variables.  For an atomic formula: \(U(v)\) becomes
\(v=v\) if \(v\) is an individual variable and \(\neg v=v\) if it is a
set variable; \(v\in v'\) becomes \(\neg v=v\) if \(v'\) is an individual
variable (urelements have no members) and the \(L^{*}\)-atomic formula
\(v\in v'\) otherwise; \(v=v'\) becomes \(\neg v=v\) if the sorts differ
(\(M\cap\mathrm{HF}(M)=\emptyset\)) and \(v=v'\) otherwise; a relation
symbol \(R(v_1,\ldots,v_k)\) for a relation of \(\mathfrak M\) becomes
\(R(v_1,\ldots,v_k)\) if all \(v_i\) are individual variables and \(\neg
v_1=v_1\) otherwise.  Connectives are kept, and \(\exists v\,\varphi\)
becomes \(\exists x\,\varphi^{\mathrm s,x}\lor\exists
a\,\varphi^{\mathrm s,a}\), with \(x\) an individual and \(a\) a set
variable and \(\varphi^{\mathrm s,x}\), \(\varphi^{\mathrm s,a}\) the
translations with \(v\) given the corresponding sort (and renamed);
\(\forall\) is treated dually.  Each clause matches the meaning in \(A\)
of the atomic formula for arguments of the given sorts, and the
quantifier clause matches because \(A=M\cup\mathrm{HF}(M)\), with
individual variables ranging over \(M\) and set variables over
\(\mathrm{HF}(M)\).  Applying this to \(\Phi_\gamma\) gives the last claim.
\end{proof}
\end{annexitem}

\begin{annexitem}{annex:p02-w02c-p02-003-wso-sentences}{thm:p02-w02c-p02-003-wso-sentences}{The \S2 sentences in \(L^{*}\)}
\begin{proof}
All facts about \(\mathrm{Nat}\), \(\mathrm{Seq}\), \(\mathrm{Val}\),
\(\mathrm{Val}'\) are from parts~1 and~2 of \cref{thm:p02-w02c-p02-003-hf}.

\emph{(M1) \(\mathrm{Mul}\).}  For every natural number \(i\) and
\(v\in M\) there is exactly one \(w\in M\) with \(w=v+^{\mathfrak M}x\), so
\cref{thm:p02-w02c-p02-003-hf-recursion} applies to \(\gamma_+\) with
parameter \(x\in M\) and start value \(x\): \(F_x(0)=x\) and
\(F_x(i+1)=F_x(i)+^{\mathfrak M}x\).  By induction on \(i\),
\(F_x(i)=((i+1)x)^{\mathfrak M}\): for \(i=0\) the term \(1x\) is \(x\);
and \(((i+1)x)^{\mathfrak M}+^{\mathfrak M}x\) is the value of the term
\(((i+1)x+x)\), which is \((i+2)x\).  If
\(\mathrm{HF}(\mathfrak M)\models\mathrm{Mul}[n,x,y]\), some \(k\) has
\(n=k\cup\{k\}\) and \(\mathrm{Rec}_{\gamma_+}[k,x,y,x]\), so \(k\) is a
natural number, \(n=k+1\geq1\) and \(y=F_x(k)=(nx)^{\mathfrak M}\).
Conversely, for a natural number \(n\geq1\) and \(y=(nx)^{\mathfrak M}\),
\(k=n-1\) witnesses \(\mathrm{Mul}[n,x,y]\).  So
\(\mathrm{Mul}[n,x,y]\) holds iff \(n\) is a natural number \(n\geq1\) and
\(y=(nx)^{\mathfrak M}\).

\emph{(M2) \(\mathrm{Add}\).}  For every set \(v\in\mathrm{HF}(M)\),
\(v\cup\{v\}\) is a finite subset of \(M\cup\mathrm{HF}(M)\), hence in
\(\mathrm{HF}(M)\), and it is the only \(w\) with \(\mathrm{Succ}(v,w)\).
So the theorem applies to \(\sigma\) with the sets as value sort:
\(\mathrm{Add}[a,i,w]\) holds iff \(i\) is a natural number and
\(w=F_a(i)\), where \(F_a(0)=a\) and \(F_a(i+1)=F_a(i)\cup\{F_a(i)\}\).
For a natural number \(a\), \(F_a(i)=a+i\) by induction on \(i\).

\emph{(M3) \(\mathrm{Times}\).}  Fix a natural number \(a\).  By (M2), for
every natural number \(i\) and every set \(v\), exactly one \(w\)
satisfies \(\tau(i,v,w,a)=\mathrm{Add}(v,a,w)\), namely \(F_v(a)\), which
is \(v+a\) when \(v\) is a natural number.  The theorem applies to
\(\tau\) with parameter \(a\) and start value \(\emptyset=0\): the function
\(H\) with \(H(0)=0\) and \(H(i+1)=F_{H(i)}(a)\) has \(H(i)=a\cdot i\), by
induction on \(i\).  So for natural numbers \(a\), \(b\),
\(\mathrm{Times}[a,b,w]\) holds iff \(w=a\cdot b\).

\emph{(M4) \(\mathrm{Prime}\).}  Let \(p\in\mathrm{HF}(M)\).  If
\(\mathrm{Nat}(p)\) fails, so do \(\mathrm{Prime}(p)\) and the claim that
\(p\) is a prime number.  Let \(p\) be a natural number.  Its members are
the numbers \(k<p\), and \(k\) has a member iff \(k\geq1\); so the second
conjunct says \(p\geq2\).  By (M3) the third says that no \(a,b<p\) have
\(a\cdot b=p\).  If \(p\geq2\) is prime, a product \(a\cdot b=p\) has
\(\{a,b\}=\{1,p\}\), so one factor is not \(<p\).  If \(p\geq2\) is not
prime, \(p=a\cdot b\) with \(2\leq a,b\), and then \(a,b\leq p/2<p\).  So
\(\mathrm{Prime}(p)\) holds iff \(p\) is a prime number.  Likewise
\(\mathrm{Pos}(n)\) holds iff \(n\) is a natural number \(n\geq1\), since
the only number without members is \(0\).

\emph{The sentences.}  The quantifier \(\forall n\) of \((6)^{*}\) ranges
over \(\mathrm{HF}(M)\), and \(\mathrm{Pos}(n)\) restricts it to the
natural numbers \(n\geq1\); for such \(n\), by (M1),
\(\forall x\,\exists y\,\mathrm{Mul}(n,y,x)\) holds iff every \(x\in M\)
equals \((ny)^{\mathfrak M}\) for some \(y\in M\).  This is the claim for
\((6)^{*}\).  For \((7)^{*}\), (M1) shows that for \(n\geq1\) the matrix
says that \(x\neq0^{\mathfrak M}\) implies
\((nx)^{\mathfrak M}\neq0^{\mathfrak M}\).  For \((8)^{*}\),
\(\mathrm{Mul}(n,x,0)\) holds iff \(n\geq1\) is a natural number and
\((nx)^{\mathfrak M}=0^{\mathfrak M}\).  For \((10)^{*}\), (M4) restricts
\(p\) to the prime numbers, and by (M1) with \(x=1^{\mathfrak M}\) the
matrix says \((p1)^{\mathfrak M}\neq0^{\mathfrak M}\); Barwise's term
\(p1\) is \(((1+1)+\cdots+1)\) with \(p\) occurrences of \(1\), which is
the term \(px\) with \(x=1\).  Each \((10)_p\) is a first-order sentence,
true in \(\mathfrak M\) iff \((p1)^{\mathfrak M}\neq0^{\mathfrak M}\).
\end{proof}
\end{annexitem}

\begin{annexitem}{annex:p02-w02c-p02-003-infinitary-rule}{thm:p02-w02c-p02-003-infinitary-rule}{No finitary complete relation for \(L_{\omega_1\omega}\)}
\begin{proof}
Let \(T\) be the countable set of sentences of \(L_{\omega_1\omega}\) of
the empty language given in part~3 of
\cref{thm:p02-w02-p02-005-bw-extended-ls}: every finite subset of \(T\)
has a model and \(T\) has none.  Let \(\bot\) be \(\neg\exists x\,(x=x)\),
which no structure satisfies since domains are nonempty.  Every model of
\(T\) satisfies \(\bot\), because \(T\) has no model.  Suppose
\(T\vdash\bot\) for a sound finitary relation \(\vdash\).  Then
\(T_0\vdash\bot\) for a finite \(T_0\subseteq T\), and by soundness every
model of \(T_0\) satisfies \(\bot\).  But \(T_0\) has a model, and no
structure satisfies \(\bot\).  So \(T\nvdash\bot\).

If derivability is defined by finite derivations in which only finitely
many premises occur, and a derivation from \(\Gamma\) whose premises all
lie in \(\Gamma_0\subseteq\Gamma\) is a derivation from \(\Gamma_0\), then
derivability is finitary, and the argument applies to it.
\end{proof}
\end{annexitem}

\begin{annexitem}{annex:p02-w02c-p02-003-wso-translation}{thm:p02-w02c-p02-003-wso-translation}{Weak second order into \(L_{\omega_1\omega}\)}
\begin{proof}
\emph{Part~1.}  Fix distinct objects \(u_1,u_2,\ldots\) that are not sets,
used as atoms.  For \(k\geq0\) let \(X_k=\{u_1,\ldots,u_k\}\) and
\(H_k=\mathrm{HF}(X_k)\), formed as in
\cref{def:p02-w02c-p02-003-wso-reading}.  Each level
\(\mathrm{HF}_n(X_k)\) is finite (induction on \(n\): the finite subsets of a
finite set), so \(H_k\) is countable, \(H_k\subseteq H_{k+1}\), and
\(H=\bigcup_k H_k\) is countable (a countable union of countable sets; the
axiom of choice is assumed).

\emph{(A) Images.}  Let \(\bar m=(m_1,\ldots,m_k)\) be pairwise distinct
elements of \(M\) and \(F=\{m_1,\ldots,m_k\}\).  Define
\(\pi_{\bar m}\) on \(X_k\cup H_k\) by \(\pi_{\bar m}(u_i)=m_i\) and
\(\pi_{\bar m}(h)=\{\pi_{\bar m}(e):e\in h\}\), by recursion on the least
\(n\) with \(h\in\mathrm{HF}_n(X_k)\).  By induction on \(n\),
\(\pi_{\bar m}\) maps \(\mathrm{HF}_n(X_k)\) bijectively onto
\(\mathrm{HF}_n(F)\): if this holds for \(n\), then \(\pi_{\bar m}\) is a
bijection from \(X_k\cup\mathrm{HF}_n(X_k)\) onto \(F\cup\mathrm{HF}_n(F)\)
(the pieces are disjoint because atoms are not sets), and a bijection
induces a bijection between the sets of finite subsets, \(h\mapsto\) its
image, which is \(\pi_{\bar m}\) on \(\mathrm{HF}_{n+1}(X_k)\).  Hence
\(\pi_{\bar m}\) is injective on \(X_k\cup H_k\), maps \(H_k\) onto
\(\mathrm{HF}(F)\subseteq\mathrm{HF}(M)\) and \(X_k\) onto \(F\subseteq M\),
and for \(e\in X_k\cup H_k\), \(h\in H_k\): \(e\in h\) iff
\(\pi_{\bar m}(e)\in\pi_{\bar m}(h)\).  If \(\bar m'\) extends \(\bar m\)
by further distinct elements, \(\pi_{\bar m'}\) extends \(\pi_{\bar m}\),
by induction on the same rank.

\emph{(B) Supports.}  Every \(b\in\mathrm{HF}(M)\) lies in \(\mathrm{HF}(F)\)
for some finite \(F\subseteq M\): if \(b\in\mathrm{HF}_{n+1}(M)\), each of
its finitely many members is in \(M\) or, by induction, in some
\(\mathrm{HF}(F_j)\) with \(F_j\) finite; with \(F\) the union of these
\(F_j\) and the members of \(b\) in \(M\), \(b\) is a finite subset of
\(F\cup\mathrm{HF}_n(F)\), since \(\mathrm{HF}_n(X)\) grows with \(X\) and
with \(n\).

\emph{(C) The translation.}  Rename bound set variables so that no set
variable is bound twice or both bound and free, and fix distinct
individual variables \(z_1,z_2,\ldots\) not occurring in the given
sentence.  For a subformula
\(\varphi\) whose free set variables are among \(a_1,\ldots,a_n\), each
\(k\), and each \(\bar h=(h_1,\ldots,h_n)\in H_k^{\,n}\), define a formula
\(\varphi^{k,\bar h}\) of \(L_{\omega_1\omega}\) for \(L\), with free
variables among the free individual variables of \(\varphi\) and
\(z_1,\ldots,z_k\).  Write \(\top\) for \(\exists x\,(x=x)\) and \(\bot\)
for \(\neg\top\).
\begin{itemize}
\item An atomic formula of \(L\) is unchanged.
\item \((t\in a_j)^{k,\bar h}\) is \(\bigvee\{t=z_i:u_i\in h_j\}\), and
      \(\bot\) if there is no such \(i\).
\item \((a_j\in a_l)^{k,\bar h}\) is \(\top\) if \(h_j\in h_l\) and \(\bot\)
      otherwise; \((a_j=a_l)^{k,\bar h}\) is \(\top\) if \(h_j=h_l\) and
      \(\bot\) otherwise.
\item \({}^{k,\bar h}\) commutes with \(\neg\), \(\land\), \(\exists x\),
      \(\forall x\); \(\forall a\) is read as \(\neg\exists a\neg\).
\item \((\exists a_{n+1}\,\psi)^{k,\bar h}\) is the disjunction, over all
      \(k'\geq k\) and \(g\in H_{k'}\), of
      \(\exists z_{k+1}\cdots\exists z_{k'}\,
      \bigl(\bigwedge_{i<j\leq k',\,j>k}z_i\neq z_j\land
      \psi^{k',(\bar h,g)}\bigr)\).
\end{itemize}
The last disjunction is indexed by a subset of \(\omega\times H\), so it is
countable, and each formula has finitely many free variables.

\emph{(D) Correctness.}  By induction on \(\varphi\), for all \(k\),
\(\bar h\), pairwise distinct \(\bar m\in M^k\) and assignments \(s\) to the
individual variables:
\(\mathrm{HF}(\mathfrak M)\models\varphi[s,a_j\mapsto\pi_{\bar m}(h_j)]\)
iff \(\mathfrak M\models\varphi^{k,\bar h}[s(m_1/z_1)\cdots(m_k/z_k)]\).
Atomic formulas of \(L\) do not contain \(z_i\) or set variables.  For
\(t\in a_j\): the value \(t^{\mathfrak M}(s)\) lies in \(M\), and the
members of \(\pi_{\bar m}(h_j)\) that lie in \(M\) are exactly the
\(m_i\) with \(u_i\in h_j\), by (A); so the left side holds iff
\(t^{\mathfrak M}(s)=m_i\) for some such \(i\).  For \(a_j\in a_l\) and
\(a_j=a_l\), (A) gives \(\pi_{\bar m}(h_j)\in\pi_{\bar m}(h_l)\) iff
\(h_j\in h_l\), and \(\pi_{\bar m}(h_j)=\pi_{\bar m}(h_l)\) iff
\(h_j=h_l\).  The Boolean cases and \(\exists x\) follow from the induction
hypothesis, since the variables \(z_i\) differ from \(x\).  For
\(\exists a_{n+1}\,\psi\): if some \(b\in\mathrm{HF}(M)\) satisfies
\(\psi\) with \(a_{n+1}\mapsto b\), take \(F\) finite with
\(b\in\mathrm{HF}(F)\) by (B), list \(F\setminus\{m_1,\ldots,m_k\}\) as
\(m_{k+1},\ldots,m_{k'}\), and let \(\bar m'=(m_1,\ldots,m_{k'})\).  By
(A), \(b=\pi_{\bar m'}(g)\) for some \(g\in H_{k'}\), and
\(\pi_{\bar m'}(h_j)=\pi_{\bar m}(h_j)\).  The induction hypothesis for
\(\psi\), \(k'\), \((\bar h,g)\) and \(\bar m'\) shows that the disjunct for
\(k',g\) holds, with \(z_{k+1},\ldots,z_{k'}\) sent to
\(m_{k+1},\ldots,m_{k'}\), which are distinct from each other and from
\(m_1,\ldots,m_k\).  Conversely, if the disjunct for \(k',g\) holds with
witnesses \(m_{k+1},\ldots,m_{k'}\), these make \(\bar m'\) pairwise
distinct, and the induction hypothesis gives that \(b=\pi_{\bar m'}(g)\)
satisfies \(\psi\), so \(\exists a_{n+1}\,\psi\) holds.

For a sentence \(\varphi\) put \(\varphi^{\natural}=\varphi^{0,()}\); (D)
with \(k=n=0\) is the claim.

\emph{Part~2.}  For \(n\geq1\) let \(\varepsilon_n\) be
\(\exists x_1\cdots\exists x_n\,(\bigwedge_{i<j\leq n}x_i\neq x_j\land
\forall y\,(y=x_1\lor\cdots\lor y=x_n))\), and for
\(S\subseteq\{1,2,\ldots\}\) let \(\psi_S=\bigvee\{\varepsilon_n:n\in S\}\),
a sentence of \(L_{\omega_1\omega}\).  Let \(\mathfrak M_n\) be the
structure for the empty language with domain \(\{1,\ldots,n\}\).  Then
\(\mathfrak M_n\models\varepsilon_m\) iff \(m=n\), so
\(\{n:\mathfrak M_n\models\psi_S\}=S\).  Suppose that for every \(S\) some
sentence \(\varphi\) of \(L^{*}\) satisfied
\(\mathrm{HF}(\mathfrak M)\models\varphi\) iff \(\mathfrak M\models\psi_S\)
for all \(\mathfrak M\).  Fix an enumeration
\(\varphi_0,\varphi_1,\ldots\) of the sentences of \(L^{*}\); there are
countably many, since they are finite strings over a countable alphabet.
Send \(S\) to the least \(i\) such that \(\varphi_i\) works for \(S\).
Then \(S=\{n:\mathrm{HF}(\mathfrak M_n)\models\varphi_i\}\) is determined by
\(i\), so the map is an injection of the power set of \(\{1,2,\ldots\}\)
into \(\omega\), contradicting Cantor's theorem.  Hence some \(\psi_S\) has
no such \(\varphi\).
\end{proof}
\end{annexitem}

\begin{annexitem}{annex:p02-w02c-p02-003-nnf}{thm:p02-w02c-p02-003-nnf}{Negation normal form}
\begin{proof}
By induction on the generation of \(\chi\), both claims together.  If
\(\chi\) is atomic, \(\chi^{+}=\chi\) and \(\chi^{\ominus}=\neg\chi\) are
atomic and negated atomic, hence in n.n.f., with the variables and symbols
of \(\chi\); \(\chi^{+}\) holds at \(s\) iff \(\chi\) does, and
\(\neg\chi\) holds iff \(\chi\) fails.  If \(\chi=\neg\chi_1\), then
\(\chi^{+}=\chi_1^{\ominus}\) holds iff \(\chi_1\) fails iff \(\chi\)
holds, and \(\chi^{\ominus}=\chi_1^{+}\) holds iff \(\chi_1\) holds iff
\(\chi\) fails.  If \(\chi=\bigwedge\Phi\), then
\(\{\varphi^{+}:\varphi\in\Phi\}\) and
\(\{\varphi^{\ominus}:\varphi\in\Phi\}\) are images of the countable set
\(\Phi\), hence countable, and their members have, by the induction
hypothesis, the free variables and symbols of the members of \(\Phi\); so
\(\chi^{+}\) and \(\chi^{\ominus}\) are formulas of \(L_{\omega_1\omega}\)
in n.n.f. with the free variables and symbols of \(\chi\).  At \(s\),
\(\chi^{+}\) holds iff every \(\varphi^{+}\) holds iff every \(\varphi\)
holds iff \(\chi\) holds; \(\chi^{\ominus}\) holds iff some
\(\varphi^{\ominus}\) holds iff some \(\varphi\) fails iff \(\chi\) fails.
If \(\chi=\bigvee\Phi\), \(\chi^{+}\) holds iff some \(\varphi\) holds iff
\(\chi\) holds, and \(\chi^{\ominus}=\bigwedge\{\varphi^{\ominus}\}\) holds
iff every \(\varphi\) fails iff \(\chi\) fails.  If \(\chi=\exists
x\,\chi_1\), then \(\chi^{+}=\exists x\,\chi_1^{+}\) holds at \(s\) iff
\(\chi_1^{+}\) holds at some \(s(b/x)\) iff \(\chi_1\) holds at some
\(s(b/x)\) iff \(\chi\) holds; \(\chi^{\ominus}=\forall x\,\chi_1^{\ominus}\)
holds iff \(\chi_1\) fails at every \(s(b/x)\) iff \(\chi\) fails.  If
\(\chi=\forall x\,\chi_1\), \(\chi^{+}=\forall x\,\chi_1^{+}\) holds iff
\(\chi_1\) holds at every \(s(b/x)\) iff \(\chi\) holds, and
\(\chi^{\ominus}=\exists x\,\chi_1^{\ominus}\) holds iff \(\chi_1\) fails
at some \(s(b/x)\) iff \(\chi\) fails.  In the last two cases the free
variables of \(\chi^{+}\), \(\chi^{\ominus}\) are those of \(\chi_1\) other
than \(x\).
\end{proof}
\end{annexitem}

\begin{annexitem}{annex:p02-w02c-p02-003-hintikka-model}{thm:p02-w02c-p02-003-hintikka-model}{A Hintikka set has a model}
\begin{proof}
\emph{(S) Substitution of a constant.}  For every formula \(\chi\) of
\(L_{\omega_1\omega}\) for \(K\cup C\), variable \(x\), constant
\(c\in C\), structure \(\mathfrak M\) and assignment \(s\):
\(\mathfrak M\models\chi(c/x)[s]\) iff
\(\mathfrak M\models\chi[s(c^{\mathfrak M}/x)]\).  The terms of \(K\cup C\)
are variables and constants.  The value of \(t(c/x)\) at \(s\) equals the
value of \(t\) at \(s(c^{\mathfrak M}/x)\): for \(t=x\) both are
\(c^{\mathfrak M}\); for another variable \(y\) both are \(s(y)\); for a
constant \(d\) both are \(d^{\mathfrak M}\).  Atomic formulas follow, and
\(\neg\), \(\bigwedge\), \(\bigvee\) commute with the substitution and
refer to the same assignment.  For \(\exists y\,\chi_1\) with \(y=x\),
\(x\) has no free occurrence, the formula is unchanged, and it holds at
\(s\) iff \(\chi_1\) holds at some \(s(b/x)\), which is
\(s(c^{\mathfrak M}/x)(b/x)\); so it holds at \(s\) iff it holds at
\(s(c^{\mathfrak M}/x)\).  For \(y\neq x\), \((\exists
y\,\chi_1)(c/x)=\exists y\,\chi_1(c/x)\) holds at \(s\) iff for some
\(b\), \(\chi_1(c/x)\) holds at \(s(b/y)\), iff (induction hypothesis)
\(\chi_1\) holds at \(s(b/y)(c^{\mathfrak M}/x)=s(c^{\mathfrak M}/x)(b/y)\)
for some \(b\), iff \(\exists y\,\chi_1\) holds at
\(s(c^{\mathfrak M}/x)\).  \(\forall y\) is treated by the same
computation with ``for every \(b\)''.  Since a constant contains no
variable, substituting constants for several distinct variables one after
another gives the same result in any order; for a map \(\sigma\) from a
finite set of variables to \(C\), \(\chi\sigma\) denotes the result.

\emph{The equivalence.}  For \(c,d\in C\) put \(c\sim d\) iff
\((c=d)\in H\).  Reflexivity is (iia).  If \((c=d)\in H\), apply (iib) to
\(\varphi(v)=(v=c)\): \(\varphi(c/v)=(c=c)\in H\) by (iia), so
\(\varphi(d/v)=(d=c)\in H\); this is symmetry.  If \((c=d)\in H\) and
\((d=e)\in H\), then \((d=c)\in H\) by symmetry, and (iib) applied to
\(\varphi(v)=(v=e)\), with \(\varphi(d/v)\in H\) and \((d=c)\in H\), gives
\((c=e)\in H\); this is transitivity.

\emph{The structure.}  Let \(M=\{[c]:c\in C\}\), the set of
\(\sim\)-classes; it is nonempty since \(C\) is.  Put
\(c^{\mathfrak M}=[c]\) and, for an \(n\)-ary \(R\in K\),
\(R^{\mathfrak M}=\{([c_1],\ldots,[c_n]):R(c_1,\ldots,c_n)\in H\}\).  Every
element of \(M\) is some \(c^{\mathfrak M}\).  Claim: if
\(R(d_1,\ldots,d_n)\in H\) and \(d_i\sim c_i\) for each \(i\), then
\(R(c_1,\ldots,c_n)\in H\).  For \(j=1,\ldots,n\) in turn, let
\(\varphi_j(v)=R(c_1,\ldots,c_{j-1},v,d_{j+1},\ldots,d_n)\); from
\(\varphi_j(d_j/v)\in H\) and \((d_j=c_j)\in H\), (iib) gives
\(\varphi_j(c_j/v)\in H\), which is \(\varphi_{j+1}(d_{j+1}/v)\).  Hence
\(\mathfrak M\models R(c_1,\ldots,c_n)\) iff \(R(c_1,\ldots,c_n)\in H\), and
\(\mathfrak M\models c=d\) iff \([c]=[d]\) iff \((c=d)\in H\).

\emph{Truth.}  By induction on the n.n.f. formula \(\varphi\): for every
map \(\sigma\) from the free variables of \(\varphi\) to \(C\), if
\(\varphi\sigma\in H\) then \(\mathfrak M\models\varphi\sigma\).  If
\(\varphi\) is atomic this is the previous paragraph.  If \(\varphi=\neg\theta\)
with \(\theta\) atomic and \(\neg\theta\sigma\in H\), then
\(\theta\sigma\notin H\) by (i), so \(\mathfrak M\not\models\theta\sigma\)
by the previous paragraph.  If \(\varphi=\bigwedge\Phi\), then
\(\varphi\sigma=\bigwedge\{\chi\sigma_\chi:\chi\in\Phi\}\), with
\(\sigma_\chi\) the restriction of \(\sigma\) to the free variables of
\(\chi\); by (iii) each \(\chi\sigma_\chi\in H\), by the induction
hypothesis each holds, so \(\varphi\sigma\) holds.  If
\(\varphi=\bigvee\Phi\), (v) gives one \(\chi\sigma_\chi\in H\), which holds.
If \(\varphi=\forall x\,\chi\), then \(\varphi\sigma=\forall
x\,(\chi\sigma)\), \(x\) not being in the domain of \(\sigma\), and by (iv)
\((\chi\sigma)(c/x)\in H\) for every \(c\in C\).  This is
\(\chi\sigma'\) for \(\sigma'=\sigma\cup\{(x,c)\}\), so by the induction
hypothesis \(\mathfrak M\models(\chi\sigma)(c/x)\), and by (S)
\(\mathfrak M\models\chi\sigma[s(c^{\mathfrak M}/x)]\) for every
assignment \(s\).  As every element of \(M\) is some \(c^{\mathfrak M}\),
\(\mathfrak M\models\forall x\,(\chi\sigma)\).  If
\(\varphi=\exists x\,\chi\), (vi) gives one \(c\) with
\((\chi\sigma)(c/x)\in H\), which holds, and by (S) \(c^{\mathfrak M}\)
witnesses \(\exists x\,(\chi\sigma)\).  A member of \(H\) is an n.n.f.
sentence \(\varphi\), the case of empty \(\sigma\).
\end{proof}
\end{annexitem}

\begin{annexitem}{annex:p02-w02c-p02-003-mex}{thm:p02-w02c-p02-003-model-existence}{Model existence}
\begin{proof}
Let \(s_0\in\mathcal S\).

\emph{A countable stock of sentences.}  Let \(A_0\) be \(s_0\) together
with all atomic sentences of \(K\cup C\); these are finite strings over
the countable set \(K\cup C\cup\{=\}\), so \(A_0\) is countable.  Given
\(A_m\), let \(A_{m+1}\) add, for each \(\chi\in A_m\) of the form
\(\bigwedge\Phi\) or \(\bigvee\Phi\), the members of \(\Phi\), and for each
\(\chi\in A_m\) of the form \(\forall x\,\varphi\) or \(\exists
x\,\varphi\), the sentences \(\varphi(c/x)\), \(c\in C\).  Each \(\chi\)
contributes countably many sentences, so each \(A_m\) is countable, and
\(A=\bigcup_m A_m\) is countable (countable unions of countable sets, by
the axiom of choice).  \(A\) contains \(s_0\) and all atomic sentences, and
if \(\bigwedge\Phi\in A\) or \(\bigvee\Phi\in A\) then \(\Phi\subseteq A\),
and if \(\forall x\,\varphi\in A\) or \(\exists x\,\varphi\in A\) then
\(\varphi(c/x)\in A\) for every \(c\in C\).  Members of \(A\) are n.n.f.
sentences: the members of \(s_0\) are, atomic sentences are, and the
members of \(\Phi\) and the \(\varphi(c/x)\) obtained from an n.n.f.
sentence are n.n.f. sentences.

\emph{Enumerations.}  Fix enumerations \(a_0,a_1,\ldots\) of \(A\) and
\(c_0,c_1,\ldots\) of \(C\), a variable \(x_0\), and an enumeration
\(\alpha_0,\alpha_1,\ldots\) of the atomic formulas of \(K\cup C\) whose
only variable is \(x_0\) (a countable set of finite strings).  Fix a
sequence \((\varphi_n,\psi_n,k_n)_{n<\omega}\) of members of
\(A\times A\times\omega\) in which every triple occurs for infinitely many
\(n\): compose a surjection of \(\omega\) onto \(\omega^{4}\) with
\((i,j,k,l)\mapsto(a_i,a_j,k)\).

\emph{The chain.}  Put \(s_0\) as given; suppose \(s_n\in\mathcal S\) is a
finite subset of \(A\).  Let \(t=s_n\) unless \(\varphi_n\in s_n\), in
which case \(t\) is given by the form of \(\varphi_n\):
\begin{itemize}
\item \(\varphi_n=\bigwedge\Phi\): \(t=s_n\cup\{a_{k_n}\}\) if
      \(a_{k_n}\in\Phi\), else \(t=s_n\); \(t\in\mathcal S\) by (C3).
\item \(\varphi_n=\forall x\,\varphi\): \(t=s_n\cup\{\varphi(c_{k_n}/x)\}\),
      in \(\mathcal S\) by (C4).
\item \(\varphi_n=\bigvee\Phi\): \(t=s_n\cup\{a_j\}\) for the least \(j\)
      with \(a_j\in\Phi\) and \(s_n\cup\{a_j\}\in\mathcal S\); such \(j\)
      exists by (C5), as \(\Phi\subseteq A\).
\item \(\varphi_n=\exists x\,\varphi\): \(t=s_n\cup\{\varphi(c_j/x)\}\) for
      the least \(j\) with this set in \(\mathcal S\), which exists by (C6).
\item \(\varphi_n=(c=d)\) with \(c,d\in C\): if \(\psi_n\in s_n\) and
      \(\psi_n=\alpha_{k_n}(c/x_0)\), then
      \(t=s_n\cup\{\alpha_{k_n}(d/x_0)\}\), in \(\mathcal S\) by (C2);
      otherwise \(t=s_n\).
\item any other form: \(t=s_n\).
\end{itemize}
Then put \(s_{n+1}=t\cup\{c_{k_n}=c_{k_n}\}\), in \(\mathcal S\) by (C1).
Each step adds at most two members of \(A\), so \(s_{n+1}\) is a finite
subset of \(A\) and \(s_n\subseteq s_{n+1}\).

\emph{\(H=\bigcup_n s_n\) is a Hintikka set for \(K\cup C\).}  Its members
are n.n.f. sentences.  (i): if \(\theta,\neg\theta\in H\), both lie in one
\(s_n\), against (C0).  (iia): \(c=c_k\) gives
\((c=c)\in s_{n+1}\) for any \(n\) with \(k_n=k\).  (iib): let
\(\varphi(v)\) be atomic with no variable but \(v\), and
\(\varphi(c/v),(c=d)\in s_m\).  Replacing \(v\) by \(x_0\) gives an atomic
\(\alpha_k\) with \(\alpha_k(e/x_0)=\varphi(e/v)\) for every \(e\in C\).
Take \(n\geq m\) with \((\varphi_n,\psi_n,k_n)=((c=d),\varphi(c/v),k)\);
then \(\varphi(d/v)\in s_{n+1}\).  (iii): if \(\bigwedge\Phi\in s_m\) and
\(a_k\in\Phi\), take \(n\geq m\) with \(\varphi_n=\bigwedge\Phi\),
\(k_n=k\).  (iv): if \(\forall x\,\varphi\in s_m\) and \(c=c_k\), take
\(n\geq m\) with \(\varphi_n=\forall x\,\varphi\), \(k_n=k\).  (v) and (vi):
take \(n\geq m\) with \(\varphi_n\) the given disjunction or existential
sentence.  In each case the required sentence lies in \(s_{n+1}\).

By \cref{thm:p02-w02c-p02-003-hintikka-model}, \(H\) has a model, which is
a model of \(s_0\subseteq H\).
\end{proof}
\end{annexitem}

\begin{annexitem}{annex:p02-w02c-p02-003-interpolation}{thm:p02-w02c-p02-003-interpolation}{Interpolation for \(L_{\omega_1\omega}\)}
\begin{proof}
The proof uses facts (a), (b), (c) proved at the start of
\cref{annex:p02-w02c-p02-003-beth} (renaming and reducts, renaming
variables, constants to variables), which do not depend on the present
theorem, and (S) of \cref{annex:p02-w02c-p02-003-hintikka-model}, whose
proof applies to any language whose terms are variables and constants.
By (a), whether \(\chi\) is valid, and whether \(\Gamma\models\theta\),
does not depend on the language chosen to contain the symbols that occur,
since every structure for a smaller language expands to one for a larger
language (interpret the further symbols arbitrarily on the nonempty
domain).

\emph{Countability.}  By induction on formulas, the set of non-logical
symbols occurring in a formula of \(L_{\omega_1\omega}\) is countable: an
atomic formula has finitely many, and the set for \(\bigwedge\Phi\) or
\(\bigvee\Phi\) is a countable union of countable sets (axiom of choice).

\emph{Step 1: relational languages.}  Let \(\varphi\), \(\psi\) be sentences
for a relational language, \(\varphi\to\psi\) valid, \(K_1\) and \(K_2\)
the sets of relation symbols occurring in \(\varphi\) and in \(\psi\),
\(K=K_1\cup K_2\) (countable), and \(C\) a countably infinite set of new
constant symbols.  For a set \(\Delta\) of sentences, \(C(\Delta)\) is the
set of members of \(C\) occurring in members of \(\Delta\).  All
structures below are for \(K\cup C\).

A \emph{pair} is \((\Gamma_1,\Gamma_2)\) with \(\Gamma_i\) a finite set of
n.n.f. sentences of \(L_{\omega_1\omega}\) for \(K_i\cup C\) and
\(C(\Gamma_1\cup\Gamma_2)\) finite.  A \emph{separator} of the pair is a
sentence \(\theta\) of \(L_{\omega_1\omega}\) whose non-logical symbols lie
in \((K_1\cap K_2)\cup C(\Gamma_1\cup\Gamma_2)\), with
\(\Gamma_1\models\theta\) and \(\Gamma_2\models\neg\theta\).  The pair is
\emph{inseparable} if it has no separator.  Let \(\mathcal S\) be the set
of all \(\Gamma_1\cup\Gamma_2\) for inseparable pairs.  \emph{Symmetry}:
\(\theta\) separates \((\Gamma_1,\Gamma_2)\) iff \(\neg\theta\) separates
\((\Gamma_2,\Gamma_1)\) with the roles of \(K_1\) and \(K_2\) exchanged,
since \(\Gamma_1\models\theta\) iff \(\Gamma_1\models\neg\neg\theta\).  So
each claim below about adding a sentence to the first member of an
inseparable pair, proved for every \(K_1,K_2\), gives the same claim for
the second member by applying it to \((\Gamma_2,\Gamma_1)\).

\emph{(E) Removing a new constant.}  Let \(\Delta\) be a set of sentences,
\(\theta\) a sentence, and \(c\in C\) a constant occurring in no member of
\(\Delta\).  As in the proof of \cref{thm:p02-w02c-p02-003-beth}, take
distinct variables \(u_0,u_1,\ldots\), let \(\pi(u_k)=u_{k+1}\) and
\(\pi(w)=w\) for every other variable \(w\), put \(y=u_0\), and let
\(\hat\theta=\theta^{\pi}[y/c]\), which replaces \(c\) by \(y\) in
\(\theta^{\pi}\); \(y\) does not occur in \(\theta^{\pi}\).  Then
\(\hat\theta\) has at most \(y\) free, and its non-logical symbols are
those of \(\theta\) other than \(c\).  Let \(\mathfrak B\models\Delta\),
\(b\) an element, and \(\mathfrak B'\) the structure that differs from
\(\mathfrak B\) only in \(c^{\mathfrak B'}=b\).  By (a) for the reduct to
the language without \(c\), \(\mathfrak B'\models\Delta\); by (b),
\(\mathfrak B'\models\theta\) iff \(\mathfrak B'\models\theta^{\pi}\);
by (c) and (a), \(\mathfrak B'\models\theta^{\pi}\) iff
\(\mathfrak B\models\hat\theta[y\mapsto b]\).  Hence:
(E1) if \(\Delta\models\theta\), then \(\Delta\models\forall y\,\hat\theta\);
(E2) if \(\Delta\models\exists x\,\chi\) for a formula \(\chi\) in which
\(c\) does not occur, and \(\Delta\cup\{\chi(c/x)\}\models\theta\), then
\(\Delta\models\exists y\,\hat\theta\): for \(\mathfrak B\models\Delta\)
pick \(b\) with \(\mathfrak B\models\chi[x\mapsto b]\); then
\(\mathfrak B'\models\chi[x\mapsto b]\) by (a), so
\(\mathfrak B'\models\chi(c/x)\) by (S), hence \(\mathfrak B'\models\theta\)
and \(\mathfrak B\models\hat\theta[y\mapsto b]\).  Since
\((\neg\theta)^{\pi}[y/c]=\neg\hat\theta\), (E1) applied to \(\neg\theta\)
gives: if \(\Delta\models\neg\theta\) then \(\Delta\models\forall
y\,\neg\hat\theta\), so \(\Delta\models\neg\exists y\,\hat\theta\) and, as
domains are nonempty, \(\Delta\models\neg\forall y\,\hat\theta\).

\emph{\(\mathcal S\) is a consistency property for \(K\cup C\).}  Let
\(s=\Gamma_1\cup\Gamma_2\in\mathcal S\) with \((\Gamma_1,\Gamma_2)\)
inseparable.  In each case below the new set is shown to be
\(\Gamma'_1\cup\Gamma'_2\) for a pair \((\Gamma'_1,\Gamma'_2)\), and a
separator \(\theta\) of the new pair is turned into a separator of
\((\Gamma_1,\Gamma_2)\), which is impossible.  The added sentences are
n.n.f. sentences of the language of the member they are added to, and they
add at most one constant to \(C(s)\), so the new pair is a pair.

(C0) Let \(\theta\) be atomic.  If \(\theta,\neg\theta\in\Gamma_1\),
\(\Gamma_1\) has no model and \(\bot=\neg\exists x\,(x=x)\) separates.  If
\(\theta\in\Gamma_1\) and \(\neg\theta\in\Gamma_2\), the relation symbol of
\(\theta\), if any, lies in \(K_1\cap K_2\), its constants in \(C(s)\), and
\(\theta\) separates.  The other two cases follow by symmetry.

(C1) Add \(c=c\) to \(\Gamma_1\).  \(\Gamma_1\cup\{c=c\}\models\theta\)
gives \(\Gamma_1\models\theta\), as \(c=c\) holds in every structure.  If
\(c\in C(s)\), \(\theta\) separates \((\Gamma_1,\Gamma_2)\).  If not, (E1)
for \(\Delta=\Gamma_1\) and the last remark of (E) for \(\Delta=\Gamma_2\)
show that \(\forall y\,\hat\theta\) separates it.

(C2) Let \(\varphi(c/v)\in\Gamma_1\), \((c=d)\in s\), and add
\(\varphi(d/v)\) to \(\Gamma_1\); \(C\) gains nothing.  If \(c^{\mathfrak
B}=d^{\mathfrak B}\), then \(\mathfrak B\models\varphi(c/v)\) iff
\(\mathfrak B\models\varphi(d/v)\), by (S).  If \((c=d)\in\Gamma_1\), then
\(\Gamma_1\models\varphi(d/v)\), so \(\Gamma_1\models\theta\) and
\(\theta\) separates the old pair.  If \((c=d)\in\Gamma_2\), then
\(\neg(c=d)\lor\theta\) separates it: a model of \(\Gamma_1\) with
\(c=d\) satisfies \(\varphi(d/v)\), hence \(\theta\); and
\(\Gamma_2\models(c=d)\land\neg\theta\).  The case
\(\varphi(c/v)\in\Gamma_2\) follows by symmetry.

(C3) Let \(\bigwedge\Phi\in\Gamma_1\) and \(\chi\in\Phi\).  Then
\(\Gamma_1\models\chi\), and \(C(s\cup\{\chi\})=C(s)\); a separator of the
new pair separates the old one.

(C4) Let \(\forall x\,\chi\in\Gamma_1\), \(c\in C\), and add
\(\chi(c/x)\).  By (S), \(\Gamma_1\models\chi(c/x)\), so
\(\Gamma_1\models\theta\).  If \(c\in C(s)\), \(\theta\) separates the old
pair; otherwise \(\forall y\,\hat\theta\) does, by (E1) and the last
remark of (E).

(C5) Let \(\bigvee\Phi\in\Gamma_1\), and suppose that for every
\(\chi\in\Phi\) the pair \((\Gamma_1\cup\{\chi\},\Gamma_2)\) has a separator
\(\theta_\chi\); choose one for each \(\chi\) (\(\Phi\) is countable).  As
\(C(s\cup\{\chi\})=C(s)\), \(\theta=\bigvee\{\theta_\chi:\chi\in\Phi\}\) is
a sentence with symbols in \((K_1\cap K_2)\cup C(s)\).  A model of
\(\Gamma_1\) satisfies some \(\chi\in\Phi\), hence \(\theta_\chi\) and
\(\theta\); a model of \(\Gamma_2\) satisfies every \(\neg\theta_\chi\),
hence \(\neg\theta\).  So \(\theta\) separates the old pair.  Hence some
\(\chi\) gives an inseparable pair.

(C6) Let \(\exists x\,\chi\in\Gamma_1\), and take \(c\in C\setminus
C(s)\), possible since \(C(s)\) is finite; \(c\) does not occur in
\(\chi\).  If \(\theta\) separates \((\Gamma_1\cup\{\chi(c/x)\},\Gamma_2)\),
then (E2) for \(\Delta=\Gamma_1\) gives \(\Gamma_1\models\exists
y\,\hat\theta\), and the last remark of (E) for \(\Delta=\Gamma_2\) gives
\(\Gamma_2\models\neg\exists y\,\hat\theta\); so \(\exists y\,\hat\theta\)
separates the old pair.

\emph{Conclusion of Step 1.}  By \cref{thm:p02-w02c-p02-003-nnf},
\(\varphi^{+}\) and \(\psi^{\ominus}\) are n.n.f. sentences of \(K_1\) and
\(K_2\) equivalent to \(\varphi\) and \(\neg\psi\).  If
\((\{\varphi^{+}\},\{\psi^{\ominus}\})\) were inseparable,
\(\{\varphi^{+},\psi^{\ominus}\}\in\mathcal S\) would have a model by
\cref{thm:p02-w02c-p02-003-model-existence}, a model of \(\varphi\) and
\(\neg\psi\), against the validity of \(\varphi\to\psi\).  So it has a
separator \(\theta\); no constant of \(C\) occurs in \(\varphi^{+}\) or
\(\psi^{\ominus}\), so the symbols of \(\theta\) lie in \(K_1\cap K_2\).
\(\varphi\models\theta\), so \(\varphi\to\theta\) is valid; a model of
\(\theta\) does not satisfy \(\psi^{\ominus}\), hence satisfies \(\psi\), so
\(\theta\to\psi\) is valid.

\emph{Step 2: function and constant symbols.}  Let \(\varphi\), \(\psi\) be
sentences for \(L\), \(\varphi\to\psi\) valid.  For a function symbol
\(f\) of arity \(n\) and a constant symbol \(c\), let \(\mathrm{Fn}_f\) be
\(\forall x_1\cdots\forall x_n\,\exists y\,(R_f(x_1,\ldots,x_n,y)\land
\forall y'\,(R_f(x_1,\ldots,x_n,y')\to y'=y))\) and \(\mathrm{Fn}_c\) be
\(\exists y\,(P_c(y)\land\forall y'\,(P_c(y')\to y'=y))\), sentences of
\(L^{r}\) (\cref{def:p02-w02c-p02-003-relational-image}).  For a sentence
\(\chi\) of \(L\), let \(\mathrm{Fn}(\chi)\) be the conjunction of
\(\mathrm{Fn}_e\) over the countably many function and constant symbols
\(e\) occurring in \(\chi\) (the empty conjunction, true in every
structure, if there is none).  Put \(\varphi_1=\varphi^{u}\land
\mathrm{Fn}(\varphi)\) and \(\psi_1=\neg\mathrm{Fn}(\psi)\lor\psi^{u}\),
relational sentences of \(L_{\omega_1\omega}\) by part~3 of
\cref{thm:p02-w02c-p02-003-translations}.  The relation symbols occurring
in \(\varphi_1\) are the \(R\in L\), \(R_f\) and \(P_c\) for \(R\), \(f\),
\(c\) occurring in \(\varphi\), and likewise for \(\psi_1\).

\(\varphi_1\to\psi_1\) is valid.  Let \(\mathfrak B\) satisfy \(\varphi_1\)
and \(\mathrm{Fn}(\psi)\), let \(L_0\) be the set of symbols occurring in
\(\varphi\) or \(\psi\), and \(\mathfrak B_0=\mathfrak B\restriction
L_0^{r}\).  Define \(\mathfrak A\) for \(L_0\) on the domain of
\(\mathfrak B\): \(R^{\mathfrak A}=R^{\mathfrak B}\),
\(f^{\mathfrak A}(\bar a)\) the unique \(b\) with
\((\bar a,b)\in R_f^{\mathfrak B}\), and \(c^{\mathfrak A}\) the unique
member of \(P_c^{\mathfrak B}\), which exist by \(\mathrm{Fn}_f\),
\(\mathrm{Fn}_c\).  Then \(R_f^{\mathfrak A^{r}}\), the graph of
\(f^{\mathfrak A}\), is \(R_f^{\mathfrak B}\), and
\(P_c^{\mathfrak A^{r}}=P_c^{\mathfrak B}\); so \(\mathfrak A^{r}=\mathfrak
B_0\).  By (a), \(\mathfrak B_0\models\varphi^{u}\); by part~3 of
\cref{thm:p02-w02c-p02-003-translations},
\(\mathfrak A\models\varphi\); so \(\mathfrak A\models\psi\), hence
\(\mathfrak B_0\models\psi^{u}\) and \(\mathfrak B\models\psi^{u}\).

By Step 1 there is a relational interpolant \(\theta_1\) of \(\varphi_1\)
and \(\psi_1\).  Let \(\theta=\theta_1^{-}\)
(\cref{def:p02-w02c-p02-003-relational-image}).  A symbol \(R\in L\),
\(R_f\) or \(P_c\) occurring in \(\theta_1\) occurs in \(\varphi_1\) and in
\(\psi_1\), so \(R\), \(f\) or \(c\) occurs in \(\varphi\) and in \(\psi\);
these are the non-logical symbols of \(\theta\).  For a structure
\(\mathfrak A\) for \(L\), \(\mathfrak A^{r}\) satisfies every
\(\mathrm{Fn}_e\).  If \(\mathfrak A\models\varphi\), then
\(\mathfrak A^{r}\models\varphi_1\) by part~3, so
\(\mathfrak A^{r}\models\theta_1\) and \(\mathfrak A\models\theta\) by
part~1 of \cref{thm:p02-w02c-p02-003-translations}.  If
\(\mathfrak A\models\theta\), then \(\mathfrak A^{r}\models\theta_1\) by
part~1, so \(\mathfrak A^{r}\models\psi_1\); as
\(\mathfrak A^{r}\models\mathrm{Fn}(\psi)\), \(\mathfrak
A^{r}\models\psi^{u}\) and \(\mathfrak A\models\psi\) by part~3.  So
\(\varphi\to\theta\) and \(\theta\to\psi\) are valid.
\end{proof}
\end{annexitem}

\begin{annexitem}{annex:p02-w02c-p02-003-beth}{thm:p02-w02c-p02-003-beth}{Beth's theorem for \(L_{\omega_1\omega}\)}
\begin{proof}
The formulas of \(L_{\omega_1\omega}\) are generated from the atomic
formulas by \(\neg\), by \(\bigwedge\) and \(\bigvee\) of countable sets,
and by \(\exists v\), \(\forall v\); each formula is obtained in a
well-founded way from its immediate subformulas, so properties of all
formulas can be proved by induction on this generation.  Three auxiliary
facts are proved first.

\emph{(a) Renaming and reducts.}  Let \(K\) and \(K'\) be languages and
\(\rho\) an injective map from the symbols of \(K\) to those of \(K'\)
that sends relation, function and constant symbols to symbols of the same
kind and arity.  For a formula \(\chi\) for \(K\), \(\rho\chi\) replaces
each symbol \(e\) of \(K\) by \(\rho(e)\).  For a structure \(\mathfrak B\)
for \(K'\), \(\mathfrak B_\rho\) is the structure for \(K\) with the domain
of \(\mathfrak B\) and \(e^{\mathfrak B_\rho}=\rho(e)^{\mathfrak B}\).
Claim: \(\mathfrak B\models\rho\chi[s]\) iff
\(\mathfrak B_\rho\models\chi[s]\) for every assignment \(s\).  For a term
\(t\), the value of \(\rho t\) in \(\mathfrak B\) at \(s\) equals the value
of \(t\) in \(\mathfrak B_\rho\) at \(s\), by induction on \(t\): a variable
\(v\) takes the value \(s(v)\) in both, and
\((\rho f)^{\mathfrak B}=f^{\mathfrak B_\rho}\),
\((\rho c)^{\mathfrak B}=c^{\mathfrak B_\rho}\).  Hence an atomic formula
\(R(t_1,\ldots,t_k)\) holds at \(s\) in \(\mathfrak B_\rho\) iff the values
of the \(t_i\) lie in \(R^{\mathfrak B_\rho}=(\rho R)^{\mathfrak B}\) iff
\(\rho R(\rho t_1,\ldots,\rho t_k)\) holds at \(s\) in \(\mathfrak B\); an
equation is treated in the same way.  The clauses for \(\neg\),
\(\bigwedge\), \(\bigvee\), \(\exists v\), \(\forall v\) refer only to the
satisfaction of immediate subformulas at \(s\) or at \(s(b/v)\) for
\(b\) in the common domain, so the claim passes from the immediate
subformulas to the formula.  When \(K\subseteq K'\) and \(\rho\) is the
inclusion, \(\rho\chi=\chi\) and \(\mathfrak B_\rho\) is the reduct
\(\mathfrak B\restriction K\): a formula for \(K\) has the same
satisfaction in \(\mathfrak B\) and in \(\mathfrak B\restriction K\).
\emph{(b) Renaming variables.}  Let \(\pi\) be an injective map from the
set of variables into itself, and let \(\chi^{\pi}\) replace every
occurrence of every variable \(v\) in \(\chi\), including those next to a
quantifier, by \(\pi(v)\).  Claim: \(\mathfrak M\models\chi^{\pi}[s]\) iff
\(\mathfrak M\models\chi[s\circ\pi]\).  For terms, the value of \(t^{\pi}\)
at \(s\) is the value of \(t\) at \(s\circ\pi\), since a variable \(v\)
becomes \(\pi(v)\) with value \(s(\pi(v))\).  Atomic formulas follow, and
the connectives are immediate.  For \(\exists v\,\chi\):
\(\mathfrak M\models\exists\pi(v)\,\chi^{\pi}[s]\) iff some \(b\) has
\(\mathfrak M\models\chi^{\pi}[s(b/\pi(v))]\) iff, by the induction
hypothesis, some \(b\) has \(\mathfrak M\models\chi[s(b/\pi(v))\circ\pi]\).
Now \(s(b/\pi(v))\circ\pi=(s\circ\pi)(b/v)\): both send \(v\) to \(b\),
and for \(w\neq v\) injectivity gives \(\pi(w)\neq\pi(v)\), so both send
\(w\) to \(s(\pi(w))\).  So the condition is
\(\mathfrak M\models\exists v\,\chi[s\circ\pi]\); \(\forall\) is dual.  If
\(\chi\) is a sentence, \(\chi^{\pi}\) is a sentence (its free variables
are the images of those of \(\chi\)), and the claim shows that the two
sentences hold in the same structures.

\emph{(c) Constants to variables.}  Let \(a_1,\ldots,a_n\) be distinct
constant symbols not in \(K\), let \(\chi\) be a formula for
\(K\cup\{a_1,\ldots,a_n\}\), and let \(x_1,\ldots,x_n\) be distinct
variables none of which occurs in \(\chi\).  Let \(\chi[\bar x/\bar a]\)
replace each occurrence of \(a_i\) by \(x_i\).  Claim: for a structure
\(\mathfrak M\) for \(K\), elements \(c_1,\ldots,c_n\), and an assignment
\(s\) with \(s(x_i)=c_i\) for each \(i\),
\((\mathfrak M,c_1,\ldots,c_n)\models\chi[s]\) iff
\(\mathfrak M\models\chi[\bar x/\bar a][s]\).  For terms, the value of
\(t[\bar x/\bar a]\) in \(\mathfrak M\) at \(s\) equals that of \(t\) in
\((\mathfrak M,\bar c)\) at \(s\), because \(a_i\) denotes \(c_i=s(x_i)\);
atomic formulas and connectives follow.  For \(\exists v\,\chi\), the
variable \(v\) is not any \(x_i\), since the \(x_i\) do not occur in
\(\chi\); so \(s(b/v)(x_i)=c_i\) and the induction hypothesis applies at
\(s(b/v)\).  If \(\chi\) is a sentence, the free variables of
\(\chi[\bar x/\bar a]\) are among \(x_1,\ldots,x_n\).

\emph{The theorem.}  Let \(P_1,P_2\) be distinct \(n\)-ary relation
symbols and \(a_1,\ldots,a_n\) distinct constant symbols, none in
\(L\cup\{P\}\).  Let \(\rho_i\) send \(P\) to \(P_i\) and fix every symbol
of \(L\), and write \(\varphi(P_i)\) for \(\rho_i\varphi(P)\) and
\(P_i\bar a\) for \(P_i(a_1,\ldots,a_n)\).  We show that
\[
 \bigl(\varphi(P_1)\land P_1\bar a\bigr)\to
 \bigl(\varphi(P_2)\to P_2\bar a\bigr)
\]
is valid.  Let \(\mathfrak B\) be a structure for
\(L\cup\{P_1,P_2,a_1,\ldots,a_n\}\) with
\(\mathfrak B\models\varphi(P_1)\), \(\mathfrak B\models P_1\bar a\) and
\(\mathfrak B\models\varphi(P_2)\), and let
\(\mathfrak M=\mathfrak B\restriction L\).  By (a),
\(\mathfrak B_{\rho_i}=(\mathfrak M,P_i^{\mathfrak B})\) satisfies
\(\varphi(P)\) for \(i=1,2\).  By the hypothesis of implicit
definability, \(P_1^{\mathfrak B}=P_2^{\mathfrak B}\).  Since
\((a_1^{\mathfrak B},\ldots,a_n^{\mathfrak B})\in P_1^{\mathfrak B}\), it
lies in \(P_2^{\mathfrak B}\), so \(\mathfrak B\models P_2\bar a\).

By \cref{thm:p02-w02c-p02-003-interpolation} there is a sentence
\(\theta_0\) of \(L_{\omega_1\omega}\) such that
\(\varphi(P_1)\land P_1\bar a\to\theta_0\) and
\(\theta_0\to(\varphi(P_2)\to P_2\bar a)\) are valid and every non-logical
symbol of \(\theta_0\) occurs on both sides.  \(P_1\) does not occur on
the right and \(P_2\) not on the left, so the symbols of \(\theta_0\) are
among the symbols of \(L\) occurring in \(\varphi(P)\) and
\(a_1,\ldots,a_n\).

The set of variables is infinite, so it contains a sequence
\(u_0,u_1,\ldots\) of distinct variables with \(u_{i-1}=x_i\) for
\(1\leq i\leq n\) (if the variables are not listed as \(v_0,v_1,\ldots\),
the axiom of choice gives such a sequence).  Let \(\pi(u_k)=u_{k+n}\) and
\(\pi(w)=w\) for every other variable \(w\).  Then \(\pi\) is injective
and no \(x_i\) is in its range, so no \(x_i\) occurs in
\(\theta_1=\theta_0^{\pi}\), and by (b) \(\theta_1\) is a sentence that
holds in exactly the structures in which \(\theta_0\) holds.  Let
\(\theta=\theta_1[\bar x/\bar a]\).  By (c), \(\theta\) is a formula for
\(L\) with free variables among \(x_1,\ldots,x_n\), and its non-logical
symbols are those of \(\theta_0\) other than the \(a_i\).

Let \((\mathfrak M,R)\models\varphi(P)\), where \(\mathfrak M\) is a
structure for \(L\), and let \(c_1,\ldots,c_n\) be elements and \(s\) an
assignment with \(s(x_i)=c_i\).  Let \(\mathfrak B\) be the structure for
\(L\cup\{P_1,P_2,a_1,\ldots,a_n\}\) expanding \(\mathfrak M\) by
\(P_1^{\mathfrak B}=P_2^{\mathfrak B}=R\) and \(a_i^{\mathfrak B}=c_i\).
Then \(\mathfrak B_{\rho_i}=(\mathfrak M,R)\), so by (a)
\(\mathfrak B\models\varphi(P_1)\) and \(\mathfrak B\models\varphi(P_2)\).
By (a) applied to the inclusion of the language of \(\theta_1\),
\(\mathfrak B\models\theta_1\) iff
\((\mathfrak M,\bar c)\models\theta_1\), and by (c) this holds iff
\(\mathfrak M\models\theta[s]\).  If \((c_1,\ldots,c_n)\in R\), then
\(\mathfrak B\models P_1\bar a\), so \(\mathfrak B\models\theta_0\) by the
first validity, hence \(\mathfrak B\models\theta_1\) and
\(\mathfrak M\models\theta[s]\).  Conversely, if
\(\mathfrak M\models\theta[s]\), then \(\mathfrak B\models\theta_1\), so
\(\mathfrak B\models\theta_0\); with \(\mathfrak B\models\varphi(P_2)\)
the second validity gives \(\mathfrak B\models P_2\bar a\), that is
\((c_1,\ldots,c_n)\in R\).  By (a) for the inclusion of \(L\) in
\(L\cup\{P\}\), \(\mathfrak M\models\theta[s]\) iff
\((\mathfrak M,R)\models\theta[s]\).  As \(c_1,\ldots,c_n\) were
arbitrary, \((\mathfrak M,R)\models\forall x_1\cdots\forall
x_n\,(P(x_1,\ldots,x_n)\leftrightarrow\theta)\).

For a countable set \(\Phi(P)\), the conjunction \(\bigwedge\Phi(P)\) is
a sentence of \(L_{\omega_1\omega}\) with the same models as \(\Phi(P)\),
so the hypothesis and the conclusion for \(\Phi(P)\) are those for this
sentence.
\end{proof}
\end{annexitem}

\begin{annexitem}{annex:p02-w02c-p02-003-hinf-sound}{thm:p02-w02c-p02-003-hinf-soundness}{Soundness}
\begin{proof}
\emph{(S\('\)) Substitution of a closed term.}  For a formula \(\chi\), a
variable \(x\), a closed term \(t\), a structure \(\mathfrak M\) and an
assignment \(s\): \(\mathfrak M\models\chi(t/x)[s]\) iff
\(\mathfrak M\models\chi[s(t^{\mathfrak M}/x)]\), where \(t^{\mathfrak M}\)
is the value of \(t\), which does not depend on the assignment.  For a term
\(u\), the value of \(u(t/x)\) at \(s\) equals the value of \(u\) at
\(s(t^{\mathfrak M}/x)\), by induction on \(u\): for \(u=x\) both are
\(t^{\mathfrak M}\); for a variable \(y\neq x\) both are \(s(y)\); for a
constant \(c\) both are \(c^{\mathfrak M}\); for \(f(u_1,\ldots,u_k)\) both
are \(f^{\mathfrak M}\) of the common values of the \(u_i\).  Atomic formulas
follow, and \(\neg\), \(\bigwedge\), \(\bigvee\) commute with the
substitution and refer to the same assignment.  For \(\exists y\,\chi_1\)
with \(y=x\), there is no free \(x\), the formula is unchanged, and since
\(s(b/x)=s(t^{\mathfrak M}/x)(b/x)\) for every \(b\) it holds at \(s\) iff
it holds at \(s(t^{\mathfrak M}/x)\).  For \(y\neq x\),
\((\exists y\,\chi_1)(t/x)=\exists y\,\chi_1(t/x)\), as \(t\) has no
variable; it holds at \(s\) iff \(\chi_1(t/x)\) holds at \(s(b/y)\) for
some \(b\), iff (induction hypothesis) \(\chi_1\) holds at
\(s(b/y)(t^{\mathfrak M}/x)=s(t^{\mathfrak M}/x)(b/y)\) for some \(b\), iff
\(\exists y\,\chi_1\) holds at \(s(t^{\mathfrak M}/x)\).  \(\forall y\) is
the same computation with ``for every \(b\)''.  Substituting closed terms
for several distinct variables one after another gives the same result in
any order, since a closed term contains no variable.

\emph{Part 1.}  Let \(\Sigma\) be the set of sentences of
\(L_{\omega_1\omega}\) for \(L\cup\mathrm{Par}\) that hold in every
structure for \(L\cup\mathrm{Par}\) satisfying every member of \(T\).  It
suffices to show that \(\Sigma\) contains the axioms and \(T\) and is closed
under (R1)--(R3), for then it contains every theorem of \(T\).  Members of
\(T\) lie in \(\Sigma\) by definition.  Let \(\mathfrak M\) satisfy every
member of \(T\).

(A1): let \(v\) give a prime sentence the value \(\mathbf t\) iff
\(\mathfrak M\) satisfies it.  By induction on the generation of a sentence
from prime sentences by \(\neg\), \(\bigwedge\), \(\bigvee\) (the immediate
subformulas of a sentence of one of these three forms are sentences), the
extended value of every sentence \(\chi\) is \(\mathbf t\) iff
\(\mathfrak M\models\chi\): the three clauses of the extension are the
satisfaction clauses.  So a tautology, true under \(v\), holds in
\(\mathfrak M\).  (A2): \(t=t\) holds.  If \(t^{\mathfrak M}=t'^{\mathfrak M}\)
and \(\mathfrak M\models\alpha(t/x)\), then by (S\('\))
\(\mathfrak M\models\alpha[s(t^{\mathfrak M}/x)]\), which is
\(\alpha[s(t'^{\mathfrak M}/x)]\), so \(\mathfrak M\models\alpha(t'/x)\).
(A3): if \(\mathfrak M\models\forall x\,\chi\), then
\(\mathfrak M\models\chi[s(t^{\mathfrak M}/x)]\), so
\(\mathfrak M\models\chi(t/x)\) by (S\('\)); if
\(\mathfrak M\models\chi(t/x)\), then \(t^{\mathfrak M}\) witnesses
\(\exists x\,\chi\).

(R1) and (R3) preserve truth in \(\mathfrak M\): if \(\theta\) holds and
\(\theta\to\chi\) holds for every \(\chi\in\Phi\), every \(\chi\in\Phi\)
holds, so \(\bigwedge\Phi\) holds.  (R2), first form: let
\(\theta\to\chi(p/x)\in\Sigma\) with \(p\) in neither \(\theta\) nor
\(\chi\), and suppose \(\mathfrak M\models\theta\).  For an element \(b\),
let \(\mathfrak M_b\) differ from \(\mathfrak M\) only in
\(p^{\mathfrak M_b}=b\).  The two structures have the same reduct to the
language without \(p\), and \(p\) occurs neither in \(\theta\), nor in
\(\chi\), nor in \(T\) (no parameter occurs in \(T\)); by fact (a) of
\cref{annex:p02-w02c-p02-003-beth}, \(\mathfrak M_b\) satisfies every
member of \(T\) and \(\theta\).  So \(\mathfrak M_b\models\chi(p/x)\), by
(S\('\)) \(\mathfrak M_b\models\chi[s(b/x)]\), and by fact (a)
\(\mathfrak M\models\chi[s(b/x)]\).  As \(b\) was arbitrary,
\(\mathfrak M\models\forall x\,\chi\).  Second form: let
\(\chi(p/x)\to\theta\in\Sigma\) with \(p\) in neither, and
\(\mathfrak M\models\exists x\,\chi\), witnessed by \(b\).  As before,
\(\mathfrak M_b\) satisfies \(T\) and, by fact (a) and (S\('\)),
\(\chi(p/x)\); so \(\mathfrak M_b\models\theta\), and
\(\mathfrak M\models\theta\) by fact (a).

\emph{Part 2.}  Let \(\Sigma'\) be the set of sentences \(\theta\) with
\(T_0\vdash_{\infty}\theta\) for some countable \(T_0\subseteq T\).  If
\(T_1\subseteq T_2\), every theorem of \(T_1\) is a theorem of \(T_2\),
since the theorems of \(T_2\) contain the axioms and \(T_1\) and are closed
under the rules.  Axioms lie in \(\Sigma'\) with \(T_0=\varnothing\), and
\(\theta\in T\) with \(T_0=\{\theta\}\).  (R1): from countable \(T_1\), \(T_2\)
for the two premises, \(T_1\cup T_2\) serves.  (R2): the theorems of a
fixed \(T_0\) are closed under (R2), and its proviso does not mention the
premises.  (R3): choose for each \(\chi\in\Phi\) a countable
\(T_\chi\subseteq T\) with \(T_\chi\vdash_{\infty}\theta\to\chi\); the
union of the countably many countable sets \(T_\chi\) is countable (axiom
of choice) and proves every premise, hence the conclusion.  So
\(\Sigma'\) contains every theorem of \(T\).
\end{proof}
\end{annexitem}

\begin{annexitem}{annex:p02-w02c-p02-003-hinf-nnf}{thm:p02-w02c-p02-003-hinf-nnf}{Derived rules and negation normal form in \(\mathbf H_{\infty}\)}
\begin{proof}
Each sentence said to be a tautology below is checked by assuming its
consequent false and computing values under the clauses of
\cref{def:p02-w02c-p02-003-hinf}.

\emph{Part 1.}  The tautology is an axiom (A1); \(n\) applications of
(R1) give \(\theta\).  This is \emph{tautological consequence}.

\emph{Part 2.}  For \(\chi\in\Psi\),
\((\chi\to\theta)\to(\neg\theta\to\neg\chi)\) is a tautology, so
\(T\vdash_{\infty}\neg\theta\to\neg\chi\) by part~1.  (R3), applied to the
countable set \(\{\neg\chi:\chi\in\Psi\}\), gives
\(T\vdash_{\infty}\neg\theta\to\bigwedge\{\neg\chi:\chi\in\Psi\}\).  The
sentence \((\neg\theta\to\bigwedge\{\neg\chi:\chi\in\Psi\})\to
(\bigvee\Psi\to\theta)\) is a tautology: if \(\bigvee\Psi\) has value
\(\mathbf t\) and \(\theta\) value \(\mathbf f\), some \(\chi\in\Psi\) has
value \(\mathbf t\), so the conjunction has value \(\mathbf f\) while
\(\neg\theta\) has value \(\mathbf t\).  Part~1 gives
\(T\vdash_{\infty}\bigvee\Psi\to\theta\).  For the second claim let
\(\top=\bigwedge\varnothing\), which has value \(\mathbf t\) under every
truth assignment and so is an axiom.  For \(\theta\in T\),
\(\theta\to(\top\to\theta)\) is a tautology, so
\(T\vdash_{\infty}\top\to\theta\); (R3) gives
\(T\vdash_{\infty}\top\to\bigwedge T\), and (R1) gives
\(T\vdash_{\infty}\bigwedge T\).

\emph{Part 3.}  For sentences \(A,B\), \(T\vdash_{\infty}A\leftrightarrow B\)
holds iff \(T\vdash_{\infty}A\to B\) and \(T\vdash_{\infty}B\to A\), by
part~1 and the tautologies \((A\to B)\to(B\to A)\to(A\leftrightarrow B)\),
\((A\leftrightarrow B)\to(A\to B)\), \((A\leftrightarrow B)\to(B\to A)\).
The immediate subformulas of a formula with finitely many free variables
and finitely many parameters again have finitely many of each (the free
variables of a member of \(\Phi\) are among those of \(\bigwedge\Phi\) or
\(\bigvee\Phi\), and those of \(\chi_1\) are among those of
\(\exists x\,\chi_1\), \(\forall x\,\chi_1\) together with \(x\)), so the
claim can be proved by induction on the generation of \(\chi\), for all
\(\sigma\) at once.  For a member \(\varphi\) of \(\Phi\),
\(\sigma_\varphi\) is the restriction of \(\sigma\) to the free variables of
\(\varphi\), so that \((\bigwedge\Phi)\sigma=\bigwedge\{\varphi\sigma_\varphi:
\varphi\in\Phi\}\).  For \(\chi=\exists x\,\chi_1\) or \(\forall x\,\chi_1\),
\(\sigma'\) is the restriction of \(\sigma\) to the free variables of
\(\chi_1\) other than \(x\), so that \(\chi\sigma\) is
\(\exists x\,(\chi_1\sigma')\) or \(\forall x\,(\chi_1\sigma')\); a parameter
\(p\) is chosen that occurs neither in \(\chi\) nor in any value of
\(\sigma\) (finitely many parameters are excluded), and
\(\sigma''=\sigma'\cup\{(x,p)\}\), so that
\((\chi_1\sigma')(p/x)=\chi_1\sigma''\).  By \cref{thm:p02-w02c-p02-003-nnf},
\(\chi^{+}\) and \(\chi^{\ominus}\) have the non-logical symbols of
\(\chi\), so \(p\) occurs in no sentence below except where it is
substituted.  All steps named ``part~1'' combine displayed theorems by a
tautology of the form \((A\to B)\to(B\to C)\to(A\to C)\) or
\((A\to B)\to(\neg B\to\neg A)\).

\emph{Atomic \(\chi\).}  \(\chi^{+}\sigma=\chi\sigma\) and
\(\chi^{\ominus}\sigma=\neg\chi\sigma\); \(A\leftrightarrow A\) is a
tautology.

\emph{\(\chi=\neg\chi_1\).}  \(\chi^{+}\sigma=\chi_1^{\ominus}\sigma\), and
the induction hypothesis gives
\(\chi_1^{\ominus}\sigma\leftrightarrow\neg\chi_1\sigma\), which is the
claim.  \(\chi^{\ominus}\sigma=\chi_1^{+}\sigma\); from
\(\chi_1^{+}\sigma\leftrightarrow\chi_1\sigma\) and the tautology
\(\chi_1\sigma\leftrightarrow\neg\neg\chi_1\sigma\), part~1 gives
\(\chi_1^{+}\sigma\leftrightarrow\neg\chi\sigma\).

\emph{\(\chi=\bigwedge\Phi\).}  \(\chi^{+}\sigma=\bigwedge\{\varphi^{+}
\sigma_\varphi\}\) and \(\chi^{\ominus}\sigma=\bigvee\{\varphi^{\ominus}
\sigma_\varphi\}\), over \(\varphi\in\Phi\).  For each \(\varphi\),
\(\chi^{+}\sigma\to\varphi^{+}\sigma_\varphi\) is a tautology and the
induction hypothesis gives \(\varphi^{+}\sigma_\varphi\to\varphi\sigma_\varphi\);
so \(T\vdash_{\infty}\chi^{+}\sigma\to\varphi\sigma_\varphi\), and (R3) gives
\(\chi^{+}\sigma\to\chi\sigma\).  Exchanging the roles of \(\varphi^{+}
\sigma_\varphi\) and \(\varphi\sigma_\varphi\) gives \(\chi\sigma\to
\chi^{+}\sigma\).  Next, \(\varphi^{\ominus}\sigma_\varphi\to\neg\varphi
\sigma_\varphi\) (induction hypothesis) and the tautology
\(\neg\varphi\sigma_\varphi\to\neg\chi\sigma\) give
\(\varphi^{\ominus}\sigma_\varphi\to\neg\chi\sigma\) for every \(\varphi\),
and part~2 gives \(\chi^{\ominus}\sigma\to\neg\chi\sigma\).  Conversely,
\(\neg\chi^{\ominus}\sigma\to\neg\varphi^{\ominus}\sigma_\varphi\) is a
tautology, and \(\neg\varphi^{\ominus}\sigma_\varphi\to\varphi\sigma_\varphi\)
follows from the induction hypothesis \(\neg\varphi\sigma_\varphi\to
\varphi^{\ominus}\sigma_\varphi\) by part~1; so
\(\neg\chi^{\ominus}\sigma\to\varphi\sigma_\varphi\) for every \(\varphi\),
(R3) gives \(\neg\chi^{\ominus}\sigma\to\chi\sigma\), and part~1 gives
\(\neg\chi\sigma\to\chi^{\ominus}\sigma\).

\emph{\(\chi=\bigvee\Phi\).}  \(\chi^{+}\sigma=\bigvee\{\varphi^{+}
\sigma_\varphi\}\) and \(\chi^{\ominus}\sigma=\bigwedge\{\varphi^{\ominus}
\sigma_\varphi\}\).  For each \(\varphi\), \(\varphi^{+}\sigma_\varphi\to
\varphi\sigma_\varphi\) (induction hypothesis) and the tautology
\(\varphi\sigma_\varphi\to\chi\sigma\) give \(\varphi^{+}\sigma_\varphi\to
\chi\sigma\), and part~2 gives \(\chi^{+}\sigma\to\chi\sigma\); exchanging
roles gives \(\chi\sigma\to\chi^{+}\sigma\).  The tautology
\(\chi^{\ominus}\sigma\to\varphi^{\ominus}\sigma_\varphi\) and
\(\varphi^{\ominus}\sigma_\varphi\to\neg\varphi\sigma_\varphi\) give
\(\chi^{\ominus}\sigma\to\neg\varphi\sigma_\varphi\) for every \(\varphi\);
(R3) gives \(\chi^{\ominus}\sigma\to\bigwedge\{\neg\varphi\sigma_\varphi\}\),
and the tautology \(\bigwedge\{\neg\varphi\sigma_\varphi\}\to\neg\chi\sigma\)
gives \(\chi^{\ominus}\sigma\to\neg\chi\sigma\).  Conversely, the tautology
\(\neg\chi\sigma\to\neg\varphi\sigma_\varphi\) and
\(\neg\varphi\sigma_\varphi\to\varphi^{\ominus}\sigma_\varphi\) give
\(\neg\chi\sigma\to\varphi^{\ominus}\sigma_\varphi\) for every \(\varphi\),
and (R3) gives \(\neg\chi\sigma\to\chi^{\ominus}\sigma\).

\emph{\(\chi=\forall x\,\chi_1\).}  \(\chi^{+}\sigma=\forall x\,
(\chi_1^{+}\sigma')\) and \(\chi^{\ominus}\sigma=\exists x\,
(\chi_1^{\ominus}\sigma')\).  (A3) gives \(\chi^{+}\sigma\to\chi_1^{+}
\sigma''\) and the induction hypothesis \(\chi_1^{+}\sigma''\to\chi_1
\sigma''\); so \(\chi^{+}\sigma\to(\chi_1\sigma')(p/x)\), and (R2), first
form, gives \(\chi^{+}\sigma\to\chi\sigma\).  The converse is the same with
the roles of \(\chi_1^{+}\) and \(\chi_1\) exchanged.  Next,
\(\chi_1^{\ominus}\sigma''\to\neg\chi_1\sigma''\) (induction hypothesis)
and \(\chi\sigma\to\chi_1\sigma''\) (A3) give, by part~1,
\(\chi_1^{\ominus}\sigma''\to\neg\chi\sigma\), that is
\((\chi_1^{\ominus}\sigma')(p/x)\to\neg\chi\sigma\); (R2), second form,
gives \(\chi^{\ominus}\sigma\to\neg\chi\sigma\).  Conversely,
\(\chi_1^{\ominus}\sigma''\to\chi^{\ominus}\sigma\) (A3) and
\(\neg\chi_1\sigma''\to\chi_1^{\ominus}\sigma''\) give
\(\neg\chi^{\ominus}\sigma\to\chi_1\sigma''\) by part~1; (R2), first form,
gives \(\neg\chi^{\ominus}\sigma\to\chi\sigma\), and part~1 gives
\(\neg\chi\sigma\to\chi^{\ominus}\sigma\).

\emph{\(\chi=\exists x\,\chi_1\).}  \(\chi^{+}\sigma=\exists x\,
(\chi_1^{+}\sigma')\) and \(\chi^{\ominus}\sigma=\forall x\,
(\chi_1^{\ominus}\sigma')\).  The induction hypothesis
\(\chi_1^{+}\sigma''\to\chi_1\sigma''\) and (A3)
\(\chi_1\sigma''\to\chi\sigma\) give \((\chi_1^{+}\sigma')(p/x)\to
\chi\sigma\), and (R2), second form, gives \(\chi^{+}\sigma\to\chi\sigma\);
the converse exchanges roles.  Next, \(\chi_1\sigma''\to\neg\chi_1^{\ominus}
\sigma''\) (from the induction hypothesis by part~1) and
\(\neg\chi_1^{\ominus}\sigma''\to\neg\chi^{\ominus}\sigma\) (from (A3)
\(\chi^{\ominus}\sigma\to\chi_1^{\ominus}\sigma''\) by part~1) give
\((\chi_1\sigma')(p/x)\to\neg\chi^{\ominus}\sigma\); (R2), second form,
gives \(\chi\sigma\to\neg\chi^{\ominus}\sigma\), and part~1 gives
\(\chi^{\ominus}\sigma\to\neg\chi\sigma\).  Conversely,
\(\neg\chi\sigma\to\neg\chi_1\sigma''\) (from (A3) \(\chi_1\sigma''\to
\chi\sigma\)) and \(\neg\chi_1\sigma''\to\chi_1^{\ominus}\sigma''\) give
\(\neg\chi\sigma\to(\chi_1^{\ominus}\sigma')(p/x)\), and (R2), first form,
gives \(\neg\chi\sigma\to\chi^{\ominus}\sigma\).

In every use of (R2) above, the sentence written \(\theta\) in the rule is one
of \(\chi\sigma\), \(\chi^{+}\sigma\), \(\chi^{\ominus}\sigma\) or the negation
of one of them, and the formula into which \(p\) is substituted is
\(\chi_1\sigma'\), \(\chi_1^{+}\sigma'\) or \(\chi_1^{\ominus}\sigma'\); by the
choice of \(p\), it occurs in none of these.
\end{proof}
\end{annexitem}

\begin{annexitem}{annex:p02-w02c-p02-003-term-mex}{thm:p02-w02c-p02-003-term-model-existence}{Model existence with function symbols}
\begin{proof}
Write \(\mathrm{CT}\) for the set of closed terms of \(L\cup\mathrm{Par}\);
it is countable (finite strings over a countable set of symbols) and
nonempty.  For a formula \(\varphi\) and a map \(\sigma\) from its free
variables to \(\mathrm{CT}\), \(\varphi\sigma\) replaces each free
occurrence of each such variable \(x\) by \(\sigma(x)\); no variable is
captured, and (S\('\)) of \cref{annex:p02-w02c-p02-003-hinf-sound}
applies.

\emph{Hintikka sets with closed terms.}  Call a set \(H\) of n.n.f.
sentences of \(L_{\omega_1\omega}\) for \(L\cup\mathrm{Par}\) a
\emph{term Hintikka set} if (H0) no atomic \(\theta\) has \(\theta\) and
\(\neg\theta\) in \(H\); (H1) \((t=t)\in H\) for every \(t\in\mathrm{CT}\);
(H2) if \(\alpha\) is atomic with at most \(x\) free and
\(\alpha(t/x),(t=t')\in H\), then \(\alpha(t'/x)\in H\); (H3) if
\(\bigwedge\Phi\in H\) then \(\Phi\subseteq H\); (H4) if
\(\forall x\,\varphi\in H\) then \(\varphi(t/x)\in H\) for every
\(t\in\mathrm{CT}\); (H5) if \(\bigvee\Phi\in H\) then some member of
\(\Phi\) is in \(H\); (H6) if \(\exists x\,\varphi\in H\) then
\(\varphi(u/x)\in H\) for some \(u\in\mathrm{CT}\).

\emph{A term Hintikka set has a model in which every element is the value
of a closed term.}  Put \(t\sim t'\) iff \((t=t')\in H\).  Reflexivity is
(H1).  If \((t=t')\in H\), (H2) for \(\alpha=(x=t)\), with
\(\alpha(t/x)=(t=t)\in H\), gives \((t'=t)\in H\).  If \((t=t'),(t'=t'')\in
H\), then \((t'=t)\in H\), and (H2) for \(\alpha=(x=t'')\), with
\(\alpha(t'/x)\in H\) and \((t'=t)\in H\), gives \((t=t'')\in H\).
\emph{Replacement}: let \(\alpha\) be atomic with free variables among the
distinct \(x_1,\ldots,x_n\), and \(t_i\sim t'_i\); if
\(\alpha(t_1,\ldots,t_n/x_1,\ldots,x_n)\in H\), then
\(\alpha(t'_1,\ldots,t'_n/x_1,\ldots,x_n)\in H\).  For \(j=1,\ldots,n\) in
turn let \(\beta_j\) be \(\alpha\) with \(x_i\) replaced by \(t'_i\) for
\(i<j\) and by \(t_i\) for \(i>j\); it is atomic with at most \(x_j\) free,
\(\beta_j(t_j/x_j)\) is the sentence reached after \(j-1\) steps, and (H2)
gives \(\beta_j(t'_j/x_j)\in H\).  Let \(M=\mathrm{CT}/{\sim}\), put
\(c^{\mathfrak M}=[c]\) for each constant symbol \(c\) of
\(L\cup\mathrm{Par}\),
\(f^{\mathfrak M}([t_1],\ldots,[t_n])=[f(t_1,\ldots,t_n)]\), and
\(R^{\mathfrak M}=\{([t_1],\ldots,[t_n]):R(t_1,\ldots,t_n)\in H\}\).  The
function is well defined: with \(\alpha\) the formula
\(f(t_1,\ldots,t_n)=f(x_1,\ldots,x_n)\), \(\alpha(\bar t/\bar x)\in H\) by
(H1), and replacement gives \((f(\bar t)=f(\bar t'))\in H\) whenever
\(t_i\sim t'_i\).  By replacement (and symmetry of \(\sim\)),
\(R(\bar t)\in H\) iff \(R(\bar t')\in H\) when \(t_i\sim t'_i\), so
\(([t_1],\ldots,[t_n])\in R^{\mathfrak M}\) iff \(R(t_1,\ldots,t_n)\in H\).
By induction on \(t\in\mathrm{CT}\), \(t^{\mathfrak M}=[t]\); so every
element is the value of a closed term, \(\mathfrak M\models R(\bar t)\) iff
\(R(\bar t)\in H\), and \(\mathfrak M\models t=t'\) iff \((t=t')\in H\).
Now, by induction on the n.n.f. formula \(\varphi\): for every \(\sigma\)
from its free variables to \(\mathrm{CT}\), if \(\varphi\sigma\in H\) then
\(\mathfrak M\models\varphi\sigma\).  Atomic \(\varphi\): shown.  Negated
atomic: (H0) and the atomic case.  \(\bigwedge\Phi\): by (H3) every
\(\chi\sigma_\chi\in H\) (\(\sigma_\chi\) the restriction of \(\sigma\)),
each holds, so the conjunction holds.  \(\bigvee\Phi\): (H5) and the
induction hypothesis.  \(\forall x\,\chi\): with \(\sigma'\) the
restriction of \(\sigma\) to the free variables of \(\chi\) other than
\(x\), \((\forall x\,\chi)\sigma=\forall x\,(\chi\sigma')\), and by (H4)
\((\chi\sigma')(t/x)\in H\) for every \(t\in\mathrm{CT}\); this is
\(\chi(\sigma'\cup\{(x,t)\})\), which holds by the induction hypothesis,
so by (S\('\)) \(\chi\sigma'\) holds at every assignment sending \(x\) to
\(t^{\mathfrak M}\); since every element is some \(t^{\mathfrak M}\),
\(\mathfrak M\models\forall x\,(\chi\sigma')\).  \(\exists x\,\chi\): (H6)
gives \(u\) with \((\chi\sigma')(u/x)\in H\), which holds, and
\(u^{\mathfrak M}\) is a witness by (S\('\)).  Members of \(H\) are the case
of empty \(\sigma\).

\emph{A term Hintikka set containing \(s_0\in\mathcal S\).}  Let \(A_0\) be
\(s_0\) together with all atomic sentences of \(L\cup\mathrm{Par}\)
(countable).  Given \(A_m\), let \(A_{m+1}\) add the members of \(\Phi\)
for each \(\bigwedge\Phi\) or \(\bigvee\Phi\) in \(A_m\), and the sentences
\(\varphi(t/x)\), \(t\in\mathrm{CT}\), for each \(\forall x\,\varphi\) or
\(\exists x\,\varphi\) in \(A_m\).  Each \(A_m\) is countable and so is
\(A=\bigcup_m A_m\) (axiom of choice); \(A\) consists of n.n.f. sentences,
contains \(s_0\) and all atomic sentences, contains \(\Phi\) when it contains
\(\bigwedge\Phi\) or \(\bigvee\Phi\), and contains every \(\varphi(t/x)\)
when it contains \(\forall x\,\varphi\) or \(\exists x\,\varphi\).  Fix
enumerations \(a_0,a_1,\ldots\) of \(A\) and \(t_0,t_1,\ldots\) of
\(\mathrm{CT}\), a variable \(x_0\), an enumeration
\(\alpha_0,\alpha_1,\ldots\) of the atomic formulas of
\(L\cup\mathrm{Par}\) with no variable other than \(x_0\), and a sequence
\((\varphi_n,\psi_n,k_n)_{n<\omega}\) in \(A\times A\times\omega\) in which
every triple occurs for infinitely many \(n\) (compose a surjection of
\(\omega\) onto \(\omega^{4}\) with \((i,j,k,l)\mapsto(a_i,a_j,k)\)).
Given \(s_n\in\mathcal S\), a finite subset of \(A\), let \(u=s_n\) unless
\(\varphi_n\in s_n\), in which case: if \(\varphi_n=\bigwedge\Phi\) and
\(a_{k_n}\in\Phi\), \(u=s_n\cup\{a_{k_n}\}\) (T3); if
\(\varphi_n=\forall x\,\varphi\), \(u=s_n\cup\{\varphi(t_{k_n}/x)\}\) (T4);
if \(\varphi_n=\bigvee\Phi\), \(u=s_n\cup\{a_j\}\) for the least \(j\) with
\(a_j\in\Phi\) and this set in \(\mathcal S\), which exists by (T5) as
\(\Phi\subseteq A\); if \(\varphi_n=\exists x\,\varphi\),
\(u=s_n\cup\{\varphi(t_j/x)\}\) for the least \(j\) with this set in
\(\mathcal S\), which exists by (T6); if \(\varphi_n\) is \((t=t')\),
\(\psi_n\in s_n\) and \(\psi_n=\alpha_{k_n}(t/x_0)\), then
\(u=s_n\cup\{\alpha_{k_n}(t'/x_0)\}\) (T2); otherwise \(u=s_n\).  Put
\(s_{n+1}=u\cup\{t_{k_n}=t_{k_n}\}\), in \(\mathcal S\) by (T1).  Each
\(s_{n+1}\) is a finite subset of \(A\) containing \(s_n\).  Let
\(H=\bigcup_n s_n\).  (H0): two members of \(H\) lie in a common \(s_n\),
and (T0) applies.  (H1): \(t=t_k\) is added at every \(n\) with \(k_n=k\).
(H2): if \(\alpha(t/x),(t=t')\in s_m\), with \(\alpha\) atomic and \(x\) its
only variable (if any), replacing \(x\) by \(x_0\) gives some \(\alpha_k\)
with \(\alpha_k(e/x_0)=\alpha(e/x)\) for every \(e\in\mathrm{CT}\); at an
\(n\geq m\) with \((\varphi_n,\psi_n,k_n)=((t=t'),\alpha(t/x),k)\),
\(\alpha(t'/x)\) is added.  (H3): if \(\bigwedge\Phi\in s_m\) and
\(\varphi\in\Phi\), then \(\varphi=a_k\) for some \(k\), and it is added at
an \(n\geq m\) with \(\varphi_n=\bigwedge\Phi\), \(k_n=k\).  (H4): for
\(t=t_k\), take \(n\geq m\) with \(\varphi_n=\forall x\,\varphi\),
\(k_n=k\).  (H5), (H6): take \(n\geq m\) with \(\varphi_n\) the given
sentence.  So \(H\) is a term Hintikka set, and its model satisfies
\(s_0\subseteq H\).
\end{proof}
\end{annexitem}

\begin{annexitem}{annex:p02-w02c-p02-003-hinf-complete}{thm:p02-w02c-p02-003-infinitary-completeness}{Completeness}
\begin{proof}
\emph{Soundness.}  Let \(T\vdash_{\infty}\varphi\) and let \(\mathfrak M\)
be a structure for \(L\) satisfying \(T\).  Expand \(\mathfrak M\) to
\(L\cup\mathrm{Par}\) by interpreting every parameter as one element of
its nonempty domain.  By fact (a) of \cref{annex:p02-w02c-p02-003-beth},
the expansion satisfies \(T\); by part~1 of
\cref{thm:p02-w02c-p02-003-hinf-soundness} it satisfies \(\varphi\), and by
fact (a) so does \(\mathfrak M\).

\emph{Completeness.}  Suppose \(T\nvdash_{\infty}\varphi\).

\emph{Reduction to a countable language.}  Let \(L_0\) be the set of
symbols of \(L\) occurring in \(\varphi\) or in a member of \(T\).  The
symbols occurring in one formula form a countable set (the countability
step of \cref{annex:p02-w02c-p02-003-interpolation}), and \(T\) is
countable, so \(L_0\) is countable (axiom of choice).  Write
\(\vdash_{\infty}^{L_0}\) for derivability in \(\mathbf H_{\infty}\) built
over \(L_0\) with the same \(\mathrm{Par}\).  The set of theorems of \(T\)
over \(L\) that are sentences for \(L_0\cup\mathrm{Par}\) contains the
axioms over \(L_0\) (a truth assignment to the prime sentences for
\(L_0\cup\mathrm{Par}\) extends to one for \(L\cup\mathrm{Par}\), and the
value of a sentence depends only on the prime sentences built into it; the
closed terms for \(L_0\cup\mathrm{Par}\) are closed terms for
\(L\cup\mathrm{Par}\)), contains \(T\), and is closed under the rules over
\(L_0\).  So \(T\vdash_{\infty}^{L_0}\theta\) implies
\(T\vdash_{\infty}\theta\), and \(T\nvdash_{\infty}^{L_0}\varphi\).  Below,
\(\vdash_{\infty}\) means \(\vdash_{\infty}^{L_0}\), and all sentences are
for \(L_0\cup\mathrm{Par}\).

\emph{The consistency property.}  For a finite set \(s\) of sentences,
\(\bigwedge s\) is a sentence.  Let \(\mathcal S\) be the set of finite
sets \(s\) of n.n.f. sentences in whose members only finitely many
parameters occur and such that \(T\nvdash_{\infty}\neg\bigwedge s\).  The
sets added in (T1)--(T6) below again have finitely many parameters: a
closed term contains finitely many, and \(\varphi\in\Phi\) or
\(\varphi(t/x)\) has the parameters of \(\bigwedge\Phi\), \(\bigvee\Phi\),
\(\forall x\,\varphi\) or \(\exists x\,\varphi\) together with those of
\(t\).  Each claim is proved by assuming the new set is not in
\(\mathcal S\) and deriving \(T\vdash_{\infty}\neg\bigwedge s\), against
\(s\in\mathcal S\); ``by part~1'' refers to
\cref{thm:p02-w02c-p02-003-hinf-nnf}.
(T0): if \(\theta,\neg\theta\in s\), then \(\neg\bigwedge s\) is a
tautology, hence an axiom.
(T1): from \(\neg\bigwedge(s\cup\{t=t\})\) and the axiom \(t=t\), by
part~1, \(\neg\bigwedge s\).
(T2): from \(\neg\bigwedge(s\cup\{\alpha(t'/x)\})\) and the axiom
\(t=t'\to\alpha(t/x)\to\alpha(t'/x)\), with \(\alpha(t/x)\), \(t=t'\in s\),
by part~1, \(\neg\bigwedge s\): an assignment making \(\bigwedge s\) true
makes \(\alpha(t/x)\), \(t=t'\), hence \(\alpha(t'/x)\), and so
\(\bigwedge(s\cup\{\alpha(t'/x)\})\) true.
(T3): if \(\bigwedge\Phi\in s\) and \(\varphi\in\Phi\),
\(\neg\bigwedge(s\cup\{\varphi\})\to\neg\bigwedge s\) is a tautology.
(T4): from \(\neg\bigwedge(s\cup\{\varphi(t/x)\})\) and the axiom
\(\forall x\,\varphi\to\varphi(t/x)\), with \(\forall x\,\varphi\in s\), by
part~1.
(T5): let \(\bigvee\Phi\in s\) and suppose
\(T\vdash_{\infty}\neg\bigwedge(s\cup\{\varphi\})\) for every
\(\varphi\in\Phi\).  By part~1, \(T\vdash_{\infty}\varphi\to\neg\bigwedge s\)
for every \(\varphi\in\Phi\); by part~2 of
\cref{thm:p02-w02c-p02-003-hinf-nnf},
\(T\vdash_{\infty}\bigvee\Phi\to\neg\bigwedge s\); as \(\bigvee\Phi\in s\),
\((\bigvee\Phi\to\neg\bigwedge s)\to\neg\bigwedge s\) is a tautology.
(T6): let \(\exists x\,\varphi\in s\) and let \(p\) be a parameter
occurring in no member of \(s\).  If
\(T\vdash_{\infty}\neg\bigwedge(s\cup\{\varphi(p/x)\})\), then by part~1
\(T\vdash_{\infty}\varphi(p/x)\to\neg\bigwedge s\); \(p\) occurs neither in
\(\varphi\) nor in \(\neg\bigwedge s\), so (R2), second form, gives
\(T\vdash_{\infty}\exists x\,\varphi\to\neg\bigwedge s\), and as
\(\exists x\,\varphi\in s\), part~1 gives \(T\vdash_{\infty}\neg\bigwedge s\).
So \(\mathcal S\) is a consistency property with closed terms for
\(L_0\cup\mathrm{Par}\).

\emph{The initial set.}  Let \(s_0=\{(\bigwedge T)^{+},\varphi^{\ominus}\}\);
\(\bigwedge T\) is a sentence because \(T\) is countable, and \(s_0\) has no
parameter.  Suppose \(T\vdash_{\infty}\neg\bigwedge s_0\).  By part~2 of
\cref{thm:p02-w02c-p02-003-hinf-nnf}, \(T\vdash_{\infty}\bigwedge T\); by
part~3 with empty \(\sigma\),
\(T\vdash_{\infty}(\bigwedge T)^{+}\leftrightarrow\bigwedge T\) and
\(T\vdash_{\infty}\varphi^{\ominus}\leftrightarrow\neg\varphi\).  The
sentence \(\bigwedge T\to((\bigwedge T)^{+}\leftrightarrow\bigwedge T)\to
(\varphi^{\ominus}\leftrightarrow\neg\varphi)\to\neg\bigwedge s_0\to\varphi\)
is a tautology: if the four antecedents are true, \((\bigwedge T)^{+}\) is
true, so \(\varphi^{\ominus}\) is false, so \(\neg\varphi\) is false.  So
\(T\vdash_{\infty}\varphi\), against the assumption.  Hence
\(s_0\in\mathcal S\).

\emph{The model.}  By \cref{thm:p02-w02c-p02-003-term-model-existence},
some structure \(\mathfrak M\) for \(L_0\cup\mathrm{Par}\) satisfies
\((\bigwedge T)^{+}\) and \(\varphi^{\ominus}\); by
\cref{thm:p02-w02c-p02-003-nnf} it satisfies every member of \(T\) and not
\(\varphi\).  Its reduct to \(L_0\), expanded to \(L\) by interpreting the
remaining symbols arbitrarily on its nonempty domain, satisfies the same
sentences for \(L_0\) by fact (a) of \cref{annex:p02-w02c-p02-003-beth}; it
is a model of \(T\) in which \(\varphi\) fails.

The two special cases: take \(T=\varnothing\); and note that
\(\neg\bigwedge\varnothing\) holds in no structure, so \(T\) has a model iff
not every model of \(T\) satisfies \(\neg\bigwedge\varnothing\), iff
\(T\nvdash_{\infty}\neg\bigwedge\varnothing\).
\end{proof}
\end{annexitem}

\begin{annexitem}{annex:p02-w02c-p02-003-hinf-uncount}{thm:p02-w02c-p02-003-hinf-uncountable}{An uncountable theory}
\begin{proof}
If \(\mathfrak M\) were a model of \(T\), then by the first sentence every
element would be some \(d_n^{\mathfrak M}\), so \(n\mapsto d_n^{\mathfrak M}\)
would map \(\omega\) onto \(M\), and by the others
\(\alpha\mapsto c_\alpha^{\mathfrak M}\) would be injective on
\(\omega_1\); composing with a choice of preimages would give an injection
of \(\omega_1\) into \(\omega\), which does not exist.  Suppose
\(T\vdash_{\infty}\neg\bigwedge\varnothing\).  By part~2 of
\cref{thm:p02-w02c-p02-003-hinf-soundness}, some countable
\(T_0\subseteq T\) has \(T_0\vdash_{\infty}\neg\bigwedge\varnothing\).  The
set \(A\) of \(\alpha\) with \(c_\alpha\) occurring in \(T_0\) is
countable; fix an injection \(g:A\to\omega\).  Let \(\mathfrak M\) have
domain \(\omega\), \(d_n^{\mathfrak M}=n\), \(c_\alpha^{\mathfrak M}=g(\alpha)\)
for \(\alpha\in A\) and \(c_\alpha^{\mathfrak M}=0\) otherwise, and every
parameter interpreted as \(0\).  Every element is some \(d_n^{\mathfrak M}\),
and for \(\neg\,c_\alpha=c_\beta\) in \(T_0\), \(\alpha\neq\beta\) both lie
in \(A\) and \(g(\alpha)\neq g(\beta)\).  So \(\mathfrak M\) satisfies every
member of \(T_0\), and by part~1 of
\cref{thm:p02-w02c-p02-003-hinf-soundness} it satisfies
\(\neg\bigwedge\varnothing\); but \(\bigwedge\varnothing\) holds in every
structure.  Hence \(T\nvdash_{\infty}\neg\bigwedge\varnothing\).
\end{proof}
\end{annexitem}

\begin{annexitem}{annex:p02-w02c-p02-003-wso-n}{thm:p02-w02c-p02-003-wso-n}{Characterizing \(\mathcal N\) in \(L^{*}\)}
\begin{proof}
An \(L^{*}\)-formula without set variables is evaluated in
\(\mathrm{HF}(\mathfrak M)\) by the clauses of first-order satisfaction in
\(\mathfrak M\), since its quantifiers range over \(M\); so the axioms of
\(\mathsf Q\) hold in \(\mathrm{HF}(\mathfrak M)\) iff they hold in
\(\mathfrak M\).  For every natural number \(i\) and \(v\in M\) exactly one
\(w\in M\) has \(w=S^{\mathfrak M}v\), so by
\cref{thm:p02-w02c-p02-003-hf-recursion}, with start value
\(0^{\mathfrak M}\), \(\exists z\,(z=0\land\mathrm{Rec}_{\gamma_S}(n,z,x))\)
holds iff \(n\) is a natural number and \(x=h(n)\), where
\(h(0)=0^{\mathfrak M}\) and \(h(k+1)=S^{\mathfrak M}h(k)\).

\emph{If \(\mathrm{HF}(\mathfrak M)\models\alpha_{\mathcal N}\), then
\(\mathfrak M\cong\mathcal N\).}  The axioms of \(\mathsf Q\) hold in
\(\mathfrak M\), and every \(x\in M\) is \(h(n)\) for some \(n\), so
\(h\colon\N\to M\) is onto.  It is one-one: if \(h(k)=h(j)\) with \(k<j\),
then \(k\) applications of the injectivity axiom \(Sx=Sy\to x=y\) give
\(h(0)=h(j-k)\), that is \(0^{\mathfrak M}=S^{\mathfrak M}h(j-k-1)\),
against \(Sx\neq0\).  By definition \(h\) preserves \(0\) and \(S\).  For
\(+\), \(h(k+j)=h(k)+^{\mathfrak M}h(j)\) by induction on \(j\): for
\(j=0\) by \(x+0=x\); and \(h(k)+^{\mathfrak M}h(j+1)
=h(k)+^{\mathfrak M}S^{\mathfrak M}h(j)
=S^{\mathfrak M}(h(k)+^{\mathfrak M}h(j))=S^{\mathfrak M}h(k+j)=h(k+j+1)\)
by \(x+Sy=S(x+y)\).  For \(\times\), \(h(k\cdot j)=h(k)\times^{\mathfrak
M}h(j)\) by induction on \(j\): for \(j=0\) by \(x\times0=0\); and
\(h(k)\times^{\mathfrak M}S^{\mathfrak M}h(j)
=(h(k)\times^{\mathfrak M}h(j))+^{\mathfrak M}h(k)
=h(k\cdot j)+^{\mathfrak M}h(k)=h(k\cdot j+k)\) by
\(x\times Sy=(x\times y)+x\) and the case of \(+\).  So \(h\) is an
isomorphism of \(\mathcal N\) onto \(\mathfrak M\).

\emph{If \(g\colon\mathcal N\cong\mathfrak M\), then
\(\mathrm{HF}(\mathfrak M)\models\alpha_{\mathcal N}\).}  The axioms of
\(\mathsf Q\) are true in \(\mathcal N\), since they state the facts about
successor, addition and multiplication listed in the proof of
\cref{thm:p02-w02c-p02-003-integers-second}, so they are true in
\(\mathfrak M\) by \cref{thm:p02-w02-p02-005-bw-iso-equiv}.  By induction on
\(k\), \(g(k)=h(k)\): \(g(0)=0^{\mathfrak M}\) and
\(g(k+1)=S^{\mathfrak M}g(k)\).  As \(g\) is onto, every \(x\in M\) is
\(h(k)\) for some \(k\), so the last conjunct holds.
\end{proof}
\end{annexitem}

\begin{annexitem}{annex:p02-w02c-p02-003-hf-arith}{thm:p02-w02c-p02-003-hf-arith}{Coding \(\mathrm{HF}(\mathcal N)\) in \(\mathcal N\)}
\begin{proof}
\emph{Graph formulas.}  For a numeral, \(\overline n^{\mathcal N}=n\) by
induction on \(n\).  Let \(f\) be primitive recursive with description
\(d\), and let \(r_1,\ldots,r_k,r\) be terms with values
\(b_1,\ldots,b_k,b\) under an assignment \(s'\) in \(\mathcal N\).  By
\cref{thm:p02-w02b-p02-034-pr-graph},
\(\mathcal N\models\Gamma_d(\overline{b_1},\ldots,\overline{b_k},\overline b)\)
iff \(f(b_1,\ldots,b_k)=b\).  The renaming to new variables and the
substitutions one at a time that produce \([f(r_1,\ldots,r_k)=r]\) are
evaluated by \cref{thm:fol-substitution}, each step replacing the value of
one variable by the value of the substituted term, which is unaffected by
the earlier steps because the renamed variables occur in none of the terms;
with \cref{lem:fol-coincidence} this gives
\[
 \mathcal N\models[f(r_1,\ldots,r_k)=r][s']
 \iff f(b_1,\ldots,b_k)=b. \tag{G}
\]
Hence \(\mathcal N\models\psi(f(r_1,\ldots,r_k))[s']\) iff
\(\mathcal N\models\psi(y)[s'(f(b_1,\ldots,b_k)/y)]\), since exactly one
value of \(y\) satisfies the graph formula.

\emph{Binary expansions.}  (B1) \(j<2^{j}\) for every \(j\), by induction.
(B2) Every \(w\in\N\) equals \(\sum_{j\in B}2^{j}\) for exactly one finite
\(B\subseteq\N\), and then \(j\in B\) iff
\(\operatorname{pty}(\lfloor w/2^{j}\rfloor)=1\).  For the second clause let
\(w=\sum_{i\in B}2^{i}\) and fix \(j\).  Distinct powers of \(2\) below
\(2^{j}\) sum to at most \(\sum_{i<j}2^{i}=2^{j}-1\), so
\(w=2^{j}q+r\) with \(q=\sum_{i\in B,\,i\geq j}2^{i-j}\) and
\(r=\sum_{i\in B,\,i<j}2^{i}<2^{j}\); thus \(\lfloor w/2^{j}\rfloor=q\),
and \(q\) is odd iff its term \(2^{0}\), present iff \(j\in B\), occurs.
This clause determines \(B\) from \(w\), which is the uniqueness.  For
existence use strong induction on \(w\): \(w=0\) is the empty sum; for
\(w>0\), \(w=2\lfloor w/2\rfloor+\operatorname{pty}(w)\) with
\(\lfloor w/2\rfloor<w\), and if \(\lfloor w/2\rfloor=\sum_{i\in B'}2^{i}\)
then \(w=\sum_{j\in B}2^{j}\) for
\(B=\{i+1:i\in B'\}\cup\{0:\operatorname{pty}(w)=1\}\).

\emph{Part~1.}  First \(\mathrm{HF}_n(\N)\subseteq\mathrm{HF}_{n+1}(\N)\)
by induction: \(\mathrm{HF}_0(\N)=\emptyset\), and if
\(\mathrm{HF}_n\subseteq\mathrm{HF}_{n+1}\) then every finite subset of
\(\N\cup\mathrm{HF}_n\) is a finite subset of \(\N\cup\mathrm{HF}_{n+1}\).
Next \(\operatorname{hc}_{n+1}\) extends \(\operatorname{hc}_n\), by
induction on \(n\): both give \(2m\) at an atom; \(\mathrm{HF}_0\) is
empty; and if \(\operatorname{hc}_{n+1}\) extends \(\operatorname{hc}_n\)
and \(a\in\mathrm{HF}_{n+1}\), the members of \(a\) lie in
\(\N\cup\mathrm{HF}_n\), so
\(\operatorname{hc}_{n+2}(a)=2\sum_{u\in a}2^{\operatorname{hc}_{n+1}(u)}+1
=2\sum_{u\in a}2^{\operatorname{hc}_{n}(u)}+1=\operatorname{hc}_{n+1}(a)\).
So \(\operatorname{hc}\) is a function, \(\operatorname{hc}(m)=2m\), and
\(\operatorname{hc}(a)=2\sum_{u\in a}2^{\operatorname{hc}(u)}+1\) for every
set \(a\); atoms receive even and sets odd values.

We show by strong induction on \(v\) that exactly one
\(z\in\N\cup\mathrm{HF}(\N)\) has \(\operatorname{hc}(z)=v\).  If \(v=2m\),
no set has value \(v\), and the atom \(m\) is the only atom with value
\(v\).  Let \(v=2w+1\), and let \(B\) be as in (B2) for \(w\).  Each
\(j\in B\) has \(j<2^{j}\leq w<v\) by (B1), so by the induction hypothesis
exactly one \(z_j\) has \(\operatorname{hc}(z_j)=j\).  Each \(z_j\) lies in
\(\N\cup\mathrm{HF}_{n_j}(\N)\) for some \(n_j\); with \(n\) the largest
\(n_j\) (or \(0\) if \(B=\emptyset\)), the finite set
\(a=\{z_j:j\in B\}\) is a subset of \(\N\cup\mathrm{HF}_n(\N)\), so
\(a\in\mathrm{HF}_{n+1}(\N)\).  The \(z_j\) are distinct, so
\(\operatorname{hc}(a)=2\sum_{j\in B}2^{j}+1=v\).  If also
\(\operatorname{hc}(b)=v\), then \(b\) is a set and
\(\sum_{u\in b}2^{\operatorname{hc}(u)}=w\).  For \(u\in b\),
\(\operatorname{hc}(u)<2^{\operatorname{hc}(u)}\leq w<v\), so by the
induction hypothesis distinct members of \(b\) have distinct values; hence
\(\{\operatorname{hc}(u):u\in b\}\) is a finite set whose powers of \(2\)
sum to \(w\), so it is \(B\) by (B2), and each \(u\in b\) is
\(z_{\operatorname{hc}(u)}\) by the induction hypothesis.  Thus \(b=a\).
This proves that \(\operatorname{hc}\) is a bijection onto \(\N\); the even
numbers are the values at atoms and the odd numbers the values at sets.

\emph{Part~2.}  Addition and multiplication are primitive recursive by
\(x+0=x\), \(x+(n+1)=S(x+n)\), \(x\cdot0=0\), \(x\cdot(n+1)=x\cdot n+x\);
\(2^{u}\) by \(2^{0}=1\), \(2^{u+1}=2\cdot2^{u}\).  For \(y\geq1\),
\(\lfloor x/y\rfloor\) is the least \(q\leq x\) with \(x<(q+1)y\), a least
search bounded by \(x\), which is primitive recursive with truncated
subtraction and the order test by the closure facts recorded in
\cref{sec:arithmetic-coding-04}; the divisors used below, \(2\) and
\(2^{u}\), are at least \(1\).  So \(\operatorname{pty}\) and
\(\chi_{\in}\) are primitive recursive.  Let \(a\) be a set and
\(v=\operatorname{hc}(a)=2w+1\) with
\(w=\sum_{u\in a}2^{\operatorname{hc}(u)}\).  Then
\(\operatorname{pty}(v)=1\) and \(\lfloor v/2\rfloor=w\).  The values
\(\operatorname{hc}(u)\), \(u\in a\), are distinct by part~1, so by (B2)
\(\chi_{\in}(\operatorname{hc}(z),v)=1\) iff
\(\operatorname{hc}(z)=\operatorname{hc}(u)\) for some \(u\in a\), iff
\(z\in a\) by part~1.

\emph{Part~3.}  By induction on \(\theta\), for all \(s\) at once.  For an
\(\Lang_A\)-term \(t\), the value of \(t^{\iota}\) under
\(s^{\mathrm{ar}}\) is the value of \(t\) under \(s\), by induction on \(t\),
since \(s^{\mathrm{ar}}(x_{2i})=s(x_i)\).  This gives the atomic case
\(r=t\).  For \(t\in a_i\): with \(m\) the value of \(t\),
\(\mathrm{HF}(\mathcal N)\models(t\in a_i)[s]\) iff \(m\in s(a_i)\), iff
\(\chi_{\in}(2m,\operatorname{hc}(s(a_i)))=1\) by part~2, as
\(\operatorname{hc}(m)=2m\); the term \((SS0)\times t^{\iota}\) has value
\(2m\) and \(x_{2i+1}\) has value \(\operatorname{hc}(s(a_i))\) under
\(s^{\mathrm{ar}}\), so by (G) this holds iff
\(\mathcal N\models(t\in a_i)^{\mathrm{ar}}[s^{\mathrm{ar}}]\).  The case
\(a_i\in a_j\) is the same with \(z=s(a_i)\).  The case \(a_i=a_j\) holds
because \(\operatorname{hc}\) is one-one.  Connectives commute with the
translation and with satisfaction.  For \(\exists x_i\,\theta\): the
renaming \(\iota\) is one-one, so
\((s(m/x_i))^{\mathrm{ar}}=s^{\mathrm{ar}}(m/x_{2i})\); hence
\(\mathrm{HF}(\mathcal N)\models\exists x_i\,\theta[s]\) iff some \(m\in\N\)
has \(\mathcal N\models\theta^{\mathrm{ar}}[s^{\mathrm{ar}}(m/x_{2i})]\),
iff \(\mathcal N\models\exists x_{2i}\,\theta^{\mathrm{ar}}[s^{\mathrm{ar}}]\).
For \(\exists a_i\,\theta\): \((s(b/a_i))^{\mathrm{ar}}
=s^{\mathrm{ar}}(\operatorname{hc}(b)/x_{2i+1})\), and as \(b\) ranges over
\(\mathrm{HF}(\N)\), \(\operatorname{hc}(b)\) ranges over exactly the odd
numbers (part~1), which are the values satisfying \(\mathrm{Odd}\) in
\(\mathcal N\); so the left side holds iff some odd \(v\) has
\(\mathcal N\models\theta^{\mathrm{ar}}[s^{\mathrm{ar}}(v/x_{2i+1})]\), iff
\(\mathcal N\models(\exists a_i\,\theta)^{\mathrm{ar}}[s^{\mathrm{ar}}]\).
The universal cases are the same with ``every'' in place of ``some''.

\emph{Part~4.}  The free variables of \(\theta^{\mathrm{ar}}\) are the images
under \(\iota\) of those of \(\theta\): the formulas \(\mathrm{Odd}\) and
\([\chi_{\in}(\cdot,\cdot)=S0]\) add only bound variables, by the renaming
step in the definition of \([f(\bar r)=r]\) and
\cref{lem:substitution-fv}.  Given \(m_1,\ldots,m_k\), choose any \(s\) with
\(s(x_{i_j})=m_j\); by \cref{lem:fol-coincidence} and part~3,
\(\mathrm{HF}(\mathcal N)\models\theta[s]\) iff
\(\mathcal N\models\theta^{\mathrm{ar}}[s^{\mathrm{ar}}]\), and
\(s^{\mathrm{ar}}(x_{2i_j})=m_j\).
\end{proof}
\end{annexitem}

\begin{annexitem}{annex:p02-w02c-p02-003-truth-clauses}{thm:p02-w02c-p02-003-truth-clauses}{Truth clauses over \(\mathcal N\)}
\begin{proof}
\emph{Part~1.}  \(h\), \(t\), the list entries \((c)_j\), recognition of
formula codes, the list of free variable indices, and the numeral
substitution function are primitive recursive, and
\(\operatorname{sb}(\ulcorner\varphi\urcorner,i,n)=\ulcorner\varphi[\overline n/x_i]\urcorner\),
by \cref{sec:arithmetic-coding-04}; a sentence is a formula whose list of
free variables is empty, so \(\operatorname{Snt}\) is primitive recursive,
and so are \(\operatorname{arg}_0,\operatorname{arg}_1\).  For \(c>0\),
\(h(c),t(c)<c\) (same section).  If \(t(c)>0\) then
\(\operatorname{arg}_0(c)\) is \(h(t(c))<t(c)<c\) or \(0\); if
\(t(c)=0\) it is \(0\).  Likewise \(\operatorname{arg}_1(c)\) is
\(h(t(t(c)))<c\) or \(0\).  So for \(c>0\) both arguments are less than
\(c\), and \(\operatorname{tv}(0)=0\) as \(h(0)=0\); the definition of
\(\operatorname{tv}\) is a course-of-values recursion with a primitive
recursive update using only smaller values, which is primitive recursive by
the construction recorded in \cref{sec:arithmetic-coding-04}.  For a code
\(c=N_k(c_0,c_1)=C(k,[c_0,c_1])\), \(h(c)=k\), and
\([c_0,c_1]=C(c_0,C(c_1,0))\) gives \(\operatorname{arg}_0(c)=c_0\) and
\(\operatorname{arg}_1(c)=c_1\); for \(N_k(c_0)\),
\(\operatorname{arg}_0=c_0\).  By induction on a closed term \(r\):
\(\ulcorner0\urcorner=N_1()\) has \(\operatorname{tv}=0\);
\(\operatorname{tv}(\ulcorner S(r)\urcorner)=\operatorname{tv}(\ulcorner r\urcorner)+1=r^{\mathcal N}+1\);
and the cases \(+\), \(\times\) are the same.

\emph{Part~2.}  Write \(c(\varphi)\) for the formula complexity of
\cref{sec:fol-syntax-semantics-04}.  By the coding, a sentence code
\(c=\ulcorner\sigma\urcorner\) has \(h(c)\) equal to the tag of the outer
constructor of \(\sigma\), and its arguments are the codes of the
immediate parts: for \(\sigma\) of the form \(r=t\), \(r\) and \(t\) are
closed terms; for \(\neg\varphi\) or \(\varphi\circ\psi\), \(\varphi\) and
\(\psi\) are sentences of smaller complexity; for \(Qx_i\,\varphi\),
\(\operatorname{arg}_0(c)=i\), \(\operatorname{arg}_1(c)=\ulcorner\varphi\urcorner\)
and every free variable of \(\varphi\) is \(x_i\).  In the last case, for
each \(n\), \(\varphi[\overline n/x_i]\) is a sentence by
\cref{lem:substitution-fv}, its complexity is \(c(\varphi)<c(Qx_i\,\varphi)\)
by \cref{lem:substitution-complexity}, its code is
\(\operatorname{sb}(\operatorname{arg}_1(c),\operatorname{arg}_0(c),n)\), and
by \cref{thm:fol-substitution} and \(\overline n^{\mathcal N}=n\),
\begin{equation*}
\begin{gathered}
 \mathcal N\models\forall x_i\,\varphi\iff
 \mathcal N\models\varphi[\overline n/x_i]\text{ for every }n,\\
 \mathcal N\models\exists x_i\,\varphi\iff
 \mathcal N\models\varphi[\overline n/x_i]\text{ for some }n.
\end{gathered}
\tag{S}
\end{equation*}
By (G) of \cref{annex:p02-w02c-p02-003-hf-arith}, \((\mathcal N,X)\) satisfies
the conjunct \(\forall c\,(T(c)\to[\operatorname{Snt}(c)=\overline1])\) iff
\(X\) consists of sentence codes, and satisfies the conjunct for tag \(k\)
iff for every sentence code \(c\) with \(h(c)=k\) the clause
\(\beta_k(c)\) holds with \(T\) read as \(X\) and each function term read as
the value of the function.

If \(X=\mathrm{TA}\), every conjunct holds: \(\mathrm{TA}\) consists of
sentence codes; for \(c=\ulcorner r=t\urcorner\), \(c\in\mathrm{TA}\) iff
\(r^{\mathcal N}=t^{\mathcal N}\) iff
\(\operatorname{tv}(\operatorname{arg}_0(c))=\operatorname{tv}(\operatorname{arg}_1(c))\)
by part~1; \(\ulcorner\bottom\urcorner\notin\mathrm{TA}\); the connective
tags hold by the truth clauses of \(\neg,\land,\lor,\to\) in
\(\mathcal N\); and the tags \(11\), \(12\) hold by (S).

Conversely let \((\mathcal N,X)\models\mathrm{TC}(T)\).  We show by strong
induction on \(c(\sigma)\) that \(\ulcorner\sigma\urcorner\in X\) iff
\(\mathcal N\models\sigma\), for every sentence \(\sigma\).  For \(r=t\) the
tag-\(5\) conjunct and part~1 give \(\ulcorner\sigma\urcorner\in X\) iff
\(r^{\mathcal N}=t^{\mathcal N}\).  For \(\bottom\) the tag-\(6\) conjunct
gives \(\ulcorner\bottom\urcorner\notin X\).  For \(\neg\varphi\), the
tag-\(7\) conjunct gives \(\ulcorner\sigma\urcorner\in X\) iff
\(\ulcorner\varphi\urcorner\notin X\), iff \(\mathcal N\not\models\varphi\)
by the induction hypothesis; the binary connectives are alike.  For
\(\forall x_i\,\varphi\), the tag-\(11\) conjunct gives
\(\ulcorner\sigma\urcorner\in X\) iff
\(\ulcorner\varphi[\overline n/x_i]\urcorner\in X\) for every \(n\), iff
\(\mathcal N\models\varphi[\overline n/x_i]\) for every \(n\) by the
induction hypothesis, iff \(\mathcal N\models\sigma\) by (S); \(\exists\)
is alike.  As \(X\) consists of sentence codes, \(X=\mathrm{TA}\).
\end{proof}
\end{annexitem}

\begin{annexitem}{annex:p02-w02c-p02-003-wso-beth-failure}{thm:p02-w02c-p02-003-wso-beth-failure}{Beth's theorem fails for weak second-order logic}
\begin{proof}
Two facts about evaluation in \(\mathrm{HF}((\mathfrak M,X))\), each by
induction on formulas.  (O) An \(L^{*}\)-formula in which \(T\) does not
occur has the same truth value under each assignment in
\(\mathrm{HF}((\mathfrak M,X))\) and in \(\mathrm{HF}(\mathfrak M)\), since
the two structures differ only in the interpretation of \(T\).  (F) A
formula of \(\Lang_A\cup\{T\}\) without set variables has the same truth
value in \(\mathrm{HF}((\mathfrak M,X))\) and in \((\mathfrak M,X)\), as
its quantifiers range over \(M\) (the first step of
\cref{annex:p02-w02c-p02-003-wso-n}).

\emph{Part~1.}  By (O), (F) and the definition,
\(\mathrm{HF}((\mathfrak M,X))\models\alpha_{\mathrm{TA}}\) iff
\(\mathrm{HF}(\mathfrak M)\models\alpha_{\mathcal N}\) and
\((\mathfrak M,X)\models\mathrm{TC}(T)\).  Suppose this holds for \(X\) and
for \(Y\).  By \cref{thm:p02-w02c-p02-003-wso-n} there is an isomorphism
\(g\colon\mathcal N\to\mathfrak M\).  It is an isomorphism of
\((\mathcal N,g^{-1}[X])\) onto \((\mathfrak M,X)\), so
\((\mathcal N,g^{-1}[X])\models\mathrm{TC}(T)\) by
\cref{thm:p02-w02-p02-005-bw-iso-equiv}, and \(g^{-1}[X]=\mathrm{TA}\) by
\cref{thm:p02-w02c-p02-003-truth-clauses}.  Likewise
\(g^{-1}[Y]=\mathrm{TA}\), so \(X=g[\mathrm{TA}]=Y\).  For the second
assertion, \(\mathrm{HF}(\mathcal N)\models\alpha_{\mathcal N}\) by
\cref{thm:p02-w02c-p02-003-wso-n} (with the identity isomorphism) and
\((\mathcal N,\mathrm{TA})\models\mathrm{TC}(T)\) by
\cref{thm:p02-w02c-p02-003-truth-clauses}.

\emph{Part~2.}  Suppose \(\theta(x)\) were such a formula, with
\(x=x_i\).  By part~1, \(\mathrm{HF}((\mathcal N,\mathrm{TA}))\) is a model
of \(\alpha_{\mathrm{TA}}\), so for every \(m\in\N\), \(m\in\mathrm{TA}\)
iff \(\mathrm{HF}((\mathcal N,\mathrm{TA}))\models\theta[m]\), iff
\(\mathrm{HF}(\mathcal N)\models\theta[m]\) by (O), iff
\(\mathcal N\models\theta^{\mathrm{ar}}[m]\), with \(m\) assigned to
\(x_{2i}\), by part~4 of \cref{thm:p02-w02c-p02-003-hf-arith}.  Let
\(\mathsf{Tr}(x_{2i})\) be \(\theta^{\mathrm{ar}}\).  For every
\(\Lang_A\)-sentence \(\sigma\), by \cref{thm:fol-substitution} and
\(\overline n^{\mathcal N}=n\),
\(\mathcal N\models\mathsf{Tr}(\overline{\ulcorner\sigma\urcorner})\) iff
\(\ulcorner\sigma\urcorner\in\mathrm{TA}\) iff \(\mathcal N\models\sigma\).
So \(\mathcal N\models\mathsf{Tr}(\overline{\ulcorner\sigma\urcorner})
\leftrightarrow\sigma\) for every sentence \(\sigma\): \(\mathsf{Tr}\) is an
arithmetical truth definition, which \cref{thm:tarski-undefinability}
excludes.
\end{proof}
\end{annexitem}

\begin{annexitem}{annex:p02-w02c-p02-003-q0-equivalence}{thm:p02-w02c-p02-003-q0-equivalence}{\(L(Q_0)\), \(\omega\)-logic and weak second-order logic}
\begin{proof}
Throughout, (F) is the fact recorded at the start of
\cref{annex:p02-w02c-p02-003-wso-beth-failure}: a formula without set
variables has the same truth value in \(\mathrm{HF}(\mathfrak A)\) as in
\(\mathfrak A\).  For a unary symbol \(D\), \(\chi^{D}\) is \(\chi\) with
every quantifier \(\exists x\), \(\forall x\), \(Q_0x\) relativized to
\(D\) (\(\exists x\,(D(x)\land\cdots)\), \(\forall x\,(D(x)\to\cdots)\),
\(Q_0x\,(D(x)\land\cdots)\)).  If \(D^{\mathfrak A}\) is nonempty and closed
under the functions and constants of a language \(K\), and \(\chi\) is a
formula of \(K\), then \(\mathfrak A\models\chi^{D}[s]\) iff
\(\mathfrak A^{D}\models\chi[s]\) for assignments \(s\) into
\(D^{\mathfrak A}\), where \(\mathfrak A^{D}\) is the structure for \(K\) on
\(D^{\mathfrak A}\) with the restricted interpretations; this is an induction
on \(\chi\), in which terms take values in \(D^{\mathfrak A}\) by closure.
In each part the new sentence will state that \(U\) is contained in the
part of the structure that carries the given model, so the relativized
reduct \(\mathfrak A^{(U)}_{L}\) is unchanged when that part is cut out.

\emph{Part~1.}  By induction on \(\varphi\); the only case beyond the
clauses of satisfaction is \(Q_0x\,\psi\).  For a set variable \(a\) not in
\(\psi^{w}\), \(\mathrm{HF}(\mathfrak M)\models\exists a\,\forall x\,
(\psi^{w}\to x\in a)[s]\) iff some \(b\in\mathrm{HF}(M)\) contains every
\(m\) with \(\mathrm{HF}(\mathfrak M)\models\psi^{w}[s(m/x)]\), that is, by
the induction hypothesis, every \(m\) with \(\mathfrak M\models\psi[s(m/x)]\).
Every member of \(\mathrm{HF}(M)\) is finite, and every finite subset of
\(M\) lies in \(\mathrm{HF}_1(M)\); so this holds iff finitely many \(m\)
satisfy \(\psi\), and its negation holds iff
\(\mathfrak M\models Q_0x\,\psi[s]\).  If \(K\) is RPC in \(L(Q_0)\) by
\(\varphi\) for \(L'\), then \(\varphi^{w}\) is an \(L'^{*}\)-sentence with
the same models, so \(K\) is RPC in weak second-order logic by
\(\varphi^{w}\).

\emph{Part~2.}  Let \(K\) be RPC in weak second-order logic by \(\varphi\)
for \(L'\).  Take new symbols: unary \(A\), \(\mathrm{St}\), \(\mathrm{Wd}\),
binary \(E\) and \(\prec\), and a unary function symbol \(\mathrm{rk}\), and
let \(\Theta\) be the conjunction of the following sentences of \(L(Q_0)\)
for \(L''=L'\cup\{A,\mathrm{St},\mathrm{Wd},E,\prec,\mathrm{rk}\}\):
\begin{enumerate}
\item[(H1)] every element lies in exactly one of \(A\), \(\mathrm{St}\),
      \(\mathrm{Wd}\); \(\exists x\,A(x)\); \(A(c)\) for each constant
      symbol \(c\in L'\), and
      \(\forall\bar x\,(A(x_1)\land\cdots\land A(x_n)\to A(f(\bar x)))\) for
      each \(n\)-ary function symbol \(f\in L'\); and
      \(\forall x\,(U(x)\to A(x))\);
\item[(H2)] \(\forall z\,\forall b\,(E(z,b)\to((A(z)\lor\mathrm{St}(z))\land
      \mathrm{St}(b)))\);
\item[(H3)] \(\forall a\,\forall b\,((\mathrm{St}(a)\land\mathrm{St}(b)\land
      \forall z\,(E(z,a)\leftrightarrow E(z,b)))\to a=b)\);
\item[(H4)] \(\forall b\,(\mathrm{St}(b)\to\neg Q_0z\,E(z,b))\);
\item[(H5)] \(\exists e\,(\mathrm{St}(e)\land\forall z\,\neg E(z,e))\), and
      \(\forall b\,\forall z\,((\mathrm{St}(b)\land(A(z)\lor\mathrm{St}(z)))
      \to\exists b'\,(\mathrm{St}(b')\land\forall w\,(E(w,b')\leftrightarrow
      (E(w,b)\lor w=z))))\);
\item[(H6)] \(\prec\) is irreflexive and transitive, any two distinct
      elements of \(\mathrm{Wd}\) are comparable, every
      \(x\) with \(\mathrm{Wd}(x)\) has \(\neg Q_0y\,(y\prec x)\),
      \(\forall b\,(\mathrm{St}(b)\to\mathrm{Wd}(\mathrm{rk}(b)))\), and
      \(\forall a\,\forall b\,((\mathrm{St}(a)\land E(a,b))\to
      \mathrm{rk}(a)\prec\mathrm{rk}(b))\).
\end{enumerate}
For an \(L'^{*}\)-formula \(\theta\), \(\theta^{h}\) is the \(L''\)-formula
obtained by renaming \(x_i\) to \(x_{2i}\) and \(a_i\) to \(x_{2i+1}\),
replacing \(t\in a\) and \(a\in b\) by \(E(t,a)\) and \(E(a,b)\), and
relativizing the individual quantifiers to \(A\) and the set quantifiers to
\(\mathrm{St}\).  The required sentence is \(\Theta\land\varphi^{h}\).

\emph{Claim 2a.}  If \(\mathfrak A\models\Theta\) and \(M=A^{\mathfrak A}\),
there is a bijection \(\pi\) of \(M\cup\mathrm{St}^{\mathfrak A}\) onto
\(M\cup\mathrm{HF}(M)\) that is the identity on \(M\), maps
\(\mathrm{St}^{\mathfrak A}\) onto \(\mathrm{HF}(M)\), and satisfies
\(E^{\mathfrak A}(z,b)\) iff \(\pi(z)\in\pi(b)\) for \(z\in M\cup
\mathrm{St}^{\mathfrak A}\) and \(b\in\mathrm{St}^{\mathfrak A}\).  For
\(x\in\mathrm{Wd}^{\mathfrak A}\) let \(\rho(x)\) be the number of its
\(\prec\)-predecessors, finite by (H6).  If \(x\prec y\) in
\(\mathrm{Wd}^{\mathfrak A}\), the predecessors of \(x\) together with \(x\)
are predecessors of \(y\) by transitivity, and \(x\) is not a predecessor of
itself, so \(\rho(x)<\rho(y)\).  Put \(d(b)=\rho(\mathrm{rk}(b))\) for
\(b\in\mathrm{St}^{\mathfrak A}\); by (H6), \(d(a)<d(b)\) whenever
\(a\in\mathrm{St}^{\mathfrak A}\) and \(E^{\mathfrak A}(a,b)\).  Define
\(\pi(m)=m\) for \(m\in M\) and, by strong induction on \(d(b)\),
\(\pi(b)=\{z\in M:E(z,b)\}\cup\{\pi(a):a\in\mathrm{St}^{\mathfrak A},
E(a,b)\}\), a finite set by (H4); by induction on \(d(b)\), \(\pi(b)\) is a
finite subset of \(M\cup\mathrm{HF}_n(M)\) for some \(n\), so it lies in
\(\mathrm{HF}(M)\).  \(\pi\) is one-one on \(\mathrm{St}^{\mathfrak A}\): if
\(\pi(a)=\pi(b)\), by strong induction on \(d(a)+d(b)\), let \(E(z,a)\).
If \(z\in M\), then \(z\in\pi(b)\), and as the elements of \(M\) are atoms,
not members of \(\mathrm{HF}(M)\), \(z\) is one of the members of \(\pi(b)\)
taken from \(M\), so \(E(z,b)\).  If \(z\in\mathrm{St}^{\mathfrak A}\), then
\(\pi(z)\in\mathrm{HF}(M)\) lies in \(\pi(b)\), so \(\pi(z)=\pi(z')\) for
some \(z'\in\mathrm{St}^{\mathfrak A}\) with \(E(z',b)\), and \(z=z'\) by the
induction hypothesis, as \(d(z)+d(z')<d(a)+d(b)\).  By symmetry \(a\) and
\(b\) have the same \(E\)-members, so \(a=b\) by (H3).  \(\pi\) maps
\(\mathrm{St}^{\mathfrak A}\) onto \(\mathrm{HF}(M)\): by induction on the
size of a finite \(Y\subseteq M\cup\pi[\mathrm{St}^{\mathfrak A}]\), some
\(b\) has \(\pi(b)=Y\); for \(Y=\emptyset\) take \(e\) of (H5); for
\(Y=Y'\cup\{y\}\), take \(b'\) with \(\pi(b')=Y'\) and \(z\in M\cup
\mathrm{St}^{\mathfrak A}\) with \(\pi(z)=y\), and the \(b\) given by (H5)
for \(b'\) and \(z\) has \(\pi(b)=Y'\cup\{\pi(z)\}=Y\).  Then by induction
on \(n\), \(\mathrm{HF}_n(M)\subseteq\pi[\mathrm{St}^{\mathfrak A}]\), since
each member of \(\mathrm{HF}_{n+1}(M)\) is a finite subset of
\(M\cup\mathrm{HF}_n(M)\).  Finally, \(E(z,b)\) implies \(\pi(z)\in\pi(b)\)
by definition; conversely, if \(\pi(z)\in\pi(b)\) with \(z\in M\) then
\(E(z,b)\) as above, and with \(z\in\mathrm{St}^{\mathfrak A}\) then
\(\pi(z)=\pi(z')\) for some \(z'\) with \(E(z',b)\), and \(z=z'\).

\emph{Claim 2b.}  If \(\mathfrak A\models\Theta\), then for every
\(L'^{*}\)-formula \(\theta\) and assignment \(s\) in
\(\mathrm{HF}(\mathfrak A^{A})\), \(\mathrm{HF}(\mathfrak A^{A})\models
\theta[s]\) iff \(\mathfrak A\models\theta^{h}[s^{h}]\), where
\(s^{h}(x_{2i})=s(x_i)\) and \(s^{h}(x_{2i+1})=\pi^{-1}(s(a_i))\).  By
(H1), \(A^{\mathfrak A}\) is closed, so \(\mathfrak A^{A}\) is a structure
for \(L'\) and \(L'\)-terms take the same values in both.  The atomic
\(L'\)-case is immediate; \(t\in a_i\) and \(a_i\in a_j\) hold by Claim~2a,
as \(\pi\) is the identity on \(M\); \(a_i=a_j\) holds as \(\pi\) is one-one;
individual quantifiers range over \(M=A^{\mathfrak A}\) on both sides; and
set quantifiers range over \(\mathrm{HF}(M)\) on the left and over
\(\mathrm{St}^{\mathfrak A}=\pi^{-1}[\mathrm{HF}(M)]\) on the right.

\emph{Claim 2c.}  Every structure \(\mathfrak B\) for \(L'\) is
\(\mathfrak A^{A}\) for some model \(\mathfrak A\) of \(\Theta\) with
\(U^{\mathfrak A}=U^{\mathfrak B}\).  Let \(W=\{w_n:n\in\N\}\) be a set of
distinct objects disjoint from \(B\cup\mathrm{HF}(B)\), and let
\(\mathfrak A\) have domain \(B\cup\mathrm{HF}(B)\cup W\), \(A=B\),
\(\mathrm{St}=\mathrm{HF}(B)\), \(\mathrm{Wd}=W\), \(E\) the membership
relation between \(B\cup\mathrm{HF}(B)\) and \(\mathrm{HF}(B)\),
\(w_m\prec w_n\) iff \(m<n\), and \(\mathrm{rk}(b)=w_{r(b)}\) for
\(b\in\mathrm{HF}(B)\), with \(r(b)\) the least \(n\) such that
\(b\in\mathrm{HF}_n(B)\), and \(\mathrm{rk}=w_0\) elsewhere; the relations of
\(L'\) are those of \(\mathfrak B\), and each function of \(L'\) is that of
\(\mathfrak B\) on arguments in \(B\) and takes a fixed value \(b_0\in B\)
otherwise.  (H1)--(H5) hold because the members of \(\mathrm{HF}(B)\) are
finite sets of elements of \(B\cup\mathrm{HF}(B)\), \(\emptyset\in
\mathrm{HF}_1(B)\), every finite subset of \(B\cup\mathrm{HF}(B)\) lies in
\(\mathrm{HF}(B)\) (it lies in \(B\cup\mathrm{HF}_n(B)\) for some \(n\), as
\(\mathrm{HF}_n(B)\subseteq\mathrm{HF}_{n+1}(B)\)), and two sets with the
same members are equal.  For (H6), \(w_n\) has \(n\) predecessors, and if
\(a\in b\) with \(a\) a set then \(b\in\mathrm{HF}_{r(b)}(B)\) is a subset of
\(B\cup\mathrm{HF}_{r(b)-1}(B)\), so \(r(a)\leq r(b)-1\).

Now if \(\mathfrak A\models\Theta\land\varphi^{h}\), Claim~2b gives
\(\mathrm{HF}(\mathfrak A^{A})\models\varphi\), so \(\mathfrak A^{A}\) is a
model of \(\varphi\), and \(\mathfrak A^{(U)}_{L}=(\mathfrak A^{A})^{(U)}_{L}\)
as \(U^{\mathfrak A}\subseteq A^{\mathfrak A}\); this also shows that
\(U^{\mathfrak A}\) is nonempty and closed.  Conversely, for a model
\(\mathfrak B\) of \(\varphi\), the \(\mathfrak A\) of Claim~2c satisfies
\(\Theta\), and \(\varphi^{h}\) by Claim~2b, with
\(\mathfrak A^{(U)}_{L}=\mathfrak B^{(U)}_{L}\).  So \(K\) is RPC in
\(L(Q_0)\) by \(\Theta\land\varphi^{h}\).

\emph{Part~3.}  Let \(K\) be RPC in \(L(Q_0)\) by \(\varphi\) for \(L'\).
Choose new symbols: unary \(D\), \(\mathrm{Nn}\), and \(\hat0,\hat S,\hat+,
\hat\times\); write \(u\mathbin{\hat<}v\) for
\(\exists z\,(\mathrm{Nn}(z)\land u\mathbin{\hat+}\hat Sz=v)\).  For each
occurrence of a subformula \(Q_0x\,\psi\) of \(\varphi\) whose free
variables are \(\bar y=(y_1,\ldots,y_k)\), choose a new \(k\)-ary relation
symbol \(R_\psi\), new \((k+1)\)-ary function symbols \(h_\psi\), \(g_\psi\)
and a new \(k\)-ary function symbol \(b_\psi\).  For a subformula \(\chi\)
of \(\varphi\), \(\chi'\) is the first-order formula obtained by replacing,
from the innermost outward, each \(Q_0x\,\psi\) by \(R_\psi(\bar y)\), and
then relativizing every quantifier to \(D\).  The sentence of
\(\omega\)-logic is the conjunction \(\varphi^{\omega}\) of: \(\exists
x\,D(x)\); \(D\) is closed under the constants and functions of \(L'\);
\(\forall x\,(U(x)\to D(x))\); \(\forall x\,(D(x)\leftrightarrow\neg
\mathrm{Nn}(x))\); \(\varphi'\); and for each \(Q_0x\,\psi\), with
\(\forall\bar y\in D\) abbreviating relativized universal quantifiers,
\begin{align*}
 \text{(W1)}\;&\forall\bar y\in D\,\bigl(R_\psi(\bar y)\to\bigl[\forall n\,(\mathrm{Nn}(n)\to
   (D(h_\psi(\bar y,n))\land\psi'(h_\psi(\bar y,n),\bar y)))\\
 &\quad{}\land\forall n\,\forall n'\,((\mathrm{Nn}(n)\land\mathrm{Nn}(n')\land
   h_\psi(\bar y,n)=h_\psi(\bar y,n'))\to n=n')\bigr]\bigr),\\
 \text{(W2)}\;&\forall\bar y\in D\,\bigl(\neg R_\psi(\bar y)\to\bigl[\forall x\,((D(x)\land\psi'(x,\bar y))\\
 &\quad{}\to(\mathrm{Nn}(g_\psi(\bar y,x))\land g_\psi(\bar y,x)\mathbin{\hat<}b_\psi(\bar y)))\\
 &\quad{}\land\forall x\,\forall x'\,((D(x)\land D(x')\land\psi'(x,\bar y)\land\psi'(x',\bar y)\\
 &\qquad{}\land g_\psi(\bar y,x)=g_\psi(\bar y,x'))\to x=x')\bigr]\bigr),
\end{align*}
where \(\psi'(t,\bar y)\) is \(\psi'\) with \(t\) substituted for \(x\).

Let \(\mathfrak A\) be an \(\omega\)-model of \(\varphi^{\omega}\), with
\(j\colon\mathcal N\cong(\mathrm{Nn}^{\mathfrak A},\hat0,\hat S,\hat+,
\hat\times)\).  As \(j\) preserves \(S\) and \(+\),
\(u\mathbin{\hat<}v\) holds of \(u,v\in\mathrm{Nn}^{\mathfrak A}\) iff
\(j^{-1}(u)<j^{-1}(v)\).  We show by induction on subformulas \(\chi\) of
\(\varphi\) that \(\mathfrak A\models\chi'[s]\) iff
\(\mathfrak A^{D}\models\chi[s]\) for assignments \(s\) into
\(D^{\mathfrak A}\).  The atomic, connective and quantifier cases follow from
the relativization fact above.  Let \(\chi=Q_0x\,\psi\) and
\(P=\{m\in D^{\mathfrak A}:\mathfrak A^{D}\models\psi[s(m/x)]\}\), which by
the induction hypothesis is the set of \(m\in D^{\mathfrak A}\) with
\(\mathfrak A\models\psi'[s(m/x)]\).  If \(R_\psi(s(\bar y))\) holds,
(W1) gives an injection \(n\mapsto h_\psi(s(\bar y),n)\) of the infinite set
\(\mathrm{Nn}^{\mathfrak A}\) into \(P\), so \(P\) is infinite.  If it fails,
(W2) gives an injection of \(P\) into
\(\{u\in\mathrm{Nn}^{\mathfrak A}:u\mathbin{\hat<}b_\psi(s(\bar y))\}\),
which is \(j\) of the finite set of numbers below \(j^{-1}(b_\psi(s(\bar y)))\),
so \(P\) is finite.  So \(\mathfrak A\models R_\psi(\bar y)[s]\) iff
\(\mathfrak A^{D}\models Q_0x\,\psi[s]\).  Hence \(\mathfrak A^{D}\) is a model
of \(\varphi\), and \(\mathfrak A^{(U)}_{L}=(\mathfrak A^{D})^{(U)}_{L}\).

Conversely let \(\mathfrak B\models\varphi\).  Let \(\mathfrak A\) have
domain \(B\cup N'\), where \(N'=\{n':n\in\N\}\) is a set of distinct new
objects disjoint from \(B\); \(D=B\), \(\mathrm{Nn}=N'\), \(\hat0=0'\),
\(\hat S(n')=(n+1)'\), \(n'\mathbin{\hat+}m'=(n+m)'\), \(n'\mathbin{\hat\times}
m'=(nm)'\); the symbols of \(L'\) as in \(\mathfrak B\) on \(B\); and every
function symbol given a fixed value in \(B\) at arguments where nothing else
is specified.  Interpret \(R_\psi(\bar c)\) as true iff
\(\mathfrak B\models Q_0x\,\psi[\bar c]\), for \(\bar c\) in \(B\); this
determines the symbols \(R_\psi\) from the inside outward, and then
\(\mathfrak A\models\chi'[s]\) iff \(\mathfrak B\models\chi[s]\) by the
induction just given (whose inductive steps other than the \(Q_0\)-case do
not use (W1), (W2)).  For each \(\bar c\) with \(R_\psi(\bar c)\) true,
the set \(P_{\bar c}\) of witnesses is infinite and the axiom of choice gives
an injection \(e_{\bar c}\colon\N\to P_{\bar c}\); put
\(h_\psi(\bar c,n')=e_{\bar c}(n)\).  For \(\bar c\) with \(R_\psi(\bar c)\)
false, \(P_{\bar c}\) is finite with \(p\) elements; choose a bijection
\(f_{\bar c}\colon P_{\bar c}\to\{0,\ldots,p-1\}\) and put
\(g_\psi(\bar c,x)=f_{\bar c}(x)'\) for \(x\in P_{\bar c}\) and
\(b_\psi(\bar c)=p'\).  Then (W1) and (W2) hold, \(\mathfrak A\) is an
\(\omega\)-model of \(\varphi^{\omega}\), and
\(\mathfrak A^{(U)}_{L}=\mathfrak B^{(U)}_{L}\).  So \(K\) is RPC in
\(\omega\)-logic by \(\varphi^{\omega}\).

\emph{Part~4.}  Let \(K\) be RPC in \(\omega\)-logic by \(\psi\) for
\(L'\).  Let \(\gamma(i,v,w)\) be \(w=\hat Sv\) with the individuals as value
sort, \(\rho(n,x)\) the \(L'^{*}\)-formula \(\exists z\,(z=\hat0\land
\mathrm{Rec}_\gamma(n,z,x))\), and \(\nu\) the conjunction of
\(\mathrm{Nn}(\hat0)\); \(\forall x\,(\mathrm{Nn}(x)\to\mathrm{Nn}(\hat Sx))\);
\(\forall x\,\forall y\,((\mathrm{Nn}(x)\land\mathrm{Nn}(y))\to
(\mathrm{Nn}(x\mathbin{\hat+}y)\land\mathrm{Nn}(x\mathbin{\hat\times}y)))\);
\(\forall x\,(\mathrm{Nn}(x)\to\exists n\,\rho(n,x))\);
\(\forall n\,\forall m\,\forall x\,((\rho(n,x)\land\rho(m,x))\to n=m)\);
\(\forall n\,\forall m\,\forall k\,\forall x\,\forall y\,\forall z\,
((\rho(n,x)\land\rho(m,y)\land\mathrm{Add}(n,m,k)\land\rho(k,z))\to
x\mathbin{\hat+}y=z)\); and the same with \(\mathrm{Times}\) and
\(\hat\times\).  For a structure \(\mathfrak A\) for \(L'\) let
\(h(0)=\hat0^{\mathfrak A}\) and \(h(k+1)=\hat S^{\mathfrak A}h(k)\).  For
each natural number \(i\) and \(v\in A\) exactly one \(w\) satisfies
\(\gamma\), so by \cref{thm:p02-w02c-p02-003-hf-recursion},
\(\mathrm{HF}(\mathfrak A)\models\rho[n,x]\) iff \(n\) is a natural number
and \(x=h(n)\).  By (M2) and (M3) of
\cref{annex:p02-w02c-p02-003-wso-sentences}, for natural numbers \(n,m\),
\(\mathrm{Add}[n,m,k]\) holds iff \(k=n+m\), and
\(\mathrm{Times}[n,m,k]\) iff \(k=nm\).

If \(\mathrm{HF}(\mathfrak A)\models\nu\), then \(h(k)\in
\mathrm{Nn}^{\mathfrak A}\) for all \(k\) by the first two conjuncts; \(h\)
maps onto \(\mathrm{Nn}^{\mathfrak A}\) by the fourth and is one-one by the
fifth; \(h(0)=\hat0\) and \(h(k+1)=\hat S h(k)\); and taking \(k=n+m\),
respectively \(k=nm\), the last two conjuncts give
\(h(n)\mathbin{\hat+}h(m)=h(n+m)\) and \(h(n)\mathbin{\hat\times}h(m)=h(nm)\).
So \(\mathrm{Nn}^{\mathfrak A}\) is closed and \(h\) is an isomorphism of
\(\mathcal N\) onto it with the restricted operations.  Conversely, if
\(\mathrm{Nn}^{\mathfrak A}\) is closed and \(g\) is such an isomorphism,
then \(g(k)=h(k)\) by induction on \(k\), and each conjunct of \(\nu\)
holds: the closure conjuncts by closure, the fourth as \(g\) is onto, the
fifth as \(g\) is one-one, and the last two because \(\rho(n,x)\),
\(\rho(m,y)\) and \(\mathrm{Add}(n,m,k)\) force \(x=g(n)\), \(y=g(m)\),
\(k=n+m\), so that \(z=g(n+m)=x\mathbin{\hat+}y\), and likewise for
\(\mathrm{Times}\).  So the models of \(\psi\land\nu\) in weak second-order
logic are, by (F), exactly the \(\omega\)-models of \(\psi\), and \(K\) is
RPC in weak second-order logic by \(\psi\land\nu\).

\emph{Conclusion.}  Parts~1 and~2 give the same RPC classes for
\(L(Q_0)\) and weak second-order logic; parts~3 and~4 with part~2 give
RPC in \(L(Q_0)\) \(\subseteq\) RPC in \(\omega\)-logic \(\subseteq\) RPC in
weak second-order logic \(\subseteq\) RPC in \(L(Q_0)\).  The choices of
\(e_{\bar c}\) and \(f_{\bar c}\) in part~3 are the only uses of the axiom of
choice.
\end{proof}
\end{annexitem}

\begin{annexitem}{annex:p02-w02c-p02-003-kq-soundness}{thm:p02-w02c-p02-003-kq-soundness}{Substitution and soundness for \(L(Q)\)}
\begin{proof}
The notation is that of
\cref{def:p02-w02c-p02-003-kq-language,def:p02-w02c-p02-003-kq-calculus};
\(\varphi[s]\) is satisfaction at an assignment \(s\), and \(s(a/x)\) is as
at the head of this annex.

\emph{Part~1.}  Coincidence is proved by induction on \(\varphi\), for all
assignments at once.  The first-order clauses are as in first-order logic.
For \((Qy)\psi\), the set \(\{b:(\mathfrak A,q)\models\psi[s(b/y)]\}\)
depends, by the induction hypothesis, only on the values of \(s\) at the
free variables of \(\psi\) other than \(y\), which are free in
\((Qy)\psi\); so whether it belongs to \(q\) depends only on these values.

For substitution, first \(u(t/x)^{\mathfrak A}[s]=
u^{\mathfrak A}[s(t^{\mathfrak A}[s]/x)]\) for every term \(u\), by
induction on \(u\).  This gives the atomic case, and the cases \(\neg\),
\(\land\) are immediate.  Let \(\varphi=(Ky)\psi\) with
\(K\in\{\exists,\forall,Q\}\).  If \(x\) is not free in \(\varphi\), then
\(\varphi(t/x)=\varphi\) and the two assignments agree on the free variables
of \(\varphi\), so coincidence gives the claim.  Otherwise \(y\neq x\),
\(x\) is free in \(\psi\), and substitutability gives that \(y\) does not
occur in \(t\) and \(t\) is substitutable for \(x\) in \(\psi\); and
\(\varphi(t/x)=(Ky)\psi(t/x)\).  Put \(e=t^{\mathfrak A}[s]\).  For each
\(b\in A\), \(t^{\mathfrak A}[s(b/y)]=e\) since \(y\) is not in \(t\), so by
the induction hypothesis \(\psi(t/x)[s(b/y)]\) holds iff
\(\psi[s(b/y)(e/x)]\) holds, and \(s(b/y)(e/x)=s(e/x)(b/y)\) since
\(x\neq y\).  Hence
\(\{b:\psi(t/x)[s(b/y)]\}=\{b:\psi[s(e/x)(b/y)]\}\), and the clause for
\(K\) (some member, every member, or membership of this set in \(q\)) gives
the same value on both sides.

\emph{Part~2.}  (0a): the truth value of an instance is computed from the
truth values of the substituted formulas by the clauses for \(\neg\) and
\(\land\), so it is that of the tautology under a valuation, which is
truth.  (0b): if \((\forall x)\varphi\) holds at \(s\), then \(\varphi\)
holds at \(s(t^{\mathfrak A}[s]/x)\), hence \(\varphi(t/x)\) at \(s\) by
part~1.  (0c) is the dual argument.  (0d): if \((\forall x)(\psi\to\varphi)\)
and \(\psi\) hold at \(s\), then for each \(b\), \(\psi\) holds at
\(s(b/x)\) by coincidence (\(x\) is not free in \(\psi\)), so \(\varphi\)
does; hence \((\forall x)\varphi\) holds at \(s\).  (0e): if
\((\forall x)(\varphi\to\psi)\) and \((\exists x)\varphi\) hold at \(s\),
some \(b\) satisfies \(\varphi\) at \(s(b/x)\), hence \(\psi\), and by
coincidence \(\psi\) holds at \(s\).  (0f) holds since \(=\) is interpreted
by identity.  (0g): if \(s(x)=s(y)\), every term in \(\alpha\) has the same
value as the corresponding term in \(\alpha'\) (induction on terms), so
\(\alpha\) and \(\alpha'\) have the same truth value.  Axiom~3: \(y\) does
not occur in \(\varphi\), so it is substitutable for \(x\); by part~1 and
coincidence (\(y\) is not free in \(\varphi\)),
\(\{b:\varphi(y/x)[s(b/y)]\}=\{b:\varphi[s(b/y)(b/x)]\}=\{b:\varphi[s(b/x)]\}\),
so \((Qx)\varphi\) and \((Qy)\varphi(y/x)\) have the same value at \(s\).

\emph{Part~3.}  By induction along a deduction \(\varphi_1,\ldots,\varphi_m\)
from \(\Sigma\), each \(\varphi_i\) holds at every assignment of
\((\mathfrak A,q)\).  Axioms do, by hypothesis (their universal closures
hold); members of \(\Sigma\) do; if \(\varphi\) and \(\varphi\to\psi\)
hold at every assignment, so does \(\psi\); if \(\varphi\) holds at every
assignment, so does \((\forall x)\varphi\).

\emph{Part~4.}  Let \(q\) be the set of uncountable subsets of \(A\).
Axioms 0 and 3 hold by part~2.  Axiom~1: the set
\(\{b:b=s(y)\text{ or }b=s(z)\}\) has at most two elements, so is not in
\(q\).  Axiom~2: if \((\forall x)(\varphi\to\psi)\) holds at \(s\), the set
defined by \(\varphi\) at \(s\) is included in that defined by \(\psi\), and
a set with an uncountable subset is uncountable.  Axiom~4: suppose
\((\exists x)(Qy)\varphi\) and \((Qx)(\exists y)\varphi\) fail at \(s\).
Then \(X=\{a:(\exists y)\varphi[s(a/x)]\}\) is countable and, for each
\(a\), \(Y_a=\{b:\varphi[s(a/x)(b/y)]\}\) is countable.  If
\((\exists x)\varphi\) holds at \(s(b/y)\), some \(a\) has
\(\varphi[s(b/y)(a/x)]\), that is \(\varphi[s(a/x)(b/y)]\) (\(x\neq y\)),
so \(a\in X\) and \(b\in Y_a\).  Thus
\(\{b:(\exists x)\varphi[s(b/y)]\}\subseteq\bigcup_{a\in X}Y_a\), a
countable union of countable sets, which is countable by the axiom of
choice; so \((Qy)(\exists x)\varphi\) fails at \(s\).  Hence every axiom
holds in \((\mathfrak A,q)\), and part~3 gives the last statement.  If
\(\Sigma\vdash\sigma\) and \(\Sigma\vdash\neg\sigma\), both would hold in a
standard model of \(\Sigma\); so a set with a standard model is consistent.
\end{proof}
\end{annexitem}

\begin{annexitem}{annex:p02-w02c-p02-003-kq-proof-lemmas}{thm:p02-w02c-p02-003-kq-proof-lemmas}{Proof-theoretic lemmas for \(L(Q)\)}
\begin{proof}
``By (0a)'' below means: by an instance of (0a) and modus ponens.

\emph{Part~1.}  Let \(\varphi_1,\ldots,\varphi_m\) be a deduction from
\(\Sigma\cup\{\sigma\}\); we show \(\Sigma\vdash\sigma\to\varphi_i\) by
induction on \(i\).  If \(\varphi_i\) is an axiom or in \(\Sigma\), use
\(\varphi_i\) and the (0a) instance \(\varphi_i\to(\sigma\to\varphi_i)\).
If \(\varphi_i=\sigma\), \(\sigma\to\sigma\) is a (0a) instance.  If
\(\varphi_i\) follows by modus ponens from \(\varphi_j\) and
\(\varphi_j\to\varphi_i\), use the (0a) instance
\((\sigma\to(\varphi_j\to\varphi_i))\to((\sigma\to\varphi_j)\to(\sigma\to\varphi_i))\)
twice.  If \(\varphi_i=(\forall x)\varphi_j\), generalization gives
\(\Sigma\vdash(\forall x)(\sigma\to\varphi_j)\), and (0d), with \(x\) not
free in the sentence \(\sigma\), gives \(\Sigma\vdash\sigma\to(\forall
x)\varphi_j\).

\emph{Part~2.}  The first two rules are generalization followed by (0d),
respectively (0e), and modus ponens.  For the third, (0a) gives
\(\Sigma\vdash\varphi\to\varphi'\) and \(\Sigma\vdash\varphi'\to\varphi\).
For \(\forall\): (0b) with \(t=x\) gives \((\forall x)\varphi\to\varphi\),
so \(\Sigma\vdash(\forall x)\varphi\to\varphi'\) by (0a), and the first rule
(\(x\) is not free in \((\forall x)\varphi\)) gives
\(\Sigma\vdash(\forall x)\varphi\to(\forall x)\varphi'\).  For \(\exists\):
(0c) with \(t=x\) gives \(\varphi'\to(\exists x)\varphi'\), so
\(\Sigma\vdash\varphi\to(\exists x)\varphi'\), and the second rule gives
\(\Sigma\vdash(\exists x)\varphi\to(\exists x)\varphi'\).  For \(Q\):
generalization gives \(\Sigma\vdash(\forall x)(\varphi\to\varphi')\), and
Axiom~2 gives \(\Sigma\vdash(Qx)\varphi\to(Qx)\varphi'\).  The converse
implications are proved in the same three ways with \(\varphi\) and
\(\varphi'\) exchanged, and (0a) combines the two implications.

\emph{Part~3.}  Write \(\psi'=\psi(z/y)\).  First,
\(\vdash(Ky)\psi\leftrightarrow(Kz)\psi'\).  For \(K=Q\) this is an
instance of Axiom~3.  For \(K=\forall\): \(z\) does not occur in \(\psi\),
so it is substitutable for \(y\), and (0b) gives \((\forall
y)\psi\to\psi'\); \(z\) is not free in \((\forall y)\psi\), so the first
rule of part~2 gives \((\forall y)\psi\to(\forall z)\psi'\).  Conversely,
the free occurrences of \(z\) in \(\psi'\) stand exactly where \(y\) was
free in \(\psi\), so they are not in the scope of a quantifier on \(y\), and
\(y\) is substitutable for \(z\) in \(\psi'\), with \(\psi'(y/z)=\psi\)
since \(z\) does not occur in \(\psi\); (0b) gives
\((\forall z)\psi'\to\psi\), and since \(y\) is not free in \((\forall
z)\psi'\), the first rule gives \((\forall z)\psi'\to(\forall y)\psi\).
For \(K=\exists\) the same argument uses (0c) and the second rule.  Next,
by induction on \(\varphi\): if the replaced occurrence is \(\varphi\)
itself, this is the first step; if \(\varphi\) is \(\neg\chi\),
\(\chi\land\theta\) or \(\theta\land\chi\) with the occurrence in \(\chi\),
the induction hypothesis \(\vdash\chi\leftrightarrow\chi'\) and (0a) give
the claim; if \(\varphi=(Kx)\chi\), part~2 with \(\Sigma=\varnothing\) gives
it.

\emph{Part~4.}  First take distinct \(u_1,\ldots,u_k\) not occurring in a
given deduction \(\varphi_1,\ldots,\varphi_m\) of \(\chi\) from \(\Sigma\),
and let \(\varphi'\) denote the result of replacing each occurrence of
\(c_i\) by \(u_i\).  Then \(\varphi'_1,\ldots,\varphi'_m\) is a deduction of
\(\chi'\) from \(\Sigma\).  Members of \(\Sigma\) contain no \(c_i\) and are
unchanged.  The replacement commutes with the connectives and quantifiers
and with substitution: \(\varphi(t/x)'=\varphi'(t'/x)\), since no \(u_i\)
equals \(x\).  So an instance of (0a) goes to an instance of (0a); for (0b)
and (0c), \(t'\) is substitutable for \(x\) in \(\varphi'\), because the
variables of \(t'\) are those of \(t\) and some \(u_i\), and no \(u_i\) is
bound in \(\varphi'\); for (0d), (0e), \(x\) is not free in \(\psi'\), whose
free variables are those of \(\psi\) and some \(u_i\neq x\); (0f) and
Axiom~1 contain no constants; for (0g), \(\alpha'\) is atomic and the
replacement of occurrences of \(x\) by \(y\) is preserved; Axioms~2 and~4
go to instances; for Axiom~3, \(y\neq u_i\), so \(y\) still does not occur
in \(\varphi'\).  Modus ponens and generalization go to the same rules.
Now let \(v_1,\ldots,v_k\) be distinct and not occurring in \(\chi\), and
choose the \(u_i\) also distinct from the \(v_j\).  Then
\(\Sigma\vdash(\forall u_1)\cdots(\forall u_k)\chi'\) by generalization,
and \(k\) applications of (0b), substituting \(v_1\) for \(u_1\), \ldots,
\(v_k\) for \(u_k\), give \(\chi(v_1/c_1,\ldots,v_k/c_k)\): each \(v_i\) is
substitutable since it does not occur in \(\chi\) and differs from the
\(u_j\).

For consistency: a deduction in \(L(Q)\) is one in \(L'(Q)\).  Conversely,
if \(\Sigma\vdash\sigma\) and \(\Sigma\vdash\neg\sigma\) in \(L'(Q)\), then
by (0a) \(\Sigma\vdash\tau\) and \(\Sigma\vdash\neg\tau\) for the
\(L\)-sentence \(\tau=(\forall z)(z=z)\); replacing in these two
deductions every new constant by a variable not occurring in them, as
above, gives deductions in \(L(Q)\) of \(\tau\) and \(\neg\tau\), which
contain no new constants.

\emph{Part~5.}  If \(\Gamma\vdash\sigma\) and \(\sigma\notin\Gamma\), then
\(\Gamma\cup\{\sigma\}\) is inconsistent by maximality, so it proves some
\(\tau\) and \(\neg\tau\); part~1 and (0a) give \(\Gamma\vdash\neg\sigma\),
and \(\Gamma\) is inconsistent, a contradiction.  Not both \(\sigma\) and
\(\neg\sigma\) are in \(\Gamma\); if \(\sigma\notin\Gamma\), the same
argument gives \(\Gamma\vdash\neg\sigma\), so \(\neg\sigma\in\Gamma\).  The
clause for \(\land\) follows from the first by (0a).
\end{proof}
\end{annexitem}

\begin{annexitem}{annex:p02-w02c-p02-003-kq-weak-completeness}{thm:p02-w02c-p02-003-kq-weak-completeness}{Weak completeness and omitting types}
\begin{proof}
Two derived facts are used.  (E) For every formula \(\theta\) and
variables \(x\neq y\), if \(\theta'\) comes from \(\theta\) by replacing
some free occurrences of \(x\) at which \(y\) would be free, then
\(\vdash x=y\to(\theta\leftrightarrow\theta')\).  By induction on
\(\theta\): for atomic \(\theta\), (0g) gives \(x=y\to(\theta\to\theta')\);
(0g) with the roles of \(x,y\) exchanged gives
\(y=x\to(\theta'\to\theta)\), and \(x=y\to y=x\) follows from the (0g)
instance \(x=y\to(x=x\to y=x)\) and (0f).  The steps for \(\neg,\land\) use
(0a).  For \((Kz)\chi\): if \(z\) is \(x\) or \(y\), no occurrence inside is
replaced (an occurrence of \(x\) inside is bound, or \(y\) would not be
free there), and \(\theta'=\theta\).  Otherwise the induction hypothesis,
generalization and (0d) (\(z\) is not free in \(x=y\)) give
\(x=y\to(\forall z)(\chi\leftrightarrow\chi')\), and
\((\forall z)(\chi\leftrightarrow\chi')\to((Kz)\chi\leftrightarrow(Kz)\chi')\)
is provable: by (0b) and (0d) for \(\forall\), by (0b), (0c), (0e) for
\(\exists\), and by (0b), (0d) and Axiom~2 in both directions for \(Q\).
(N) \(\vdash\neg(\forall x)\psi\to(\exists x)\neg\psi\): (0c) gives
\(\neg\psi\to(\exists x)\neg\psi\), (0a) gives
\(\neg(\exists x)\neg\psi\to\psi\), the first rule of part~2 of
\cref{thm:p02-w02c-p02-003-kq-proof-lemmas} gives
\(\neg(\exists x)\neg\psi\to(\forall x)\psi\), and (0a) contraposes.

\emph{Part~1.}  Let \(C=\{c_0,c_1,\ldots\}\) be new constants and
\(L^{C}=L\cup C\).  By part~4 of \cref{thm:p02-w02c-p02-003-kq-proof-lemmas},
\(\Sigma\) is consistent in \(L^{C}(Q)\).  \(L^{C}\) is countable, so its
sentences can be listed as \(\chi_0,\chi_1,\ldots\), and the pairs
\((c,n)\in C\times N\) as \(p_0,p_1,\ldots\) (when \(N\) is empty, step
(c) below is skipped).  We build \(\Gamma_k=\Sigma\cup F_k\) with \(F_k\)
finite and \(\Gamma_k\) consistent, starting with \(F_0=\varnothing\).
Given \(\Gamma_k\):

(a) Let \(\Gamma'=\Gamma_k\cup\{\chi_k\}\) if this is consistent and
\(\Gamma'=\Gamma_k\cup\{\neg\chi_k\}\) otherwise.  If both were
inconsistent, part~1 of \cref{thm:p02-w02c-p02-003-kq-proof-lemmas} and
(0a) would give \(\Gamma_k\vdash\neg\chi_k\) and
\(\Gamma_k\vdash\neg\neg\chi_k\).

(b) If \(\chi_k=(\exists x)\varphi\), choose \(c\in C\) not occurring in
\(\Gamma'\) (only finitely many members of \(C\) occur in \(F\)) and let
\(\Gamma''=\Gamma'\cup\{(\exists x)\varphi\to\varphi(c/x)\}\); otherwise
\(\Gamma''=\Gamma'\).  Suppose \(\Gamma''\) is inconsistent.  As in (a),
\(\Gamma'\vdash\neg((\exists x)\varphi\to\varphi(c/x))\), so by (0a)
\(\Gamma'\vdash(\exists x)\varphi\) and \(\Gamma'\vdash\neg\varphi(c/x)\).
Take \(v\) not occurring in \(\varphi\); part~4 of
\cref{thm:p02-w02c-p02-003-kq-proof-lemmas} gives
\(\Gamma'\vdash\neg\varphi(v/x)\), generalization and (0b) with \(t=x\)
(substitutable, since \(v\) stands exactly where \(x\) was free) give
\(\Gamma'\vdash\neg\varphi\), (0a) gives
\(\Gamma'\vdash\varphi\to\neg(\exists x)\varphi\), and the second rule of
part~2 gives \(\Gamma'\vdash\neg(\exists x)\varphi\): \(\Gamma'\) is
inconsistent, a contradiction.

(c) Let \(p_k=(c,n)\).  We find \(\sigma\in\Sigma_n\) with
\(\Gamma''\cup\{\neg\sigma(c/x_n)\}\) consistent and put \(\Gamma_{k+1}\)
equal to this set.  Suppose there is none.  Let \(\delta\) be the
conjunction of \(c=c\) and the members of \(\Gamma''-\Sigma\), and let
\(d_1,\ldots,d_m\) be the members of \(C\) other than \(c\) that occur in
\(\delta\).  Choose distinct variables \(y_1,\ldots,y_m\), different from
\(x_n\) and not occurring in \(\delta\), and by part~3 of
\cref{thm:p02-w02c-p02-003-kq-proof-lemmas} an alphabetic variant
\(\delta^{\circ}\) of \(\delta\) in which neither \(x_n\) nor any \(y_i\)
occurs, with \(\vdash\delta\leftrightarrow\delta^{\circ}\).  Let
\(\varphi(x_n)\) be \((\exists y_1)\cdots(\exists y_m)\delta^{\circ}_1\),
where \(\delta^{\circ}_1\) replaces \(c\) by \(x_n\) and each \(d_i\) by
\(y_i\) in \(\delta^{\circ}\); it is a formula of \(L(Q)\) with only \(x_n\)
free.  \(\Gamma''\vdash\delta^{\circ}\), and \(m+1\) applications of (0c),
each with a closed term, give \(\Gamma''\vdash(\exists x_n)\varphi\); so
\((\exists x_n)\varphi\) is consistent with \(\Sigma\), and by hypothesis
some \(\sigma\in\Sigma_n\) makes \((\exists x_n)(\varphi\land\neg\sigma)\)
consistent with \(\Sigma\).  By our assumption \(\Gamma''\vdash\sigma(c/x_n)\);
repeated use of part~1 of \cref{thm:p02-w02c-p02-003-kq-proof-lemmas} and
(0a) give \(\Sigma\vdash\delta^{\circ}\to\sigma(c/x_n)\).  The constants
\(c,d_1,\ldots,d_m\) do not occur in \(\Sigma\); with distinct variables \(u\),
\(w_1,\ldots,w_m\) not occurring in this formula and different from \(x_n\)
and the \(y_i\), part~4 gives
\(\Sigma\vdash\delta^{\circ}(u/c,\bar w/\bar d)\to\sigma(u/x_n)\).
Generalization on \(u\) and (0b) with \(t=x_n\) (substitutable: \(x_n\) is
not bound in \(\delta^{\circ}\), and \(u\) stands in \(\sigma(u/x_n)\) where
\(x_n\) was free), then the same with each \(w_i\) and \(t=y_i\), give
\(\Sigma\vdash\delta^{\circ}_1\to\sigma\).  Since no \(y_i\) is free in
\(\sigma\), the second rule of part~2 gives \(\Sigma\vdash\varphi\to\sigma\);
with \(\bot=\neg(\forall z)(z=z)\), (0a) gives
\(\Sigma\vdash(\varphi\land\neg\sigma)\to\bot\), the second rule gives
\(\Sigma\vdash(\exists x_n)(\varphi\land\neg\sigma)\to\bot\), and (0f) with
generalization gives \(\Sigma\vdash\neg\bot\).  So
\(\Sigma\cup\{(\exists x_n)(\varphi\land\neg\sigma)\}\) is inconsistent, a
contradiction.

Let \(\Gamma^{*}=\bigcup_k\Gamma_k\).  It is consistent, since a deduction
uses finitely many members, all in some \(\Gamma_k\); maximal, since each
sentence \(\chi_k\) or its negation is in \(\Gamma^{*}\); \(C\) is a set of
witnesses for it, by (b); and for each \(c\in C\) and \(n\in N\) some
\(\sigma\in\Sigma_n\) has \(\neg\sigma(c/x_n)\in\Gamma^{*}\), by (c).  Part~5
of \cref{thm:p02-w02c-p02-003-kq-proof-lemmas} is used below without
mention.

\emph{The model (Keisler's Lemma~2.2).}  For closed terms \(s,t,u\) of
\(L^{C}\), \(\Gamma^{*}\) contains \(t=t\), \(s=t\to t=s\),
\(s=t\to(t=u\to s=u)\), and \(s=t\to(\alpha\leftrightarrow\alpha')\) for
every sentence \(\alpha\) of \(L^{C}(Q)\) and \(\alpha'\) obtained by
replacing one occurrence of \(s\) by \(t\): these come from (0f), (0g) and
(E), written with variables at the relevant places, by generalization and
(0b) with closed terms.  For each closed term \(t\), (0c) applied to
\(t=t\) gives \((\exists x)(t=x)\in\Gamma^{*}\), so some \(c\in C\) has
\(t=c\in\Gamma^{*}\).  Let \(c\sim d\) iff \(c=d\in\Gamma^{*}\), an
equivalence relation on \(C\), and \(A=\{[c]:c\in C\}\).  Interpret a
constant \(e\) of \(L\) as \([c]\) where \(e=c\in\Gamma^{*}\), a function
symbol by \(f^{\mathfrak A}([c_1],\ldots,[c_n])=[d]\) where
\(f(c_1,\ldots,c_n)=d\in\Gamma^{*}\), and a relation symbol by
\(([c_1],\ldots,[c_n])\in R^{\mathfrak A}\) iff
\(R(c_1,\ldots,c_n)\in\Gamma^{*}\); the equality facts make these
independent of the representatives, and \(c_i\in C\) is interpreted as
\([c_i]\).  By induction on terms, \(t^{\mathfrak A}=[d]\) iff
\(t=d\in\Gamma^{*}\); hence an atomic sentence \(R(t_1,\ldots,t_n)\) or
\(t_1=t_2\) is in \(\Gamma^{*}\) iff it is true in \(\mathfrak A\).  For a
formula \(\psi\) with only \(x\) free, let
\(S_\psi=\{[c]:\psi(c/x)\in\Gamma^{*}\}\), and let
\[
 q=\{S_\psi:\psi\text{ has only one free variable }x\text{ and }(Qx)\psi\in\Gamma^{*}\}.
\]
By (E), if \(c\sim c'\) then \(\psi(c/x)\in\Gamma^{*}\) iff
\(\psi(c'/x)\in\Gamma^{*}\), so \([c]\in S_\psi\) iff \(\psi(c/x)\in\Gamma^{*}\).

(\(\star\)) If \(\psi(c/x)\in\Gamma^{*}\) for every \(c\in C\), then
\((\forall x)\psi\in\Gamma^{*}\).  Otherwise \(\neg(\forall
x)\psi\in\Gamma^{*}\), so \((\exists x)\neg\psi\in\Gamma^{*}\) by (N), and a
witness gives \(\neg\psi(c/x)\in\Gamma^{*}\).

We show, by induction on the number of occurrences of
\(\neg,\land,\exists,\forall,Q\), that for every sentence \(\sigma\) of
\(L^{C}(Q)\), \((\mathfrak A,q)\models\sigma\) iff \(\sigma\in\Gamma^{*}\);
\(\psi(c/x)\) has as many occurrences as \(\psi\).  Atomic sentences and
the steps for \(\neg,\land\) are done.  For each \(\psi\) with only \(x\)
free, part~1 of \cref{thm:p02-w02c-p02-003-kq-soundness} and the induction
hypothesis give \((\mathfrak A,q)\models\psi[[c]]\) iff
\(\psi(c/x)\in\Gamma^{*}\), so the set defined by \(\psi\) is \(S_\psi\).
\((\exists x)\psi\): it is true iff some \(\psi(c/x)\in\Gamma^{*}\); if so,
(0c) gives \((\exists x)\psi\in\Gamma^{*}\), and conversely a witness gives
some \(\psi(c/x)\in\Gamma^{*}\).  \((\forall x)\psi\): it is true iff every
\(\psi(c/x)\in\Gamma^{*}\); (0b) and (\(\star\)) give the equivalence with
\((\forall x)\psi\in\Gamma^{*}\).  \((Qx)\psi\): if
\((Qx)\psi\in\Gamma^{*}\), then \(S_\psi\in q\) and \((Qx)\psi\) is true.
If \((Qx)\psi\) is true, \(S_\psi\in q\), so \(S_\psi=S_\theta\) for some
\(\theta\) with only \(y\) free and \((Qy)\theta\in\Gamma^{*}\); then
\(\psi(c/x)\in\Gamma^{*}\) iff \(\theta(c/y)\in\Gamma^{*}\) for every
\(c\in C\).  Take \(u\) not occurring in \(\psi,\theta\).  The sentence
\(\psi(c/x)\leftrightarrow\theta(c/y)\), which is
\((\psi(u/x)\leftrightarrow\theta(u/y))(c/u)\), is in \(\Gamma^{*}\) for
every \(c\), so by (\(\star\))
\((\forall u)(\psi(u/x)\leftrightarrow\theta(u/y))\in\Gamma^{*}\).  By (0b),
(0d) and Axiom~2, \((Qu)\theta(u/y)\to(Qu)\psi(u/x)\in\Gamma^{*}\); by
Axiom~3, \((Qy)\theta\leftrightarrow(Qu)\theta(u/y)\) and
\((Qx)\psi\leftrightarrow(Qu)\psi(u/x)\) are in \(\Gamma^{*}\); hence
\((Qx)\psi\in\Gamma^{*}\).

Every axiom of \(L^{C}(Q)\) has a provable universal closure
(generalization), which is in \(\Gamma^{*}\) and so true; hence every axiom
holds in \((\mathfrak A,q)\), and \(\Sigma\subseteq\Gamma^{*}\) holds.  The
reduct to \(L\), with the same \(q\), gives every sentence of \(L(Q)\) the
same value; so it is a weak model of \(\Sigma\) in which all the axioms of
\(L(Q)\) hold.  \(A\) is the image of the countable set \(C\).  For
\([c]\in A\) and \(n\in N\), some \(\sigma\in\Sigma_n\) has
\(\neg\sigma(c/x_n)\in\Gamma^{*}\), so \(\sigma\) fails at \([c]\) (by part~1
of \cref{thm:p02-w02c-p02-003-kq-soundness}); hence \(\Sigma_n\) is
omitted.

\emph{Part~2.}  If \(\Sigma\) is consistent, part~1 with \(N=\varnothing\)
gives the model.  If such a model exists and \(\Sigma\vdash\sigma\),
\(\Sigma\vdash\neg\sigma\), both hold in it by part~3 of
\cref{thm:p02-w02c-p02-003-kq-soundness}, which is impossible.

\emph{Part~3.}  One direction is part~3 of
\cref{thm:p02-w02c-p02-003-kq-soundness}.  If \(\Sigma\not\vdash\sigma\),
then \(\Sigma\cup\{\neg\sigma\}\) is consistent: otherwise it proves some
\(\tau\) and \(\neg\tau\), and part~1 of
\cref{thm:p02-w02c-p02-003-kq-proof-lemmas} and (0a) give
\(\Sigma\vdash\sigma\).  Part~2 gives a weak model of
\(\Sigma\cup\{\neg\sigma\}\) in which all the axioms hold and \(\sigma\)
fails.
\end{proof}
\end{annexitem}

\begin{annexitem}{annex:p02-w02c-p02-003-kq-finite-unions}{thm:p02-w02c-p02-003-kq-finite-unions}{Finite unions of small definable sets}
\begin{proof}
By part~1 of \cref{thm:p02-w02c-p02-003-kq-soundness}, replacing the
distinguished variable of a defining formula by a variable not occurring in
it does not change the defined set; so any finitely many definable sets
can be given defining formulas with a common distinguished variable \(x\)
and with their parameter variables distinct from each other and from the
variables used below.  Fix an assignment \(s\) giving the parameters their
values.

(a) Let \(X,Y\) be defined at \(s\) by \(\varphi,\psi\) with \(X\subseteq
Y\) and \(X\in q\).  Then \((\forall x)(\varphi\to\psi)\) and
\((Qx)\varphi\) hold at \(s\), and Axiom~2 gives \((Qx)\psi\) at \(s\), that
is \(Y\in q\).

(b) For \(a_1,a_2\in A\), Axiom~1 at \(s(a_1/y)(a_2/z)\) says that
\(\{a_1,a_2\}\), the set defined by \(x=y\lor x=z\), is not in \(q\).  A
definable set with at most two elements is included in such a set
(\(A\neq\varnothing\)), so it is not in \(q\) by (a); in particular
\(\varnothing\), defined by \(\neg x=x\), is not in \(q\).

(c) Let \(X,Y\) be defined at \(s\) by \(\varphi,\psi\) and suppose
\(X\cup Y\in q\).  By (b), \(A\) has two distinct elements \(a_1,a_2\).  Let
\(y,y_1,y_2\) be new variables, \(s'=s(a_1/y_1)(a_2/y_2)\) and
\(\theta=(y=y_1\land\varphi)\lor(y=y_2\land\psi)\).  The set defined by
\((\exists y)\theta\) at \(s'\) is \(X\cup Y\), so \((Qx)(\exists y)\theta\)
holds at \(s'\).  The instance
\((Qx)(\exists y)\theta\to(\exists y)(Qx)\theta\lor(Qy)(\exists x)\theta\)
of Axiom~4 (with the variables named \(x,y\) there exchanged) holds.  The
set \(\{c:(\exists x)\theta[s'(c/y)]\}\) is included in \(\{a_1,a_2\}\), so
it is not in \(q\) by (a) and (b).  Hence some \(c\) has
\(\{b:\theta[s'(c/y)(b/x)]\}\in q\).  This set is \(X\) if \(c=a_1\), \(Y\)
if \(c=a_2\), and \(\varnothing\) otherwise; the last is excluded by (b).

The final statement follows from (c) by induction on the number of sets
and from (a).
\end{proof}
\end{annexitem}

\begin{annexitem}{annex:p02-w02c-p02-003-kq-chains}{thm:p02-w02c-p02-003-kq-chains}{Elementary chains of weak models}
\begin{proof}
\emph{Part~1.}  Elementary extension includes the atomic formulas, so each
\(\mathfrak A_\alpha\) is a substructure of \(\mathfrak A_\beta\) for
\(\alpha<\beta\), and the union structure \(\mathfrak A\) is defined; \(q\)
is a set of subsets of \(A\).  We show by induction on
\(\varphi(x_1,\ldots,x_n)\): for all \(\alpha<\gamma\) and
\(\bar a\in A_\alpha\), \((\mathfrak A_\alpha,q_\alpha)\models\varphi[\bar
a]\) iff \((\mathfrak A,q)\models\varphi[\bar a]\).  Atomic formulas:
\(\mathfrak A_\alpha\) is a substructure of \(\mathfrak A\).  \(\neg,\land\):
immediate.  \((\exists x)\psi\): a witness in \(A_\alpha\) is a witness in
\(A\) by the induction hypothesis; a witness \(b\in A\) lies in some
\(A_\beta\) with \(\beta\ge\alpha\), satisfies \(\psi\) in
\((\mathfrak A_\beta,q_\beta)\) by the induction hypothesis, so
\((\exists x)\psi[\bar a]\) holds there, hence in
\((\mathfrak A_\alpha,q_\alpha)\) by elementarity.  \((\forall x)\psi\):
if it holds in \((\mathfrak A,q)\), it holds in
\((\mathfrak A_\alpha,q_\alpha)\) by the induction hypothesis; if it holds
in \((\mathfrak A_\alpha,q_\alpha)\) and \(b\in A_\beta\), \(\beta\ge\alpha\),
it holds in \((\mathfrak A_\beta,q_\beta)\), so \(\psi[b,\bar a]\) holds
there and in \((\mathfrak A,q)\).  \((Qx)\psi\): let
\(S=\{b\in A:(\mathfrak A,q)\models\psi[b,\bar a]\}\).  For
\(\beta\ge\alpha\), the induction hypothesis gives
\(S\cap A_\beta=\{b\in A_\beta:(\mathfrak A_\beta,q_\beta)\models\psi[b,\bar
a]\}\), and elementarity gives
\((\mathfrak A_\alpha,q_\alpha)\models(Qx)\psi[\bar a]\) iff
\((\mathfrak A_\beta,q_\beta)\models(Qx)\psi[\bar a]\) iff
\(S\cap A_\beta\in q_\beta\).  If the left side holds, \(S\cap A_\beta\in
q_\beta\) for all \(\beta\ge\alpha\), so \(S\in q\).  If \(S\in q\), there is
\(\beta_0\) with \(S\cap A_\beta\in q_\beta\) for \(\beta\ge\beta_0\); for
\(\beta=\max(\alpha,\beta_0)\) this gives the left side.  This is the
argument of \cite[\S2, Lemma~2.5, pp.~17--18]{Keisler1970}.

\emph{Part~2.}  Let \(\varphi\) be an axiom of \(L^{*}(Q)\), let
\(c_{a_1},\ldots,c_{a_n}\) be the new constants in it, and let \(\psi\)
replace them by distinct variables \(v_1,\ldots,v_n\) not occurring in
\(\varphi\).  By the check in part~4 of
\cref{thm:p02-w02c-p02-003-kq-proof-lemmas}, \(\psi\) is an axiom of
\(L(Q)\), so it holds in \((\mathfrak A,q)\) at every assignment, in
particular at \(s(a_1/v_1)\cdots(a_n/v_n)\); since \(\varphi\) is
\(\psi(c_{a_1}/v_1,\ldots,c_{a_n}/v_n)\) and \(c_{a_i}\) denotes \(a_i\),
part~1 of \cref{thm:p02-w02c-p02-003-kq-soundness} gives
\((\mathfrak A^{*},q)\models\varphi[s]\) for every \(s\)
\cite[\S2, Lemma~2.6, p.~18]{Keisler1970}.

\emph{Part~3.}  A formula of \(L^{*}(Q)\) is
\(\psi(c_{a_1}/v_1,\ldots,c_{a_n}/v_n)\) with \(\psi\) in \(L(Q)\) and
\(v_i\) not in it.  By part~1 of \cref{thm:p02-w02c-p02-003-kq-soundness},
for \(\bar b\in A\) it holds at \(\bar b\) in \((\mathfrak A^{*},q)\) iff
\(\psi[\bar b,\bar a]\) holds in \((\mathfrak A,q)\), iff it holds in
\((\mathfrak B,r)\), iff the formula holds at \(\bar b\) in
\((\mathfrak B',r)\).
\end{proof}
\end{annexitem}

\begin{annexitem}{annex:p02-w02c-p02-003-kq-main-lemma}{thm:p02-w02c-p02-003-kq-main-lemma}{Keisler's Main Lemma}
\begin{proof}
By part~2 of \cref{thm:p02-w02c-p02-003-kq-chains}, all the axioms of
\(L^{*}(Q)\) hold in \((\mathfrak A^{*},q)\), so
\cref{thm:p02-w02c-p02-003-kq-finite-unions} applies there, with
\(L^{*}\) for \(L\).  A formula with one free variable can have that
variable renamed without changing the defined set, as in that proof; the
truth of \((Qy)\psi\) then does not change, by Axiom~3.

\emph{Part~1.}  Let \(L'=L^{*}\cup\{c\}\) with \(c\) new, and let
\(\Gamma\) consist of (a) all sentences of \(L^{*}(Q)\) true in
\((\mathfrak A^{*},q)\); (b) \(\varphi(c/x)\), written \(\varphi(c)\); (c)
\(\neg\psi(c)\) for every \(\psi(y)\) of \(L^{*}(Q)\) with one free variable
and \((\mathfrak A^{*},q)\models\neg(Qy)\psi(y)\).

(1) \emph{For \(\theta(y)\) of \(L^{*}(Q)\) with one free variable,
\(\theta(c)\) is consistent with \(\Gamma\) iff
\((\mathfrak A^{*},q)\models(Qu)(\theta(u)\land\varphi(u))\), with \(u\) not
occurring in \(\theta,\varphi\).}  If the right side fails,
\(\neg(\theta(c)\land\varphi(c))\) is of type (c), and with \(\varphi(c)\) it
gives \(\Gamma\vdash\neg\theta(c)\).  Conversely, suppose
\(\Gamma\vdash\neg\theta(c)\).  A finite part of \(\Gamma\) suffices; let
\(\Gamma_1\) be the set of type (a) and \(\neg\psi_1(c),\ldots,\neg\psi_n(c)\)
the sentences of type (c) used, with \(n>0\) (add \(\neg\psi(c)\) for
\(\psi(y)=\neg y=y\), which is of type (c) since \(\varnothing\notin q\)).
By part~1 of \cref{thm:p02-w02c-p02-003-kq-proof-lemmas} and (0a),
\(\Gamma_1\vdash(\varphi(c)\land\theta(c))\to\psi_1(c)\lor\cdots\lor\psi_n(c)\).
\(c\) does not occur in \(\Gamma_1\), so part~4 of that lemma, with
\(u\) also not occurring in the \(\psi_i\) (by Axiom~3 the choice of \(u\)
in (1) does not matter), gives
\(\Gamma_1\vdash(\varphi(u)\land\theta(u))\to\psi_1(u)\lor\cdots\lor\psi_n(u)\).
\((\mathfrak A^{*},q)\) is a weak model of \(\Gamma_1\) and of all the
axioms of \(L^{*}(Q)\), so by part~3 of
\cref{thm:p02-w02c-p02-003-kq-soundness} the set defined by
\(\varphi(u)\land\theta(u)\) is included in the union of the sets defined by
the \(\psi_i\), none of which is in \(q\).  By
\cref{thm:p02-w02c-p02-003-kq-finite-unions} it is not in \(q\): the right
side fails.

(2) \emph{\(\Gamma\) is consistent.}  The sets defined by \(\varphi(x)\)
and by \(u=u\land\varphi(u)\) are equal, and the first is in \(q\); so (1)
with \(\theta(y)=(y=y)\) shows that \(c=c\) is consistent with \(\Gamma\).

\(L^{*}\) is countable, since \(A\) and \(L\) are.  List the formulas
\(\psi(y)\) of \(L^{*}(Q)\) with one free variable and
\((\mathfrak A^{*},q)\models\neg(Qy)\psi\) as \(\psi_0,\psi_1,\ldots\)
(the list is not empty, by (b) of
\cref{thm:p02-w02c-p02-003-kq-finite-unions}), renaming the free variable of
\(\psi_n\) to \(x_n\), and let
\(\Sigma_n=\{\psi_n(x_n)\}\cup\{\neg x_n=c_a:a\in A\}\).  We verify the
hypothesis of part~1 of \cref{thm:p02-w02c-p02-003-kq-weak-completeness} for
\(\Gamma\) in \(L'(Q)\).  Let \(\theta\) be a formula of \(L'(Q)\) with only
\(x_n\) free and \((\exists x_n)\theta\) consistent with \(\Gamma\); write
\(\theta=\theta_0(c/u)\) with \(\theta_0\) in \(L^{*}(Q)\) and \(u\) new.
By (1),
\((\mathfrak A^{*},q)\models(Qu)(\varphi(u)\land(\exists x_n)\theta_0)\).  The
defined set is the union of the sets defined by
\(\varphi(u)\land(\exists x_n)(\theta_0\land\psi_n(x_n))\) and
\(\varphi(u)\land(\exists x_n)(\theta_0\land\neg\psi_n(x_n))\), so by
\cref{thm:p02-w02c-p02-003-kq-finite-unions} one of these is in \(q\).

If the second is, (1) shows that \((\exists
x_n)(\theta\land\neg\psi_n(x_n))\) is consistent with \(\Gamma\), and
\(\psi_n\in\Sigma_n\).

If the first is, \((Qu)(\exists x_n)\chi\) holds for
\(\chi=\varphi(u)\land\theta_0\land\psi_n(x_n)\).  By Axiom~4,
\((\exists x_n)(Qu)\chi\) or \((Qx_n)(\exists u)\chi\) holds.  The set
defined by \((\exists u)\chi\) is included in that defined by \(\psi_n\),
which is not in \(q\); so the second fails by
\cref{thm:p02-w02c-p02-003-kq-finite-unions}, and some \(a\in A\) has
\((\mathfrak A^{*},q)\models(Qu)(\varphi(u)\land\theta_0(c_a/x_n)\land\psi_n(c_a))\),
hence \((Qu)(\varphi(u)\land\theta_0(c_a/x_n))\) by the same lemma.  By (1),
\(\theta(c_a/x_n)\) is consistent with \(\Gamma\).  Since
\((\theta\land\neg\neg x_n=c_a)(c_a/x_n)\) follows from \(\theta(c_a/x_n)\)
by (0f) and (0a), (0c) shows that \((\exists x_n)(\theta\land\neg\neg
x_n=c_a)\) is consistent with \(\Gamma\), and \(\neg x_n=c_a\in\Sigma_n\).

So \(\Gamma\) has a countable weak model \((\mathfrak B',r)\) of \(L'\) in
which all the axioms of \(L'(Q)\) hold and which omits every \(\Sigma_n\).
For \(a\neq a'\), \(\neg c_a=c_{a'}\) is of type (a), so \(a\mapsto
c_a^{\mathfrak B'}\) is injective, and replacing \(\mathfrak B'\) by an
isomorphic copy (which satisfies the same formulas, by induction on
formulas) we may assume \(c_a^{\mathfrak B'}=a\).  Let \(\mathfrak B\) be the
reduct to \(L\); then \(A\subseteq B\), and \(\mathfrak B^{*}\) is the
reduct of \(\mathfrak B'\) to \(L^{*}\) together with constants for the
elements of \(B-A\).  For \(\varphi(\bar x)\) of \(L(Q)\) and \(\bar a\in
A\), \(\varphi(\bar c_{\bar a})\) or \(\neg\varphi(\bar c_{\bar a})\) is of
type (a), so \((\mathfrak A,q)\models\varphi[\bar a]\) iff
\((\mathfrak B,r)\models\varphi[\bar a]\): \((\mathfrak A,q)\prec
(\mathfrak B,r)\).  Let \(b=c^{\mathfrak B'}\).  By (b),
\((\mathfrak B^{*},r)\models\varphi[b]\).  For \(a\in A\), Axiom~1 holds in
\((\mathfrak A^{*},q)\), so \(\neg(c=c_a\lor c=c_a)\) is of type (c), and
\(b\neq a\): conclusion (i).  If \((\mathfrak A^{*},q)\models\neg(Qy)\psi\),
then \(\psi\) is some \(\psi_n\) up to the name of its free variable, and
since \((\mathfrak B',r)\) omits \(\Sigma_n\), every \(b'\in B\) satisfying
\(\psi_n\) equals some \(a\in A\): conclusion (ii)
\cite[\S2, Lemma~2.7, pp.~19--23]{Keisler1970}.

\emph{Part~2.}  If \((\mathfrak A^{*},q)\models\neg(Qx)(x=x)\), no
definable set is in \(q\) (by (a) of
\cref{thm:p02-w02c-p02-003-kq-finite-unions}), and
\((\mathfrak B,r)=(\mathfrak A,q)\) works, both sides of the equivalence
being false.  Otherwise list the formulas \(\varphi(x)\) of \(L^{*}(Q)\)
with \((\mathfrak A^{*},q)\models(Qx)\varphi\) as \(\varphi_0,\varphi_1,\ldots\)
(repeating if the list is finite).  Put \((\mathfrak A_0,q_0)=(\mathfrak
A,q)\).  Given a countable \((\mathfrak A_n,q_n)\succ(\mathfrak A,q)\), all
the axioms of \(L(Q)\) hold in it (their universal closures are sentences
true in \((\mathfrak A,q)\)), and by part~3 of
\cref{thm:p02-w02c-p02-003-kq-chains} \((\mathfrak A_n^{*},q_n)\models
(Qx)\varphi_n\).  Part~1 gives a countable \((\mathfrak A_{n+1},q_{n+1})\succ
(\mathfrak A_n,q_n)\) with some \(b_n\in A_{n+1}-A_n\) satisfying
\(\varphi_n\), and with (ii) for the language of \(\mathfrak A_n^{*}\).  Let
\((\mathfrak B,r)=\bigcup_n(\mathfrak A_n,q_n)\), countable and an
elementary extension of each \((\mathfrak A_n,q_n)\) by part~1 of
\cref{thm:p02-w02c-p02-003-kq-chains}.  If \((Qx)\varphi\) holds in
\((\mathfrak A^{*},q)\), then \(\varphi=\varphi_n\) for some \(n\), and
\(b_n\in B-A\) satisfies it in \((\mathfrak B^{*},r)\) by part~3 of
\cref{thm:p02-w02c-p02-003-kq-chains}.  If \(\neg(Qx)\varphi\) holds there,
it holds in each \((\mathfrak A_n^{*},q_n)\), so (ii) at stage \(n\) gives
\(\{a\in A_{n+1}:\varphi[a]\}\subseteq A_n\), and by induction on \(n\)
every element of \(A_n\) satisfying \(\varphi\) is in \(A\); every \(b\in B\)
lies in some \(A_n\), and satisfaction there agrees with satisfaction in
\((\mathfrak B^{*},r)\), so no \(b\in B-A\) satisfies \(\varphi\)
\cite[\S2, Lemma~2.8, pp.~23--24]{Keisler1970}.
\end{proof}
\end{annexitem}

\begin{annexitem}{annex:p02-w02c-p02-003-keisler-completeness}{thm:p02-w02c-p02-003-keisler-completeness}{Keisler's Completeness Theorem}
\begin{proof}
If \(\Sigma\) has a standard model, it is consistent by part~4 of
\cref{thm:p02-w02c-p02-003-kq-soundness}.  Let \(\Sigma\) be consistent.
By part~2 of \cref{thm:p02-w02c-p02-003-kq-weak-completeness} it has a
countable weak model \((\mathfrak A_0,q_0)\) in which all the axioms of
\(L(Q)\) hold.  By transfinite recursion (choosing at each successor step
one of the extensions given by part~2 of
\cref{thm:p02-w02c-p02-003-kq-main-lemma}) build an elementary chain
\((\mathfrak A_\alpha,q_\alpha)\), \(\alpha<\omega_1\), of countable weak
models in which all the axioms of \(L(Q)\) hold, such that: (1) for limit
\(\alpha\), \((\mathfrak A_\alpha,q_\alpha)=\bigcup_{\beta<\alpha}
(\mathfrak A_\beta,q_\beta)\); (2) for every \(\alpha\) and every formula
\(\varphi(x)\) of the language of \(\mathfrak A_\alpha^{*}\),
\((\mathfrak A_\alpha^{*},q_\alpha)\models(Qx)\varphi\) iff some \(a\in
A_{\alpha+1}-A_\alpha\) satisfies \(\varphi\) in
\((\mathfrak A_{\alpha+1},q_{\alpha+1})\) with each \(c_{a'}\),
\(a'\in A_\alpha\), read as \(a'\).  At a limit, part~1 of
\cref{thm:p02-w02c-p02-003-kq-chains} keeps the chain elementary; the union
is countable, as a countable union of countable sets (axiom of choice), and
the axioms hold in it, being true in \((\mathfrak A_0,q_0)\).

Let \((\mathfrak B,r)=\bigcup_{\alpha<\omega_1}(\mathfrak A_\alpha,q_\alpha)\).
By part~1 of \cref{thm:p02-w02c-p02-003-kq-chains},
\((\mathfrak A_0,q_0)\prec(\mathfrak B,r)\), so \((\mathfrak B,r)\) is a
weak model of \(\Sigma\).  It suffices to show (3): for every formula
\(\varphi(x_1,\ldots,x_n)\) of \(L(Q)\) and \(b_1,\ldots,b_n\in B\),
\((\mathfrak B,r)\models\varphi[\bar b]\) iff \(\mathfrak B\models\varphi[\bar
b]\).  By induction on \(\varphi\): atomic formulas and the steps for
\(\neg,\land,\exists,\forall\) are immediate, the clauses being the same.
Let \(\varphi=(Qx_0)\psi\) and choose \(\alpha<\omega_1\) with \(\bar b\in
A_\alpha\); let
\(S=\{b_0\in B:(\mathfrak B,r)\models\psi[b_0,\bar b]\}\), which by the
induction hypothesis is \(\{b_0:\mathfrak B\models\psi[b_0,\bar b]\}\).

Suppose \((\mathfrak B,r)\models(Qx_0)\psi[\bar b]\).  For every
\(\beta\) with \(\alpha\le\beta<\omega_1\), elementarity gives
\((\mathfrak A_\beta,q_\beta)\models(Qx_0)\psi[\bar b]\), so (2) gives some
\(a_\beta\in A_{\beta+1}-A_\beta\) with \(\psi[a_\beta,\bar b]\) in
\((\mathfrak A_{\beta+1},q_{\beta+1})\), hence in \((\mathfrak B,r)\).  For
\(\beta<\beta'\), \(a_\beta\in A_{\beta+1}\subseteq A_{\beta'}\) while
\(a_{\beta'}\notin A_{\beta'}\), so the \(a_\beta\) are distinct, and there
are \(\omega_1\) of them.  So \(S\) is uncountable and
\(\mathfrak B\models(Qx_0)\psi[\bar b]\).

Suppose \((\mathfrak B,r)\models\neg(Qx_0)\psi[\bar b]\).  Then for
\(\beta\ge\alpha\), \((\mathfrak A_\beta,q_\beta)\models\neg(Qx_0)\psi[\bar
b]\), and by (2) no \(a\in A_{\beta+1}-A_\beta\) satisfies \(\psi[a,\bar b]\)
there, nor, by elementarity, in \((\mathfrak B,r)\).  If \(b_0\in B-A_\alpha\),
let \(\beta'\) be least with \(b_0\in A_{\beta'}\); then \(\beta'>\alpha\),
\(\beta'\) is not a limit by (1), and \(\beta'=\beta+1\) with
\(\beta\ge\alpha\) and \(b_0\in A_{\beta+1}-A_\beta\).  Hence
\(S\subseteq A_\alpha\), which is countable, and
\(\mathfrak B\not\models(Qx_0)\psi[\bar b]\).  This completes (3), so
\(\mathfrak B\) is a standard model of \(\Sigma\)
\cite[\S2, pp.~24--26]{Keisler1970}.

\(B\) is a union of \(\omega_1\) countable sets, so has power at most
\(\omega_1\).  If \(\Sigma\cup\{(Qx)(x=x)\}\) is consistent, the model just
built for it is uncountable, so of power exactly \(\omega_1\).  If a
sentence \(\varphi\) is not provable, \(\{\neg\varphi\}\) is consistent (as
in part~3 of \cref{thm:p02-w02c-p02-003-kq-weak-completeness} with
\(\Sigma=\varnothing\)) and has a standard model in which \(\varphi\) fails;
the converse is part~4 of \cref{thm:p02-w02c-p02-003-kq-soundness}.
\end{proof}
\end{annexitem}

\begin{annexitem}{annex:p02-w02c-p02-003-keisler-compactness}{thm:p02-w02c-p02-003-keisler-compactness}{Keisler's Compactness Theorem}
\begin{proof}
If \(\Sigma\) were inconsistent, the two deductions of some \(\sigma\) and
\(\neg\sigma\) would use only a finite subset \(\Sigma_0\), which would be
inconsistent; but \(\Sigma_0\) has a standard model and is consistent by
part~4 of \cref{thm:p02-w02c-p02-003-kq-soundness}.  So \(\Sigma\) is
consistent and \cref{thm:p02-w02c-p02-003-keisler-completeness} gives a
standard model \cite[\S2, Corollary~2.11, p.~27]{Keisler1970}.
\end{proof}
\end{annexitem}

\begin{annexitem}{annex:p02-w02c-p02-003-los-vaught}{thm:p02-w02c-p02-003-los-vaught}{Categoricity test for completeness}
\begin{proof}
Suppose \(T\) is not complete: for some sentence \(\sigma\),
\(T\not\models\sigma\) and \(T\not\models\neg\sigma\).  Then
\(T\cup\{\neg\sigma\}\) and \(T\cup\{\sigma\}\) have models, and these are
models of \(T\), hence infinite.  By \cref{thm:upward-ls}, applied to
each of the two sets with the cardinal
\(\kappa\geq\max(\abs{L},\aleph_0)\), they have models
\(\mathfrak A\) and \(\mathfrak B\) of cardinality exactly \(\kappa\).
Both are models of \(T\), so by \(\kappa\)-categoricity there is an
isomorphism \(\mathfrak A\cong\mathfrak B\), and by
\cref{thm:p02-w02-p02-005-bw-iso-equiv} they satisfy the same sentences.
But \(\mathfrak A\models\neg\sigma\) and \(\mathfrak B\models\sigma\).  So
\(T\) is complete.
\end{proof}
\end{annexitem}

\begin{annexitem}{annex:p02-w02c-p02-003-simple-extension}{thm:p02-w02c-p02-003-simple-extension}{Simple algebraic extensions}
\begin{proof}
Two facts about \(K[x]\) are used.  \emph{Division:} for \(q,m\in K[x]\)
with \(m\) monic of degree \(d\geq1\) there are \(s,r\) with \(q=sm+r\)
and \(r=0\) or \(\deg r<d\); by induction on \(\deg q\), subtracting
\(c\,x^{\deg q-d}m\) with \(c\) the leading coefficient of \(q\) when
\(\deg q\geq d\).  \emph{Coprimality:} if \(m\) is irreducible (not a
product of two polynomials of smaller positive degree) and \(m\) does not
divide \(q\), there are \(u,v\in K[x]\) with \(um+vq=1\).  For the set
\(J=\{um+vq:u,v\in K[x]\}\) contains a nonzero polynomial \(e\) of least
degree, which may be taken monic.  Dividing \(m\) and \(q\) by \(e\), the
remainders lie in \(J\) and have smaller degree, so they are \(0\) and
\(e\) divides \(m\) and \(q\).  As \(m\) is irreducible, \(e\) is \(1\) or
\(e=c^{-1}m\) with \(c\) the leading coefficient of \(m\); the second
would make \(m\) divide \(q\).  So \(1=e\in J\).

\emph{Part~1.}  Among the nonzero \(q\in K[x]\) with \(q(a)=0\), take one
of least degree and divide by its leading coefficient: \(m\) is monic,
\(m(a)=0\), and \(\deg m\geq1\) since a nonzero constant has no root.  If
\(q(a)=0\), divide: \(q=sm+r\), so \(r(a)=0\) and \(\deg r<\deg m\), so
\(r=0\).  If \(m=gh\) with \(g,h\) of smaller positive degree, then
\(g(a)h(a)=0\), so \(g(a)=0\) or \(h(a)=0\) in the field \(E\),
contradicting minimality; so \(m\) is irreducible.

\emph{Part~2.}  \(K[a]\) contains \(0,1\) and is closed under \(+\),
\(-\) and \(\cdot\), since \(q\mapsto q(a)\) preserves them.  If
\(q(a)\neq0\), \(m\) does not divide \(q\), so \(um+vq=1\) for some
\(u,v\); at \(a\) this gives \(v(a)q(a)=1\), so \(q(a)\) has its inverse
\(v(a)\in K[a]\).

\emph{Part~3.}  For \(q\in K[x]\) let \(q^\sigma\in\sigma[K][x]\) apply
\(\sigma\) to the coefficients; \(q\mapsto q^\sigma\) preserves sums and
products.  \(m^\sigma\) is monic of degree \(\geq1\), so it has a root
\(b\in G\).  Put \(\tau(q(a))=q^\sigma(b)\).  This is well defined: if
\(q(a)=q'(a)\), then \(m\) divides \(q-q'\) by part~1, so \(m^\sigma\)
divides \(q^\sigma-q'^\sigma\), and \(q^\sigma(b)=q'^\sigma(b)\).
\(\tau\) preserves \(+\) and \(\cdot\) and extends \(\sigma\) (constant
polynomials).  It is injective: if \(q(a)\neq0\), then \(um+vq=1\) as in
part~2, so \(u^\sigma(b)m^\sigma(b)+v^\sigma(b)q^\sigma(b)=1\), and since
\(m^\sigma(b)=0\), \(q^\sigma(b)\neq0\).

\emph{Part~4.}  Let \(b\) lie in a field extending \(E\) and be algebraic
over \(E\), and let \(m\) be as in part~1 for \(E\) and \(b\).  Since \(E\)
is algebraically closed, \(m\) has a root \(c\in E\); dividing by the
monic \(x-c\), \(m=s\cdot(x-c)+r\) with \(r\) constant and \(r=m(c)=0\).
As \(m\) is irreducible, \(s\) has degree \(0\); being monic, \(s=1\), so
\(m=x-c\), and \(m(b)=0\) gives \(b=c\in E\).
\end{proof}
\end{annexitem}

\begin{annexitem}{annex:p02-w02c-p02-003-steinitz}{thm:p02-w02c-p02-003-steinitz}{Steinitz's classification}
\begin{proof}
\emph{Prime fields.}  In a field of prime characteristic \(p\), the
elements \(k\cdot1\) (\(k<p\)) form a subfield isomorphic to the integers
modulo \(p\): they are closed under \(+\) and \(\cdot\), which act
modulo \(p\); \(k\cdot1\neq0\) for \(0<k<p\) by the definition of the
characteristic, and a finite subring without zero divisors (\(p\) being
prime, \(k\cdot1\cdot l\cdot1=(kl)\cdot1\neq0\) for \(0<k,l<p\)) has
inverses, since multiplication by a nonzero element is injective, hence
onto.  In characteristic \(0\), \(k\cdot1\neq0\) for \(k\geq1\), and the
elements \((k\cdot1)(l\cdot1)^{-1}\) and their negatives form a subfield
isomorphic to \(\mathbb Q\).  These subfields are contained in every
subfield, so they are the prime fields, and \(k\cdot1\mapsto k\cdot1\)
(with quotients) is an isomorphism between the prime fields \(P\) and
\(P'\) of two fields of the same characteristic.  \(P\) is countable.

\emph{Elements of \(P(X)\).}  If \(X\subseteq E\), the set of quotients
\(q(\bar x)/r(\bar x)\), with \(q,r\in P[t_1,\ldots,t_k]\), \(\bar x\) from
\(X\) and \(r(\bar x)\neq0\), contains \(P\cup X\), is closed under the
field operations (common denominators) and lies in every subfield
containing \(P\cup X\); so it is \(P(X)\).

\emph{Part~1.}  Let \(\mathcal X\) be the family of subsets of \(E\)
algebraically independent over \(P\), ordered by inclusion.  It contains
\(\emptyset\).  The union of a chain in \(\mathcal X\) lies in
\(\mathcal X\), since a failure of independence involves finitely many
elements, which all lie in one member of the chain.  By Zorn's lemma
there is a maximal \(X\in\mathcal X\).  Let \(a\in E\setminus X\).  By
maximality there are distinct \(x_1,\ldots,x_k\in X\) and a nonzero \(q\)
with \(q(x_1,\ldots,x_k,a)=0\).  Write
\(q=\sum_{j\leq d}q_j(t_1,\ldots,t_k)\,t_{k+1}^j\).  Some \(q_j\) is
nonzero, so \(q_j(\bar x)\neq0\) by independence of \(X\), and \(a\) is a
root of the nonzero polynomial \(\sum_jq_j(\bar x)y^j\) over \(P(X)\).
Elements of \(X\) are algebraic over \(P(X)\) trivially.

\emph{Part~2.}  \(\abs X\leq\abs E\).  Conversely, each element of
\(P(X)\) is given by a pair of polynomials over \(P\) and a finite
sequence from \(X\), and there are at most \(\max(\abs X,\aleph_0)\) such
data (finite sequences from a set of that cardinality, which is infinite,
\(P\) being countable).  So \(\abs{P(X)}\leq\max(\abs X,\aleph_0)\), and
for the same reason there are at most \(\max(\abs X,\aleph_0)\) nonzero
polynomials over \(P(X)\).  Each has at most finitely many roots in
\(E\): a nonzero polynomial of degree \(d\) has at most \(d\) roots, by
induction on \(d\), since a root \(c\) gives \(q=s\cdot(x-c)\) by
division, and every other root is a root of \(s\).  By part~1 every
element of \(E\) is such a root, so
\(\abs E\leq\max(\abs X,\aleph_0)\).  If \(E\) is uncountable this gives
\(\abs E\leq\abs X\), so \(\abs X=\abs E\).

\emph{Part~3.}  Let \(X\), \(Y\) be transcendence bases of \(E\), \(G\)
over their prime fields \(P\), \(P'\), and \(\beta\colon X\to Y\) a
bijection.  Let \(\pi\colon P\to P'\) be the isomorphism of prime fields.
Define \(\sigma_0\colon P(X)\to G\) by
\(\sigma_0(q(\bar x)/r(\bar x))=q^\pi(\beta\bar x)/r^\pi(\beta\bar x)\).
Since \(r\neq0\) as a polynomial and \(Y\) is independent,
\(r^\pi(\beta\bar x)\neq0\).  The value does not depend on the
representation, and \(\sigma_0\) is injective: \(q(\bar x)/r(\bar
x)=q'(\bar x)/r'(\bar x)\) iff \(qr'-q'r\) vanishes at \(\bar x\) iff (by
independence of \(X\)) \(qr'=q'r\) as polynomials iff
\(q^\pi r'^\pi=q'^\pi r^\pi\) iff (by independence of \(Y\)) the values at
\(\beta\bar x\) agree.  \(\sigma_0\) preserves the field operations, as
they are computed on representatives.

Let \(\mathcal S\) be the set of pairs \((K,\sigma)\) with \(K\) a subfield
of \(E\) containing \(P(X)\) and \(\sigma\colon K\to G\) an embedding
extending \(\sigma_0\), ordered by extension.  A chain has as upper bound
the union of its fields with the union of its maps.  By Zorn's lemma there
is a maximal \((K,\sigma)\).  If \(a\in E\setminus K\), then \(a\) is
algebraic over \(P(X)\subseteq K\) by part~1, hence over \(K\), and by
part~3 of \cref{thm:p02-w02c-p02-003-simple-extension} \(\sigma\) extends
to the field \(K[a]\supsetneq K\), contrary to maximality.  So \(K=E\) and
\(\sigma\colon E\to G\) is an embedding.  Its image \(\sigma[E]\) is a
subfield of \(G\) isomorphic to \(E\), hence algebraically closed, and it
contains \(\sigma_0[P(X)]=P'(Y)\).  Every \(b\in G\) is algebraic over
\(P'(Y)\) by part~1 for \(G\), so over \(\sigma[E]\), so \(b\in\sigma[E]\)
by part~4 of \cref{thm:p02-w02c-p02-003-simple-extension}.  So \(\sigma\)
is onto and \(E\cong G\).  If \(\abs E=\abs G=\kappa\) is uncountable, any
transcendence bases have cardinality \(\kappa\) by part~2, and the first
claim applies.  The models of \(\mathrm{ACF}_p\) are exactly the
algebraically closed fields of characteristic \(p\), so
\(\mathrm{ACF}_p\) is \(\kappa\)-categorical.
\end{proof}
\end{annexitem}

\begin{annexitem}{annex:p02-w02c-p02-003-fta}{thm:p02-w02c-p02-003-fta}{Fundamental theorem of algebra}
\begin{proof}
Let \(f(z)=z^n+a_{n-1}z^{n-1}+\cdots+a_0\) with \(n\geq1\) and
\(a_j\in\mathbb C\).  Put \(S=\abs{a_0}+\cdots+\abs{a_{n-1}}\) and
\(R=1+2S\).

\emph{Step~1: a global minimum of \(\abs f\).}  For \(\abs z\geq R\),
\(\abs z\geq1\), so \(\abs z^j\leq\abs z^{n-1}\) for \(j\leq n-1\) and
\[
 \abs{f(z)}\geq\abs z^n-\sum_{j<n}\abs{a_j}\abs z^j
 \geq\abs z^{n-1}(\abs z-S)\geq\abs z-S\geq1+S>\abs{a_0}=\abs{f(0)}.
\]
The disc \(D=\{z:\abs z\leq R\}\) is closed and bounded in
\(\mathbb C=\mathbb R^2\), and \(z\mapsto\abs{f(z)}\) is continuous, being
built from sums, products and the modulus.  By the completeness of
\(\mathbb R\), a continuous real function on a closed bounded subset of
\(\mathbb R^2\) attains its minimum (every sequence in \(D\) has a
convergent subsequence with limit in \(D\), by bisection on each
coordinate; take a sequence along which the values tend to the infimum).
Let \(z_0\in D\) be a minimum point on \(D\).  Since \(0\in D\),
\(\abs{f(z_0)}\leq\abs{f(0)}<\abs{f(z)}\) for \(\abs z>R\), so \(z_0\) is
a minimum point on \(\mathbb C\).

\emph{Step~2: the minimum is \(0\).}  Suppose \(f(z_0)\neq0\).  Expanding
\(f(z_0+w)\) in powers of \(w\) and dividing by \(f(z_0)\) gives
\(g(w)=f(z_0+w)/f(z_0)=1+\sum_{j=1}^nb_jw^j\), with \(b_n=1/f(z_0)\neq0\).
Let \(k\) be the least \(j\geq1\) with \(b_j\neq0\).  Write
\(-1/b_k=r(\cos\theta+i\sin\theta)\) with \(r>0\), and put
\(\omega=r^{1/k}(\cos(\theta/k)+i\sin(\theta/k))\); by the addition
formulas for cosine and sine, \((\cos\alpha+i\sin\alpha)^k=\cos
k\alpha+i\sin k\alpha\), so \(\omega^k=-1/b_k\).  For real \(t\) with
\(0<t\leq1\),
\[
 g(t\omega)=1-t^k+t^{k+1}\sum_{j=k+1}^nb_j\omega^jt^{j-k-1},
\]
so that
\[
 \abs{g(t\omega)}\leq1-t^k+Ct^{k+1}=1-t^k(1-Ct),
\]
with \(C=\sum_{j>k}\abs{b_j}\abs\omega^j\).  For \(0<t<\min(1,1/C)\)
(any \(0<t<1\) if \(C=0\)), \(\abs{g(t\omega)}<1\), that is,
\(\abs{f(z_0+t\omega)}<\abs{f(z_0)}\), contrary to Step~1.  So
\(f(z_0)=0\).

\(\mathbb C\) is a field in which \(k\cdot1=k\neq0\) for \(k\geq1\), so it
satisfies the field axioms and the characteristic axioms of
\(\mathrm{ACF}_0\); by Steps~1 and~2 it satisfies the axiom for each
degree \(n\).
\end{proof}
\end{annexitem}

\begin{annexitem}{annex:p02-w02c-p02-003-acf}{thm:p02-w02c-p02-003-acf}{Algebraically closed fields}
\begin{proof}
\emph{Part~1.}  \(\mathrm{ACF}_p\) has no finite model.  Let \(F\) be a
finite field with elements \(a_1,\ldots,a_n\); \(n\geq2\) since
\(0\neq1\).  The polynomial
\(q(x)=(x-a_1)\cdots(x-a_n)+1\) is monic of degree \(n\), and expanding
the product gives coefficients \(y_0,\ldots,y_{n-1}\in F\) with
\(q(x)=x^n+y_{n-1}x^{n-1}+\cdots+y_0\).  For each \(i\),
\(q(a_i)=0+1=1\neq0\), so \(q\) has no root in \(F\), and the axiom of
\(\mathrm{ACF}_p\) for degree \(n\) fails in \(F\).

\(\mathrm{ACF}_p\) is \(\aleph_1\)-categorical: its models in Barwise's
language \(\{+,\cdot,0,1\}\) are the algebraically closed fields of
characteristic \(p\), and any two of cardinality \(\aleph_1\) are
isomorphic by part~3 of \cref{thm:p02-w02c-p02-003-steinitz}.  The language is finite,
so \(\aleph_1\geq\max(\abs{L},\aleph_0)\), and
\cref{thm:p02-w02c-p02-003-los-vaught} gives completeness.  Any two
algebraically closed fields \(F_1,F_2\) of characteristic \(p\) are models
of \(\mathrm{ACF}_p\); for each \(\sigma\), \(\mathrm{ACF}_p\models\sigma\)
or \(\mathrm{ACF}_p\models\neg\sigma\), and both fields follow the same
alternative, so \(F_1\equiv F_2\).

\emph{Part~2.}  Whether a sentence belongs to \(\mathrm{ACF}_0\) can be
decided by inspecting its form: it is one of the finitely many field
axioms, an axiom of the displayed form for some \(n\), or
\(\neg(1+\cdots+1=0)\) with some number of summands.  By the strong
completeness theorem \cref{thm:fol-completeness},
\(\mathrm{ACF}_0\models\sigma\) iff \(\sigma\) is derivable from
\(\mathrm{ACF}_0\) in the natural deduction calculus, and derivations are
finite objects that can be listed one after another and checked.  Given
\(\sigma\), run through all derivations and stop at the first whose
premises lie in \(\mathrm{ACF}_0\) and whose conclusion is \(\sigma\) or
\(\neg\sigma\); answer yes in the first case and no in the second.  By
part~1 one of \(\sigma\), \(\neg\sigma\) follows from \(\mathrm{ACF}_0\),
so by completeness the search stops.  If the answer is yes, then
\(\mathrm{ACF}_0\models\sigma\) by soundness, \cref{thm:fol-soundness}.  If it is no, then
\(\mathrm{ACF}_0\models\neg\sigma\); an algebraically closed field of
characteristic \(0\) exists by \cref{thm:p02-w02c-p02-003-fta} and
satisfies \(\neg\sigma\), so
\(\mathrm{ACF}_0\not\models\sigma\).

\emph{Part~3.}  By \cref{thm:p02-w02c-p02-003-fta}, \(\mathbb C\) is a
model of \(\mathrm{ACF}_0\).  If \(\mathrm{ACF}_0\models\sigma\) then
\(\mathbb C\models\sigma\).  If \(\mathrm{ACF}_0\not\models\sigma\) then
\(\mathrm{ACF}_0\models\neg\sigma\) by part~1, so
\(\mathbb C\models\neg\sigma\).
\end{proof}
\end{annexitem}

\begin{annexitem}{annex:p02-w02c-p02-003-lagrange}{thm:p02-w02c-p02-003-lagrange}{Lagrange's four-square theorem}
\begin{proof}
\emph{Euler's identity.}  For integers \(a_1,\ldots,a_4,b_1,\ldots,b_4\),
\[
 (a_1^2+a_2^2+a_3^2+a_4^2)(b_1^2+b_2^2+b_3^2+b_4^2)=c_1^2+c_2^2+c_3^2+c_4^2,
\]
where \(c_1=a_1b_1+a_2b_2+a_3b_3+a_4b_4\),
\(c_2=a_1b_2-a_2b_1+a_3b_4-a_4b_3\),
\(c_3=a_1b_3-a_2b_4-a_3b_1+a_4b_2\) and
\(c_4=a_1b_4+a_2b_3-a_3b_2-a_4b_1\).  Expanding the right side, each
square \(a_j^2b_l^2\) occurs exactly once, and each cross term
\(2a_ja_{j'}b_lb_{l'}\) occurs twice with opposite signs.  So a product of
two sums of four squares is one.

\emph{Reduction to primes.}  \(0=0^2+0^2+0^2+0^2\),
\(1=1^2+0^2+0^2+0^2\), and every \(n\geq2\) is a product of primes, so by
Euler's identity it suffices to treat a prime \(p\).  For \(p=2\),
\(2=1^2+1^2+0^2+0^2\).  Let \(p\) be odd.

\emph{A multiple of \(p\).}  The \((p+1)/2\) numbers \(x^2\),
\(0\leq x\leq(p-1)/2\), are pairwise incongruent modulo \(p\): if
\(x^2\equiv x'^2\) with \(x\neq x'\), then \(p\) divides
\((x-x')(x+x')\), but \(0<\abs{x-x'}<p\) and \(0<x+x'<p\).  Likewise the
\((p+1)/2\) numbers \(-1-y^2\), \(0\leq y\leq(p-1)/2\), are pairwise
incongruent.  There are only \(p\) residues, so some \(x^2\equiv-1-y^2\),
and \(x^2+y^2+1^2+0^2=mp\) with \(m\geq1\) and
\(mp\leq2\bigl(\frac{p-1}2\bigr)^2+1<p^2\), so \(m<p\).

\emph{Descent.}  Let \(m_0\) be the least \(m\geq1\) such that \(mp\) is a
sum of four squares; \(m_0<p\).  Suppose \(m_0>1\) and
\(m_0p=x_1^2+x_2^2+x_3^2+x_4^2\).  If \(m_0\) is even, the number of odd
\(x_j\) is even; renumber so that \(x_1\equiv x_2\) and \(x_3\equiv x_4\)
modulo \(2\).  Then
\[
 \tfrac{m_0}2p=\bigl(\tfrac{x_1+x_2}2\bigr)^2+\bigl(\tfrac{x_1-x_2}2\bigr)^2
 +\bigl(\tfrac{x_3+x_4}2\bigr)^2+\bigl(\tfrac{x_3-x_4}2\bigr)^2
\]
with integer terms, against minimality.  If \(m_0\) is odd, choose
\(y_j\equiv x_j\pmod{m_0}\) with \(\abs{y_j}\leq(m_0-1)/2\).  Then
\(\sum y_j^2\equiv\sum x_j^2\equiv0\pmod{m_0}\), so \(\sum y_j^2=rm_0\),
and \(rm_0\leq4\bigl(\frac{m_0-1}2\bigr)^2<m_0^2\), so \(r<m_0\).  If
\(r=0\), every \(y_j=0\), so \(m_0\) divides every \(x_j\), \(m_0^2\)
divides \(m_0p\), and \(m_0\) divides \(p\), impossible for
\(1<m_0<p\).  So \(1\leq r<m_0\).  By Euler's identity with \(a=x\),
\(b=y\), \((m_0p)(rm_0)=c_1^2+\cdots+c_4^2\).  Modulo \(m_0\),
\(c_1\equiv\sum x_j^2\equiv0\), and substituting \(x_j\) for \(y_j\) makes
\(c_2,c_3,c_4\) vanish identically, so they are \(\equiv0\) too.  Hence
\(rp=\sum(c_j/m_0)^2\) is a sum of four squares with \(r<m_0\), against
minimality.  So \(m_0=1\).

\emph{The formula \(D\).}  In \(\mathbb Z\), a value of the right side of
\(D(x_i)\) is a sum of squares, hence non-negative; and every
non-negative integer is such a value by what was just proved.  So
\(\mathbb Z\models D(x_i)[s]\) iff \(s(x_i)\geq0\).
\end{proof}
\end{annexitem}

\begin{annexitem}{annex:p02-w02c-p02-003-ring-coding}{thm:p02-w02c-p02-003-ring-coding}{The translation is arithmetically definable}
\begin{proof}
Write \(N_k(\ldots)\) for the constructor codes of
\cref{sec:arithmetic-coding-04}.  The function \(d(i)=\ulcorner
D(x_i)^{\circ}\urcorner\) is primitive recursive, being a fixed
composition of the constructors \(N_0\), \(N_3\), \(N_4\), \(N_5\),
\(N_{12}\) applied to \(i,i+1,\ldots,i+4\).  Define \(\mathrm{tr}(c)\) by
course-of-values recursion on \(c\), reading the tag and arguments of
\(c\):
\begin{itemize}
\item \(N_0(i)\mapsto N_0(i)\); \(N_1()\mapsto N_1()\);
      \(N_2(r)\mapsto N_3(\mathrm{tr}(r),N_2(N_1()))\);
      \(N_k(r,t)\mapsto N_k(\mathrm{tr}(r),\mathrm{tr}(t))\) for
      \(k=3,4,5,8,9,10\);
\item \(N_6()\mapsto N_6()\); \(N_7(\varphi)\mapsto
      N_7(\mathrm{tr}(\varphi))\);
\item \(N_{11}(i,\varphi)\mapsto N_{11}(i,N_{10}(d(i),\mathrm{tr}(\varphi)))\);
      \(N_{12}(i,\varphi)\mapsto N_{12}(i,N_8(d(i),\mathrm{tr}(\varphi)))\);
\item every other \(c\) to \(0\).
\end{itemize}
Each clause uses \(\mathrm{tr}\) only at proper child codes, which are
smaller than \(c\), and otherwise primitive recursive operations (tag and
argument extraction, constructors, \(d\)).  By the course-of-values
principle of \cref{sec:arithmetic-coding-04}, \(\mathrm{tr}\) is
primitive recursive.  By induction on terms and formulas,
\(\mathrm{tr}(\ulcorner t\urcorner)=\ulcorner(t^{\mathbb Z})^{\circ}\urcorner\)
and \(\mathrm{tr}(\ulcorner\chi\urcorner)=\ulcorner(\chi^{\mathbb
Z})^{\circ}\urcorner\): the clause for \(S(r)\) produces the code of
\(+(r',S(0))\), which is \(((S\,r)^{\mathbb Z})^{\circ}\) when \(r'\) is
the code of \((r^{\mathbb Z})^{\circ}\), since \((t^{\mathbb
Z}+1)^{\circ}\) is \(+((t^{\mathbb Z})^{\circ},S(0))\); the quantifier
clauses produce the codes of \(\forall x_i\,(D(x_i)^{\circ}\to
(\chi^{\mathbb Z})^{\circ})\) and \(\exists x_i\,(D(x_i)^{\circ}\land
(\chi^{\mathbb Z})^{\circ})\); the remaining clauses copy the constructor.
So \(\mathrm{tr}(\ulcorner\chi\urcorner)=\ulcorner\chi^{\mathbb
Z}\urcorner'\).

The function \(\operatorname{Snt}\) (\(1\) on the codes of
\(\Lang_A\)-sentences, \(0\) elsewhere) is primitive recursive, since a
sentence is a formula with empty free-variable list and both tests are
primitive recursive by \cref{sec:arithmetic-coding-04}.  So the function
\(e(u,v)=\operatorname{Snt}(u)\cdot\chi_=(v,\mathrm{tr}(u))\), with
\(\chi_=\) the characteristic function of equality, is primitive
recursive, and \(e(u,v)=1\) iff \(u\) codes a sentence \(\sigma\) and
\(v=\ulcorner\sigma^{\mathbb Z}\urcorner'\).  By
\cref{thm:p02-w02b-p02-034-pr-graph} there is a graph formula
\(\Gamma(x_1,x_2,y)\) with \(\mathcal N\models\Gamma(\bar a,\bar b,\bar
c)\) iff \(e(a,b)=c\).  The formula \(\Gamma(x_1,x_2,S(0))\) defines the
relation.
\end{proof}
\end{annexitem}

\begin{annexitem}{annex:p02-w02c-p02-003-integers}{thm:p02-w02c-p02-003-integers}{The ring of integers is not decidable}
\begin{proof}
\emph{Part~1.}  Let \(\mathbb N\) here denote the set of non-negative
integers inside \(\mathbb Z\).  By \cref{thm:p02-w02c-p02-003-lagrange},
\(\mathbb Z\models D(x_i)[s]\) iff \(s(x_i)\in\mathbb N\); this depends
only on \(s(x_i)\), the bound variables of \(D(x_i)\) being distinct from
\(x_i\).  For an
\(\Lang_A\)-term \(t\) let \(t^{\mathbb Z}\) be the term of
\(\{+,\cdot,0,1\}\) with \(0\mapsto0\), \(S\,t\mapsto t^{\mathbb Z}+1\),
\(t+t'\mapsto t^{\mathbb Z}+t'^{\mathbb Z}\),
\(t\times t'\mapsto t^{\mathbb Z}\cdot t'^{\mathbb Z}\).  For an assignment
\(s\) into \(\mathbb N\), the value of \(t^{\mathbb Z}\) in \(\mathbb Z\) is
the value of \(t\) in \(\mathcal N\), by induction on \(t\), since \(0\),
successor, addition and multiplication of non-negative integers are the
operations of \(\mathcal N\).  For formulas, \(\chi^{\mathbb Z}\) replaces
terms in atomic formulas as above, commutes with the connectives, and sends
\(\exists x\,\chi\) to \(\exists x\,(D(x)\land\chi^{\mathbb Z})\) and
\(\forall x\,\chi\) to \(\forall x\,(D(x)\to\chi^{\mathbb Z})\), as in
\cref{def:p02-w02c-p02-003-ring-coding}.  By induction on
\(\chi\), for every assignment \(s\) into \(\mathbb N\),
\(\mathcal N\models\chi[s]\) iff \(\mathbb Z\models\chi^{\mathbb Z}[s]\):
atomic formulas by the term clause; connectives directly; and
\(\mathbb Z\models\exists x\,(D(x)\land\chi^{\mathbb Z})[s]\) iff some
\(a\in\mathbb N\) has \(\mathbb Z\models\chi^{\mathbb Z}[s(a/x)]\) iff (by
the induction hypothesis, \(s(a/x)\) being into \(\mathbb N\)) some
\(a\in\mathbb N\) has \(\mathcal N\models\chi[s(a/x)]\); the universal case
is dual.  For a sentence \(\sigma\) this is the claim.

\emph{Part~2.}  Let \(\rho(u,v)\) define the relation of the hypothesis in
\(\mathcal N\), and suppose some formula \(\omega(v)\) defined \(W\) in
\(\mathcal N\).  Let \(\tau(u)\) be \(\exists v\,(\rho(u,v)\land\omega(v))\).
For every first-order \(\Lang_A\)-sentence \(\sigma\), with
\(c=\ulcorner\sigma\urcorner\):
\(\mathcal N\models\tau(\bar c)\) iff \(\ulcorner\sigma^{\mathbb Z}\urcorner'\in W\)
iff \(\mathbb Z\models\sigma^{\mathbb Z}\) iff \(\mathcal N\models\sigma\)
by part~1.  So \(\tau\) is an arithmetical truth definition, which
\cref{thm:tarski-undefinability} excludes.  Hence \(W\) is not definable
in \(\mathcal N\).  If \(W\) were computably enumerable, it would be
definable by \cref{thm:p02-w02b-p02-034-ce-sigma-one}; so it is not.  A
decision procedure for the theory of \(\mathbb Z\) would list the members
of \(W\) by testing the codes \(0,1,2,\ldots\) in turn, so in Barwise's
sense there is none.

\emph{Part~3.}  Let \(T\) be true in \(\mathbb Z\), complete, and with
consequence set \(C=\{\ulcorner\varphi\urcorner':T\models\varphi\}\).  If
\(T\models\varphi\) then \(\mathbb Z\models\varphi\), since \(\mathbb Z\)
is a model of \(T\).  If \(\mathbb Z\models\varphi\), then
\(T\models\varphi\) or \(T\models\neg\varphi\) by completeness, and the
second would give \(\mathbb Z\models\neg\varphi\); so \(T\models\varphi\).
Thus \(C=W\), which is not computably enumerable by part~2.
\end{proof}
\end{annexitem}

\begin{annexitem}{annex:p02-w02c-p02-003-ring-transfer}{thm:p02-w02c-p02-003-ring-transfer}{Consequence in the ring language and in \(\Lang_A\)}
\begin{proof}
For a structure \(\mathfrak M\) of \(\{+,\cdot,0,1\}\), let
\(\mathfrak M^{*}\) be the \(\Lang_A\)-structure with the universe of
\(\mathfrak M\), \(0^{\mathfrak M^{*}}=0^{\mathfrak M}\),
\(+^{\mathfrak M^{*}}=+^{\mathfrak M}\),
\(\times^{\mathfrak M^{*}}=\cdot^{\mathfrak M}\), and \(S^{\mathfrak M^{*}}\)
the constant function with value \(1^{\mathfrak M}\).  For an
\(\Lang_A\)-structure \(\mathfrak K\), let \(\mathfrak K'\) be the structure
of \(\{+,\cdot,0,1\}\) with the universe of \(\mathfrak K\),
\(0^{\mathfrak K'}=0^{\mathfrak K}\), \(+^{\mathfrak K'}=+^{\mathfrak K}\),
\(\cdot^{\mathfrak K'}=\times^{\mathfrak K}\) and
\(1^{\mathfrak K'}=S^{\mathfrak K}(0^{\mathfrak K})\).

\emph{Terms.}  For every term \(t\) of \(\{+,\cdot,0,1\}\) and every
assignment \(s\),
\((t^{\circ})^{\mathfrak M^{*}}[s]=t^{\mathfrak M}[s]\) and
\((t^{\circ})^{\mathfrak K}[s]=t^{\mathfrak K'}[s]\).  By induction on
\(t\): variables and \(0\) are unchanged by \({}^{\circ}\) and have the
same values; \(1^{\circ}=S(0)\) has value
\(S^{\mathfrak M^{*}}(0^{\mathfrak M})=1^{\mathfrak M}\) in
\(\mathfrak M^{*}\) and \(S^{\mathfrak K}(0^{\mathfrak K})=1^{\mathfrak K'}\)
in \(\mathfrak K\); \((t+t')^{\circ}=t^{\circ}+t'^{\circ}\) and
\((t\cdot t')^{\circ}=t^{\circ}\times t'^{\circ}\), whose values agree by
the induction hypothesis and the choice of the operations.

\emph{Formulas.}  The map \({}^{\circ}\) changes only terms and commutes
with \(\bottom\), the connectives and the quantifiers.  By induction on a
formula \(\varphi\) of \(\{+,\cdot,0,1\}\),
\(\mathfrak M\models\varphi[s]\) iff \(\mathfrak M^{*}\models\varphi^{\circ}[s]\),
and \(\mathfrak K'\models\varphi[s]\) iff
\(\mathfrak K\models\varphi^{\circ}[s]\): atomic formulas by the term case,
the connectives by the induction hypothesis, and the quantifiers because
each pair of structures has the same universe, so the same assignments.

\emph{Consequence.}  Write \(\Sigma^{\circ}=\{\psi^{\circ}:\psi\in\Sigma\}\).
Suppose \(\Sigma\models\varphi\) and \(\mathfrak K\models\Sigma^{\circ}\).
Then \(\mathfrak K'\models\Sigma\), so \(\mathfrak K'\models\varphi\), so
\(\mathfrak K\models\varphi^{\circ}\).  Suppose
\(\Sigma^{\circ}\models\varphi^{\circ}\) and \(\mathfrak M\models\Sigma\).
Then \(\mathfrak M^{*}\models\Sigma^{\circ}\), so
\(\mathfrak M^{*}\models\varphi^{\circ}\), so \(\mathfrak M\models\varphi\).
\end{proof}
\end{annexitem}

\begin{annexitem}{annex:p02-w02c-p02-003-sharp-enumeration}{thm:p02-w02c-p02-003-sharp-enumeration}{An explicit enumeration of \(T^{\sharp}\)}
\begin{proof}
\emph{Part~1.}  \(V\) is primitive recursive by the paragraph on certified
axiom leaves in \cref{sec:arithmetic-coding-04}; \(\mathrm{tr}\) is
primitive recursive by \cref{thm:p02-w02c-p02-003-ring-coding};
\(\mathsf{Ax}^{\rm d}_{\PA}\) is primitive recursive by
\cref{thm:p05-direct-pa-axioms}; equality, sign, addition and
multiplication are primitive recursive.  So \(\mathsf{Lf}\) is primitive
recursive, by composition.  The proof of \cref{thm:arithmetization-syntax}
shows that the derivation checks, the open lists, the root and sentence
tests and the length test on \(W\) are primitive recursive, and that a loop
of length \(\ell(W)\) applying a primitive recursive test to each pair
\((a_{j_v},s_v)\) keeps the relation primitive recursive; there the test is
\(V(e_T,\cdot,\cdot)\), here it is \(\mathsf{Lf}\), and nothing else in that
argument depends on the test.  So \(\mathsf{Pf}^{+}\) is primitive
recursive.  Sentence recognition is primitive recursive
(\cref{thm:arithmetization-syntax}), so \(\mathsf{En}^{\sharp}\), a product
of the sentence test at \(x\) and the characteristic function of
\(\mathsf{Pf}^{+}\) at \((p,\mathrm{tr}(x))\), is primitive recursive.

\emph{Part~2.}  The members of \((T^{+})^{\circ}\) are the sentences
\(\tau^{\circ}\) with \(\tau\in T\), whose codes form \(C_T\), and the
sentences \((\alpha^{\mathbb Z})^{\circ}\) with \(\alpha\in\PA^{\rm d}\),
whose codes are \(\ulcorner\alpha^{\mathbb Z}\urcorner'=
\mathrm{tr}(\ulcorner\alpha\urcorner)\) by
\cref{thm:p02-w02c-p02-003-ring-coding}.

Suppose \(\mathsf{Pf}^{+}(p,\ulcorner\psi\urcorner)\) with
\(p=N_{16}(N_{15}(R),W)\).  The row induction in the proof of
\cref{thm:arithmetization-syntax} reconstructs from \(R\) an
\(\mathsf{ND}_{\mathsf{FOL}}\) derivation of \(\psi\) whose open assumption occurrences are
exactly the leaves in the root open list.  Let \(a=a_{j_v}\) be the code of
such a leaf, with \(\mathsf{Lf}(a,s_v)=1\).  If \(V(e_T,a,s_v)=1\), then
\(\mathcal N\models\mathsf{Comp}(\bar e_T,\bar a,\bar s_v)\), so the partial
function of \(e_T\) is defined at \(a\) by \cref{thm:bounded-verification},
and \(a\in C_T\); the leaf is \(\tau^{\circ}\) for some \(\tau\in T\), since
codes determine formulas.  Otherwise \(a=\mathrm{tr}(s_v)\) and
\(\mathsf{Ax}^{\rm d}_{\PA}(s_v)=1\); by \cref{thm:p05-direct-pa-axioms},
\(s_v=\ulcorner\alpha\urcorner\) for some \(\alpha\in\PA^{\rm d}\), and the
leaf is \((\alpha^{\mathbb Z})^{\circ}\).  Hence every open leaf lies in
\((T^{+})^{\circ}\), and \((T^{+})^{\circ}\derives\psi\).

Conversely, let a derivation of \(\psi\) with open assumptions in
\((T^{+})^{\circ}\) be given.  Encode it in postorder as in the converse
direction of the proof of \cref{thm:arithmetization-syntax}; all row
checks, parent counts and the root test succeed.  For each root open leaf,
with formula \(\theta\), choose \(s_v\) as follows.  If \(\theta=\tau^{\circ}\)
with \(\tau\in T\), then \(\ulcorner\theta\urcorner\in C_T\), the partial
function of \(e_T\) is defined there, and \cref{thm:bounded-verification}
gives \(s_v\) with \(V(e_T,\ulcorner\theta\urcorner,s_v)=1\).  Otherwise
\(\theta=(\alpha^{\mathbb Z})^{\circ}\) with \(\alpha\in\PA^{\rm d}\); put
\(s_v=\ulcorner\alpha\urcorner\), so that
\(\mathsf{Ax}^{\rm d}_{\PA}(s_v)=1\) and
\(\mathrm{tr}(s_v)=\ulcorner\theta\urcorner\).  In both cases
\(\mathsf{Lf}(\ulcorner\theta\urcorner,s_v)=1\), and every leaf of
\((T^{+})^{\circ}\) is a sentence.  With \(W\) the list of these \(s_v\) in
the order of the root open list (\(W=0\) if there is none),
\(\mathsf{Pf}^{+}(N_{16}(N_{15}(R),W),\ulcorner\psi\urcorner)\).

\emph{Part~3.}  Let \(\sigma\) be an \(\Lang_A\)-sentence.  Then
\(\sigma\in T^{\sharp}\) iff \(T^{+}\models\sigma^{\mathbb Z}\) iff
\((T^{+})^{\circ}\models(\sigma^{\mathbb Z})^{\circ}\)
(\cref{thm:p02-w02c-p02-003-ring-transfer}) iff
\((T^{+})^{\circ}\derives(\sigma^{\mathbb Z})^{\circ}\)
(\cref{thm:fol-completeness,thm:fol-soundness}) iff
\(\exists p\,\mathsf{Pf}^{+}(p,\ulcorner(\sigma^{\mathbb Z})^{\circ}\urcorner)\)
(part~2) iff \(\exists p\,\mathsf{Pf}^{+}(p,\mathrm{tr}(\ulcorner\sigma\urcorner))\)
iff \(\exists p\,\mathsf{En}^{\sharp}(p,\ulcorner\sigma\urcorner)\).  If
\(x\) is not the code of an \(\Lang_A\)-sentence, no \(p\) satisfies
\(\mathsf{En}^{\sharp}(p,x)\).  So the set of codes of \(T^{\sharp}\) is
\(\{x:\exists p\,\mathsf{En}^{\sharp}(p,x)\}\).

The function \(g(x,p)=Z(\chi_{\mathsf{En}^{\sharp}}(p,x))\) is primitive
recursive by part~1, with \(Z(0)=1\) and \(Z(n+1)=0\).  The class of partial
recursive functions of \cref{def:p04-partial-recursive} contains the zero,
successor and projection functions and is closed under composition and
primitive recursion, so it contains every primitive recursive function
(\cref{sec:turing-machines-02}), by induction on the
generation of that class; in particular it contains \(g\), and, being closed
under unbounded minimization, the partial function
\(h(x)=\mu p\,[g(x,p)=0]\).  Since \(g\) is total, \(h(x)\) is defined iff
\(g(x,p)=0\) for some \(p\), iff \(\mathsf{En}^{\sharp}(p,x)\) for some
\(p\).  So the domain of \(h\) is the set of codes of \(T^{\sharp}\).  By
\cref{thm:p04-tm-partial-recursive}, \(h\) is computed by a Turing machine,
so it is partial computable (\cref{def:partial-computable}), and its
domain is computably enumerable (\cref{sec:turing-machines-03}).
\end{proof}
\end{annexitem}

\begin{annexitem}{annex:p02-w02c-p02-003-sharp-provability}{thm:p02-w02c-p02-003-sharp-provability}{Proof existence and D1--D3 for \(\mathsf{Prv}^{\sharp}\)}
\begin{proof}
\emph{Part~1.}  If \(\alpha\in\PA^{\rm d}\), then
\(\alpha^{\mathbb Z}\in T^{+}\), so \(T^{+}\models\alpha^{\mathbb Z}\) and
\(\alpha\in T^{\sharp}\).  If \(\PA\derives\psi\), then
\(\PA^{\rm d}\derives\psi\) by \cref{thm:p05-direct-pa-axioms}, so
\(T^{\sharp}\derives\psi\).

\emph{Codes of \(\mathcal A^{\sharp}\).}  Each proper child code is less
than its parent (\cref{sec:arithmetic-coding-04}), so
\(\mathrm{pd}(k+1,c)>\mathrm{pd}(k,c)\) and, by induction on \(k\),
\(\mathrm{pd}(k,c)\ge c+k\).  Hence if \(a=\mathrm{pd}(k,x)\) then
\(k\le a\) and \(x\le a\), and the bounded sums in \(\mathrm{cr}\) reach
every such pair: \(\mathrm{cr}(a)=1\) iff \(a=\mathrm{pd}(k,x)\) for some
\(k,x\) with \(\mathsf{En}^{\sharp}(k,x)\).  If \(\mathsf{En}^{\sharp}(k,x)\),
then \(x\) is a sentence code and so is \(\mathrm{pd}(k,x)\).  So \(\mathrm{cr}\)
is the characteristic function of the set of codes of \(\mathcal A^{\sharp}\),
and it is primitive recursive, being built by composition and bounded sums
from \(\mathrm{pd}\) (a primitive recursion on \(k\) through the constructor
\(N_8\)) and \(\mathsf{En}^{\sharp}\)
(\cref{thm:p02-w02c-p02-003-sharp-enumeration}).  This is why a
description \(\mathsf{Cr}\) of it exists.

\emph{\(T^{\sharp}\) and \(\PA^{\rm d}\cup\mathcal A^{\sharp}\) have the
same consequences.}  Each member of either set follows from the other set.
\(\PA^{\rm d}\subseteq T^{\sharp}\).  Let \(\theta\in\mathcal A^{\sharp}\),
\(\ulcorner\theta\urcorner=\mathrm{pd}(k,\ulcorner\sigma\urcorner)\) with
\(\mathsf{En}^{\sharp}(k,\ulcorner\sigma\urcorner)\).  Then
\(\sigma\in T^{\sharp}\) by \cref{thm:p02-w02c-p02-003-sharp-enumeration},
and \(\theta\) is derived from \(\sigma\) by \(k\) conjunction
introductions with the premise-free proof of \(0=0\).  Let
\(\sigma\in T^{\sharp}\).  By \cref{thm:p02-w02c-p02-003-sharp-enumeration}
there is \(k\) with \(\mathsf{En}^{\sharp}(k,\ulcorner\sigma\urcorner)\);
the sentence \(\theta\) with code \(\mathrm{pd}(k,\ulcorner\sigma\urcorner)\)
is in \(\mathcal A^{\sharp}\), and \(\sigma\) is derived from it by \(k\)
right conjunction eliminations.  By soundness
(\cref{thm:fol-soundness}) each set is true in every model of the other, so
the two sets have the same models, and by
\cref{thm:fol-completeness,thm:fol-soundness} they derive the same
sentences.

\emph{The numerical meaning of \(\chi_{\sharp}\).}  The proof of
\cref{thm:p05-selected-direct-pa-graph} in \cref{annex:p5.checker-description}
identifies each factor of \(\chi_{19}(p,x)\) separately: the outer tests
hold iff \(p=N_{19}(N_{15}(R),0)\) with exact arities; \(\mathsf{Tree}(R)=1\)
iff \(R\) passes the twenty-rule row checks and the parent-count test;
\(\mathsf{Sent}(\mathsf{root}(R))E(x,\mathsf{root}(R))=1\) iff the root is a
sentence with code \(x\); and the description \(\mathsf{Ax}_{19}\) computes
the characteristic function of \(\PA^{\rm d}\).  The description
\(\chi_{\sharp}\) has the same factors except that its final product applies
\(\mathsf{Ax}_{\sharp}\), whose value at \(a\) is
\(\mathsf{sg}(\mathsf{Ax}^{\rm d}_{\PA}(a)+\mathrm{cr}(a))\), equal to one
iff \(a\) codes a member of \(\PA^{\rm d}\cup\mathcal A^{\sharp}\).  A
product of zero-one values is one iff each factor is one.  Hence
\(f_{d_{\sharp}}(p,x)=1\) iff \(p\) satisfies items~1--3 of
\cref{def:p05-standard-pa-numerical-proof} with root code \(x\) and every
formula in the final open list lies in \(\PA^{\rm d}\cup\mathcal A^{\sharp}\),
and \(f_{d_{\sharp}}(p,x)=0\) otherwise.  The two directions of the proof of
\cref{thm:p05-standard-pa-numerical-proof} (\cref{annex:p5.standard-pa-checker})
use the axiom set only to say which final open leaves are allowed: the row
induction turns such a \(p\) into an \(\mathsf{ND}_{\mathsf{FOL}}\) derivation of the root
with open leaves in the allowed set, and the postorder encoding turns such a
derivation into such a \(p\).  With \(\PA^{\rm d}\cup\mathcal A^{\sharp}\)
as the allowed set, for every \(\Lang_A\)-sentence \(\sigma\),
\(\exists p\,f_{d_{\sharp}}(p,\ulcorner\sigma\urcorner)=1\) iff
\(\PA^{\rm d}\cup\mathcal A^{\sharp}\derives\sigma\).

By \cref{thm:p05-pa-pr-graphs}, \(\Gamma_{d_{\sharp}}\) is \(\Sigma^0_1\)
and \(\mathcal N\models\Gamma_{d_{\sharp}}(\bar p,\bar x,\bar1)\) iff
\(f_{d_{\sharp}}(p,x)=1\).  The \(\Sigma^0_1\) class is closed under unbounded
existential quantification, so \(\mathsf{Prv}^{\sharp}\) is \(\Sigma^0_1\),
and \(\mathcal N\models\mathsf{Prv}^{\sharp}(\overline{\ulcorner\sigma
\urcorner})\) iff some \(p\) has \(f_{d_{\sharp}}(p,\ulcorner\sigma\urcorner)=1\)
iff \(\PA^{\rm d}\cup\mathcal A^{\sharp}\derives\sigma\) iff
\(T^{\sharp}\derives\sigma\).

\emph{Part~2.}  If \(T^{\sharp}\derives\sigma\), part~1 gives a standard
\(p\) with \(f_{d_{\sharp}}(p,\ulcorner\sigma\urcorner)=1\).  The numeral
clause of \cref{thm:p05-pa-pr-graphs} gives
\(\PA\derives\forall y\,(\Gamma_{d_{\sharp}}(\bar p,
\overline{\ulcorner\sigma\urcorner},y)\leftrightarrow y=\bar1)\), hence
\(\PA\derives\mathsf{Prf}^{\sharp}(\bar p,\overline{\ulcorner\sigma\urcorner})\),
and existential introduction gives
\(\PA\derives\mathsf{Prv}^{\sharp}(\overline{\ulcorner\sigma\urcorner})\).

\emph{Part~3.}  Work in \(\PA\) with the unique graph outputs of
\cref{thm:p05-pa-pr-graphs}, as in \cref{annex:p5.checker-internal-mp};
every function application below abbreviates its graph.  Assume
\(\chi_{\sharp}(p,N_{10}(u,v))=1\) and \(\chi_{\sharp}(q,u)=1\).  Unfolding
the two descriptions as there, the zero-one product property gives the exact
tag-19 decodings \(p=N_{19}(N_{15}(R_p),0)\), \(q=N_{19}(N_{15}(R_q),0)\), the
tests \(\mathsf{Tree}(R_p)=\mathsf{Tree}(R_q)=1\), the root equalities with
\(N_{10}(u,v)\) and \(u\), the two sentence tests, and the value one for the
final product of each record; by bounded-product induction every factor
\(\mathsf{Ax}_{\sharp}(f_{R}(e(\mathsf{open}(R),j)))\), \(j<\ell(\mathsf{open}(R))\),
is one for \(R=R_p\) and \(R=R_q\).  Every factor except the final product is
the same description in \(\chi_{\sharp}\) as in \(\chi_{19}\).

Let \(a=\ell(R_q)\), \(b=\ell(R_p)\) and
\[
 R'=\mathsf{cat}\bigl(R_q,\mathsf{cat}(\mathsf{SR}(a,R_p),
 [[8,v,[D(a),D(a+b)],0,0]])\bigr).
\]  The paragraphs of
\cref{annex:p5.checker-internal-mp} headed ``The offset equations'', ``The
open-occurrence invariant'' and ``The parent-count invariant'' prove in
\(\PA\), from the row-bit and tree hypotheses on \(R_q,R_p\) just obtained,
the lookup equations for \(R'\), that all its row bits are one, its exact
parent-count product, and
\[
 \mathsf{open}(R')=\mathsf{cat}(\mathsf{open}(R_q),
                       \mathsf{sh}(a,\mathsf{open}(R_p))).
\]
No axiom description occurs in those paragraphs.  As in the paragraph
headed ``The last row and direct axiom tests'' there, the old roots occupy
positions \(D(a)\) and \(D(a+b)\) with formulas \(u\) and \(N_{10}(u,v)\),
the appended row passes its tag-8 check, and
\(\mathsf{Sent}(v)=1\) by the implication-subformula sentence lemma.

The final product for \(R'\).  At an entry of \(\mathsf{open}(R')\) from the
first segment, the leaf formula is the \(R_q\) formula at the same index; at
an entry \(a+j\) from the second segment it is \(f_{R_p}(j)\), by the lookup
equations.  So each \(\mathsf{Ax}_{\sharp}\) call in the new product has
exactly the argument of a call already shown to have value one, and since
the graph of the description \(\mathsf{Ax}_{\sharp}\) has unique outputs,
it has value one.  By the product/concatenation and map-entry equations the
new product is the product of the two old ones, hence one; the empty open
list gives the initial value one.  Therefore
\(\chi_{\sharp}(N_{19}(N_{15}(R'),0),v)=1\).  The success guard of
\(\mathsf{mp}_{\sharp}\) holds by the two premises and constructor
injectivity, so its specified graph output is this code.  We have proved in
\(\PA\), with universal closure,
\(\chi_{\sharp}(p,N_{10}(u,v))=1\land\chi_{\sharp}(q,u)=1\to\exists r\,
(\Gamma_{d^{\sharp}_{\rm mp}}(p,q,r)\land\chi_{\sharp}(r,v)=1)\), which is the
displayed assertion.

For D2, fix \(\varphi,\psi\).  The code of \(\varphi\to\psi\) is
\(N_{10}(\ulcorner\varphi\urcorner,\ulcorner\psi\urcorner)\), and finite
evaluation of its description gives this equation in \(\PA\).  Assume
\(\mathsf{Prv}^{\sharp}(\overline{\ulcorner\varphi\to\psi\urcorner})\) and
\(\mathsf{Prv}^{\sharp}(\overline{\ulcorner\varphi\urcorner})\); eliminate
the two existentials with fresh \(p,q\), apply the assertion just proved at
\(u=\overline{\ulcorner\varphi\urcorner}\),
\(v=\overline{\ulcorner\psi\urcorner}\), eliminate the existential in
\(r\), introduce
\(\mathsf{Prv}^{\sharp}(\overline{\ulcorner\psi\urcorner})\), and discharge
the two assumptions by implication introduction.

\emph{Part~4.}  By \cref{def:p05-selected-direct-pa-graph} and the
case-distinction convention of \cref{def:p05-checker-description-language},
\(\chi_{19}(p,x)=I(b,M(Y,P_{19}),0)\), where \(b\) is the value of the outer
tests,
\[
 Y=M\bigl(M(\mathsf{Tree}(R),\mathsf{Sent}(\mathsf{root}(R))),
       E(x,\mathsf{root}(R))\bigr),
\]
and \(P_{19}\) is the value at \(n=\ell(\mathsf{open}(R))\) of the fold
\[
 F_{19}(0)=1,\qquad F_{19}(j+1)=M(F_{19}(j),\mathsf{Ax}_{19}(g(j))),
\]
where \(g(j)=f_R(e(\mathsf{open}(R),j))\);
and \(\chi_{\sharp}(p,x)=I(b,M(Y,P_{\sharp}),0)\) with the same
\(b,Y,g,n\) and the fold \(F_{\sharp}\) using \(\mathsf{Ax}_{\sharp}\).  In \(\PA\), the graph equations give
\(A(x,y)=x+y\) (\cref{ex:p05-pa-pair-cons-graphs}) and \(M(x,y)=x\times y\):
by induction on \(y\), \(M(x,0)=0=x\times0\) and
\(M(x,Sy)=A(M(x,y),x)=x\times y+x=x\times Sy\).  Also \(Z(0)=1\),
\(Z(Sn)=0\), \(\mathsf{sg}(x)=Z(Z(x))\), and
\(I(c,X,W)=\mathsf{sg}(c)\times X+Z(c)\times W\).

Assume \(\chi_{19}(p,x)=1\).  If \(b=0\), then \(\mathsf{sg}(b)=Z(1)=0\) and
\(\chi_{19}(p,x)=0\), which contradicts \(0\ne S0\).  So \(b=Sm\) for some
\(m\), \(Z(b)=0\), \(\mathsf{sg}(b)=1\), and \(Y\times P_{19}=1\).  In
\(\PA\), \(x\times y=1\) implies \(x=1\) and \(y=1\): if \(x=0\) or \(y=0\)
the product is \(0\); otherwise \(y=Sy'\) and \(x\times y=x\times y'+x\ge x\),
so \(x\le1\), hence \(x=1\) and then \(y=1\).  So \(Y=1\) and \(P_{19}=1\).
Now prove by induction on \(j\) that \(F_{19}(j)=1\) implies
\(F_{\sharp}(j)=1\), the parameters \(p,x\) being fixed.  At \(0\) both values
are \(1\).  If \(F_{19}(j+1)=F_{19}(j)\times\mathsf{Ax}_{19}(g(j))=1\), then
\(F_{19}(j)=1\) and \(\mathsf{Ax}_{19}(g(j))=1\); by the induction hypothesis
\(F_{\sharp}(j)=1\); and with \(c=\mathsf{Cr}(g(j))\), the unique output
supplied by \cref{thm:p05-pa-pr-graphs},
\(\mathsf{Ax}_{\sharp}(g(j))=\mathsf{sg}(1+c)=\mathsf{sg}(Sc)=Z(0)=1\).
So \(F_{\sharp}(j+1)=1\times1=1\).  At \(j=n\), \(P_{\sharp}=1\), hence
\(\chi_{\sharp}(p,x)=\mathsf{sg}(b)\times(Y\times1)+Z(b)\times0=1\).  This
proves \(\PA\derives\forall p\,\forall x\,(\Gamma_{d_{19}}(p,x,\bar1)\to
\Gamma_{d_{\sharp}}(p,x,\bar1))\), that is, part~4.

\emph{Part~5.}  Fix \(\varphi\) and let \(\sigma\) be the sentence
\(\mathsf{Prv}^{\sharp}(\overline{\ulcorner\varphi\urcorner})\), that is
\(\exists p\,\Gamma_{d_{\sharp}}(p,\overline{\ulcorner\varphi\urcorner},\bar1)\).
By \cref{thm:p05-literal-graph-evidence}, whose first assertion is proved by
induction over an arbitrary fixed graph description, the formula
\(\Gamma_{d_{\sharp}}\) and each of its numeral instances have a positive
formation already in the fixed marked negation normal form; the outer
unbounded existential quantifier is kept by the normal-form conversion.
So \(\sigma\) is a fixed closed \(\Sigma^0_1\) sentence in its original
syntax, equal to its own normal form, and
\cref{thm:p05-pa-sigma-evidence,thm:p05-original-sigma-evidence} give
\(\PA\derives\sigma\to\mathsf{Prv}^{19}_{\PA}(\overline{\ulcorner\sigma
\urcorner})\).  Part~4, instantiated at
\(x=\overline{\ulcorner\sigma\urcorner}\) and followed by existential
elimination and introduction, gives
\(\PA\derives\mathsf{Prv}^{19}_{\PA}(\overline{\ulcorner\sigma\urcorner})
\to\mathsf{Prv}^{\sharp}(\overline{\ulcorner\sigma\urcorner})\).  Chaining
the two implications gives D3.

Finally, every theorem of \(\PA\) is a theorem of \(T^{\sharp}\) by
part~1, so the conclusions of parts~2, 3 and~5 hold with \(T^{\sharp}\) in
place of \(\PA\).
\end{proof}
\end{annexitem}

\begin{annexitem}{annex:p02-w02c-p02-003-integers-second}{thm:p02-w02c-p02-003-integers-second}{A truth of \(\mathbb Z\) not implied by \(T\)}
\begin{proof}
\emph{\(\PA^{\rm d}\) is true in \(\mathcal N\).}  By
\cref{thm:p05-direct-pa-axioms}, each member of \(\PA^{\rm d}\) is the
universal closure of an axiom of \(\mathsf Q\) or of an induction instance.
Each axiom of \(\mathsf Q\) states a property of successor, addition or
multiplication of natural numbers (zero is not a successor, successor is
injective, every non-zero number is a successor, and the recursion
equations for \(+\) and \(\times\)), and each holds in \(\N\).  For an
instance of induction with parameters \(\vec z\) fixed in \(\N\), suppose
\(\varphi(0,\vec z)\) and \(\forall x\,(\varphi(x,\vec z)\to
\varphi(Sx,\vec z))\) hold in \(\mathcal N\).  By induction on \(m\) in the
metatheory, \(\varphi(m,\vec z)\) holds for every \(m\in\N\), so
\(\forall x\,\varphi(x,\vec z)\) holds.

\emph{Part~1.}  \(T^{\sharp}\supseteq\PA^{\rm d}\) by part~1 of
\cref{thm:p02-w02c-p02-003-sharp-provability}, and \(\PA^{\rm d}\) contains
the seven codes \(q_1,\ldots,q_7\) of \cref{sec:p05-standard-pa-checker}.
Every member of \(T^{+}\) is true in \(\mathbb Z\): members of \(T\) by
hypothesis, and \(\alpha^{\mathbb Z}\) for \(\alpha\in\PA^{\rm d}\) by
part~1 of \cref{thm:p02-w02c-p02-003-integers}, since
\(\mathcal N\models\alpha\).  So \(\mathbb Z\) is a model of \(T^{+}\).  If
\(\sigma\in T^{\sharp}\), then \(T^{+}\models\sigma^{\mathbb Z}\), so
\(\mathbb Z\models\sigma^{\mathbb Z}\), and by part~1 of
\cref{thm:p02-w02c-p02-003-integers} \(\mathcal N\models\sigma\).  Thus
\(\mathcal N\) is a model of \(T^{\sharp}\).  By soundness
(\cref{thm:fol-soundness}), every sentence derivable from \(T^{\sharp}\) is
true in \(\mathcal N\); \(\sigma\) and \(\neg\sigma\) are not both true
there, so \(T^{\sharp}\) is consistent.  \(T^{\sharp}\) is computably
enumerable by part~3 of \cref{thm:p02-w02c-p02-003-sharp-enumeration}.

\emph{Part~2.}  By part~1, \(T^{\sharp}\) is a consistent computably
enumerable extension of \(\mathsf Q\).  By
\cref{thm:p02-w02c-p02-003-sharp-provability}, \(\mathsf{Prv}^{\sharp}\) is
a specified formula with the proof-existence meaning required in
\cref{sec:second-incompleteness-01} and satisfies D1--D3 for \(T^{\sharp}\).
With this formula as \(\mathsf{Prov}_{T^{\sharp}}\), \(\mathsf{Con}(T^{\sharp})\)
is \(\mathsf{Con}^{\sharp}\), and \cref{thm:godel-second} gives
\(T^{\sharp}\notderives\mathsf{Con}^{\sharp}\).

\((\mathsf{Con}^{\sharp})^{\mathbb Z}\) is true in \(\mathbb Z\).  By part~1
of \cref{thm:p02-w02c-p02-003-sharp-provability},
\(\mathcal N\models\mathsf{Prv}^{\sharp}(\overline{\ulcorner\bottom_A
\urcorner})\) iff \(T^{\sharp}\derives\bottom_A\), where \(\bottom_A\) is
\(0=S0\).  \(\bottom_A\) is false in \(\mathcal N\), a model of
\(T^{\sharp}\), so by soundness \(T^{\sharp}\notderives\bottom_A\).  Hence
\(\mathcal N\models\mathsf{Con}^{\sharp}\), and by part~1 of
\cref{thm:p02-w02c-p02-003-integers},
\(\mathbb Z\models(\mathsf{Con}^{\sharp})^{\mathbb Z}\).

\(T^{+}\not\models(\mathsf{Con}^{\sharp})^{\mathbb Z}\).  Otherwise
\(\mathsf{Con}^{\sharp}\in T^{\sharp}\) by the definition of
\(T^{\sharp}\); a member of \(T^{\sharp}\) is derivable from \(T^{\sharp}\), so
\(T^{\sharp}\derives\mathsf{Con}^{\sharp}\), which was just excluded.  Since
\(T\subseteq T^{+}\), every model of \(T^{+}\) is a model of \(T\), so
\(T\models(\mathsf{Con}^{\sharp})^{\mathbb Z}\) would give
\(T^{+}\models(\mathsf{Con}^{\sharp})^{\mathbb Z}\); hence
\(T\not\models(\mathsf{Con}^{\sharp})^{\mathbb Z}\).
\end{proof}
\end{annexitem}

\begin{annexitem}{annex:p02-w02c-p02-003-arbitrary-representative}{thm:p02-w02c-p02-003-arbitrary-representative}{A representative with a provable consistency sentence}
\begin{proof}
\emph{Distinct numerals.}  If \(m\ne n\), then
\(\mathsf Q\derives\neg\,\bar m=\bar n\).  By induction on \(\min(m,n)\):
if, say, \(m=0<n\), then \(\bar n=S\overline{n-1}\), and \(0=S\overline{n-1}\)
contradicts the axiom \(\forall x\,\neg\,Sx=0\) after symmetry of equality;
if \(m=m'+1\) and \(n=n'+1\), then \(S\bar m'=S\bar n'\) gives
\(\bar m'=\bar n'\) by the axiom
\(\forall x\,\forall y\,(Sx=Sy\to x=y)\), contradicting the induction
hypothesis.

\emph{\(\rho'\) is \(\Sigma^0_1\).}  The negation-normal-form conversion of
\(\rho\land\neg\,x=\bar c\) is the conjunction of that of \(\rho\), which is
in the \(\Sigma^0_1\) class, and the negated equality, which is a bounded
formula in negation normal form; the class is closed under \(\land\).

\emph{Numeralwise representation.}  If \(\mathsf{Proof}_S(p,x)\) holds, then
\(x\ne c\): otherwise \(S\derives\bottom_A\) by
\cref{thm:arithmetization-syntax}, and with the axiom
\(\forall x\,\neg\,Sx=0\) of \(\mathsf Q\subseteq S\) and symmetry of
equality this makes \(S\) inconsistent.  So \(\mathsf Q\derives\rho(\bar p,\bar x)\)
and \(\mathsf Q\derives\neg\,\bar x=\bar c\), hence
\(\mathsf Q\derives\rho'(\bar p,\bar x)\).  If \(\mathsf{Proof}_S(p,x)\)
fails, then \(\mathsf Q\derives\neg\rho(\bar p,\bar x)\), hence
\(\mathsf Q\derives\neg\rho'(\bar p,\bar x)\).

\emph{Meaning and D1.}  \(\mathsf Q\) is true in \(\mathcal N\), so by
soundness the two clauses give
\(\mathcal N\models\rho'(\bar p,\bar x)\) iff \(\mathsf{Proof}_S(p,x)\).
With \cref{thm:arithmetization-syntax},
\(S\derives\sigma\) iff
\(\mathcal N\models\exists p\,\rho'(p,\overline{\ulcorner\sigma\urcorner})\).
If \(S\derives\sigma\), a standard \(p\) has
\(\mathsf{Proof}_S(p,\ulcorner\sigma\urcorner)\), so
\(\mathsf Q\derives\rho'(\bar p,\overline{\ulcorner\sigma\urcorner})\), and
existential introduction gives
\(S\derives\exists p\,\rho'(p,\overline{\ulcorner\sigma\urcorner})\), since
\(\mathsf Q\subseteq S\).

\emph{The consistency sentence is provable.}  Assume
\(\exists p\,\rho'(p,\bar c)\).  Eliminate the existential with a fresh
\(p\); conjunction elimination gives \(\neg\,\bar c=\bar c\), and equality
introduction gives \(\bar c=\bar c\), hence \(\bottom\).  Negation
introduction gives \(\derives\neg\exists p\,\rho'(p,\bar c)\), so
\(S\derives\neg\exists p\,\rho'(p,\bar c)\).

\emph{Failure of D2 or D3.}  The sentence \(\neg\exists p\,\rho'(p,\bar c)\)
is \(\mathsf{Con}(S)\) of \cref{def:formal-consistency-sentence} for the
specified formula \(\exists p\,\rho'\), which has the proof-existence
meaning of \cref{sec:second-incompleteness-01} and satisfies D1.  If it also
satisfied D2 and D3, \cref{thm:godel-second} would give
\(S\notderives\mathsf{Con}(S)\), contradicting the last paragraph.
\end{proof}
\end{annexitem}
